% Options for packages loaded elsewhere
\PassOptionsToPackage{unicode}{hyperref}
\PassOptionsToPackage{hyphens}{url}
\PassOptionsToPackage{dvipsnames,svgnames,x11names}{xcolor}
%
\documentclass[
  10pt,
]{article}
\usepackage{amsmath,amssymb}
\usepackage{iftex}
\ifPDFTeX
  \usepackage[T1]{fontenc}
  \usepackage[utf8]{inputenc}
  \usepackage{textcomp} % provide euro and other symbols
\else % if luatex or xetex
  \usepackage{unicode-math} % this also loads fontspec
  \defaultfontfeatures{Scale=MatchLowercase}
  \defaultfontfeatures[\rmfamily]{Ligatures=TeX,Scale=1}
\fi
\usepackage{lmodern}
\ifPDFTeX\else
  % xetex/luatex font selection
  \setmainfont[]{DejaVu Serif}
  \setsansfont[]{DejaVu Sans}
  \setmonofont[]{DejaVu Sans Mono}
  \setmathfont[]{Latin Modern Math}
\fi
% Use upquote if available, for straight quotes in verbatim environments
\IfFileExists{upquote.sty}{\usepackage{upquote}}{}
\IfFileExists{microtype.sty}{% use microtype if available
  \usepackage[]{microtype}
  \UseMicrotypeSet[protrusion]{basicmath} % disable protrusion for tt fonts
}{}
\makeatletter
\@ifundefined{KOMAClassName}{% if non-KOMA class
  \IfFileExists{parskip.sty}{%
    \usepackage{parskip}
  }{% else
    \setlength{\parindent}{0pt}
    \setlength{\parskip}{6pt plus 2pt minus 1pt}}
}{% if KOMA class
  \KOMAoptions{parskip=half}}
\makeatother
\usepackage{xcolor}
\usepackage[margin=2.0cm]{geometry}
\usepackage{longtable,booktabs,array}
\usepackage{calc} % for calculating minipage widths
% Correct order of tables after \paragraph or \subparagraph
\usepackage{etoolbox}
\makeatletter
\patchcmd\longtable{\par}{\if@noskipsec\mbox{}\fi\par}{}{}
\makeatother
% Allow footnotes in longtable head/foot
\IfFileExists{footnotehyper.sty}{\usepackage{footnotehyper}}{\usepackage{footnote}}
\makesavenoteenv{longtable}
\setlength{\emergencystretch}{3em} % prevent overfull lines
\providecommand{\tightlist}{%
  \setlength{\itemsep}{0pt}\setlength{\parskip}{0pt}}
\setcounter{secnumdepth}{-\maxdimen} % remove section numbering
% doc_header.tex -- LaTeX preamble additions used by build_docs.sh for every
% Markdown -> LaTeX -> PDF conversion of the provenance of the notebook 3+1super-loop.nb.
\usepackage{etoolbox}
% longtable may not be loaded by pandoc when a document has no table
\usepackage{longtable}
% bold vectors: \bm -> \symbf under unicode-math (XeLaTeX), bm package otherwise
\makeatletter
\@ifpackageloaded{unicode-math}{\providecommand{\bm}[1]{\symbf{#1}}}{\usepackage{bm}}
\makeatother
% smaller type in long tables so that wide audit tables fit the text width;
% inside a table a line break is allowed after an underscore of plain text
% (pandoc writes it as \_), so that long identifiers wrap inside their column
% rows are set a little apart, so that symbols with indices do not touch
\AtBeginEnvironment{longtable}{\footnotesize\renewcommand{\arraystretch}{1.45}%
  \let\sptunderscore\_\def\_{\sptunderscore\allowbreak}}
\setlength{\LTcapwidth}{\textwidth}
% allow line breaks after slashes and hyphens inside table cells
\usepackage{xurl}
% code blocks: break over-long lines instead of running into the margin.
% Plain ``` blocks are typeset by pandoc as {verbatim}, highlighted ones
% (```bash) as {Highlighting}; both are re-declared as breakable fancyvrb
% environments.  The continuation line is marked with a small hooked arrow.
\usepackage{fvextra}
\RecustomVerbatimEnvironment{verbatim}{Verbatim}{breaklines,breakanywhere,fontsize=\small}
\makeatletter
\@ifundefined{Highlighting}{}{%
  \RecustomVerbatimEnvironment{Highlighting}{Verbatim}{breaklines,breakanywhere,fontsize=\small,commandchars=\\\{\}}}
\makeatother
% inline code (pandoc writes it as \texttt{...}): never hyphenate -- an inserted
% hyphen could be read as part of a file name -- but allow a line break after
% an underscore, slash, comma, colon or equals sign, so that long paths and
% column lists wrap inside the text width and inside table columns.
% Inside a formula (e.g. the name of a notebook symbol in an equation) the same
% command must first leave math mode, which \text of amsmath does.
\ExplSyntaxOn
\tl_new:N \l__spt_code_tl
\cs_new_protected:Npn \__spt_texttt:n #1
  {
    \tl_set:Nn \l__spt_code_tl {#1}
    \regex_replace_all:nnN { ([\/\,\:\=]) } { \1 \c{allowbreak} } \l__spt_code_tl
    \regex_replace_all:nnN { \c{_} } { \c{_} \c{allowbreak} } \l__spt_code_tl
    \group_begin: \ttfamily \hyphenchar \font = -1 \scan_stop: \l__spt_code_tl \group_end:
  }
\RenewDocumentCommand \texttt { m }
  {
    \mode_if_math:TF { \text { \__spt_texttt:n {#1} } } { \__spt_texttt:n {#1} }
  }
\ExplSyntaxOff
% A figure is written in the Markdown source as an image paragraph followed by
% a caption paragraph ("**Figure S1.** ..."), so that the caption is also
% visible on GitHub.  An image that starts a paragraph is centered.
\usepackage{graphicx}% pandoc loads it only for documents that contain an image
\NewCommandCopy\sptincludegraphics\includegraphics
\RenewDocumentCommand\includegraphics{O{} m}{%
  \ifvmode\begin{center}\sptincludegraphics[#1]{#2}\end{center}%
  \else\sptincludegraphics[#1]{#2}\fi}
% keep figures where they are referenced
\usepackage{float}
\floatplacement{figure}{H}
% running header with the document title
\usepackage{fancyhdr}
\pagestyle{fancy}
\fancyhf{}
\fancyhead[L]{\small PROJECT spacetimes, notebook 3+1super-loop.nb}
\fancyhead[R]{\small\thepage}
\renewcommand{\headrulewidth}{0.4pt}
% 16x16 matrices: tighter columns so that a full matrix fits the text width
\setlength{\arraycolsep}{3.2pt}
% The documents of this project contain long tagged symbols, file paths and
% identifiers that cannot be hyphenated.  A paragraph is therefore allowed to
% have looser lines before TeX gives up and lets a line stick out of the margin.
\AtBeginDocument{\tolerance=2000 \emergencystretch=5em \hbadness=2000}
\ifLuaTeX
  \usepackage{selnolig}  % disable illegal ligatures
\fi
\IfFileExists{bookmark.sty}{\usepackage{bookmark}}{\usepackage{hyperref}}
\IfFileExists{xurl.sty}{\usepackage{xurl}}{} % add URL line breaks if available
\urlstyle{same}
\hypersetup{
  pdftitle={Provenance P\_3+1super-loop: the Wolfram Language notebook 3+1super-loop.nb, Part I (the symbols) and Part II (the encyclopedia of Test 2): the library and item {[}1{]}},
  colorlinks=true,
  linkcolor={blue},
  filecolor={Maroon},
  citecolor={blue},
  urlcolor={blue},
  pdfcreator={LaTeX via pandoc}}

\title{Provenance P\_3+1super-loop: the Wolfram Language notebook
3+1super-loop.nb, Part I (the symbols) and Part II (the encyclopedia of
Test 2): the library and item {[}1{]}}
\usepackage{etoolbox}
\makeatletter
\providecommand{\subtitle}[1]{% add subtitle to \maketitle
  \apptocmd{\@title}{\par {\large #1 \par}}{}{}
}
\makeatother
\subtitle{PROJECT spacetimes, prepared for Patrick Lee Nash, Ph.D.}
\author{}
\date{18 November 2026}

\begin{document}
\maketitle

\hypertarget{what-this-document-is}{%
\section{0. What this document is}\label{what-this-document-is}}

This is the provenance of the notebook \texttt{3+1super-loop.nb}, which
stands in the top folder of the repository of PROJECT spacetimes. A
\textbf{provenance} is the record of where every part of a piece of work
comes from: what was asked, what was done, how it was done, what was
computed, which errors were met and how they were removed, and which
files hold the evidence. This document describes \textbf{Part I} of the
notebook, the Wolfram Language form of every symbol of the project
(sections 5 to 9), and the finished chapters of \textbf{Part II}, the
encyclopedia of the programs of Test 2 (sections 10 to 12): chapter
II.1, the library of the programs, and chapter II.2, item {[}1{]} of
Test 2. The remaining items of Part II are added to the notebook in
later stages; their provenance will be added to this document at the
same time.

The document is written for a physics student who has not seen the work
before. Every term and every symbol is defined in words in section 2
before it is used. Nothing in this document depends on any other
document: every equation, symbol and definition that it uses is written
here.

\textbf{Status (18 November 2026).} The notebook was evaluated in
Wolfram 15.0.1 (Mac, Apple silicon), in the front end of the application
Wolfram.app, by the menu command Evaluation \textgreater{} Evaluate
Notebook, in a fresh kernel: 82 input cells, 42 output cells, 35 printed
test lines, \textbf{0 messages} (no error or warning of any kind),
\textbf{35 of 35 tests passed} (15 of Part I, 20 of Part II), evaluation
time 71 seconds.

\hypertarget{the-task}{%
\section{1. The task}\label{the-task}}

The author asked for one stand-alone Wolfram Language notebook, named
\texttt{3+1super-loop.nb}, extremely well organized, clearly written and
thoughtfully commented for a student who is learning the work, with
provenance, using only the Wolfram Language, and evaluated in
Mathematica (on his own computer, in
\texttt{/Applications/Wolfram.app}). In his words, restated here:

\begin{enumerate}
\def\labelenumi{\arabic{enumi}.}
\tightlist
\item
  First, read and understand the author's note on how a symbolized
  variable is created by a program (the technique of section 5).
\item
  Then port every unique symbol of the Test 2 sub-project, that is,
  every symbol of the symbol registry of Test 2 (the registry is
  described in section 7), to the notebook, using only the Wolfram
  Language.
\item
  Create non-trivial verification tests that include sums, arithmetic
  and algebra, and execute them in Mathematica. If errors occur, find
  out why (also with the help of the internet), fix them, and document
  the fixes with provenance.
\item
  Then append to the notebook, as an ``encyclopedia'', the entire code
  of Test 2 with all its tests and verifications, ported to the Wolfram
  Language (Part II; done in this stage: the library and item {[}1{]};
  the items {[}2{]} to {[}11{]} follow).
\item
  Evaluate the notebook after every addition, fix every error at once,
  describe and discuss the fixes in the provenance, and push every
  finished item to the repository.
\end{enumerate}

Section 15 lists the decisions that were taken where the task left a
choice.

\hypertarget{terms-and-symbols-used-in-this-document}{%
\section{2. Terms and symbols used in this
document}\label{terms-and-symbols-used-in-this-document}}

\hypertarget{terms-of-the-wolfram-language}{%
\subsection{2.1 Terms of the Wolfram
Language}\label{terms-of-the-wolfram-language}}

\begin{itemize}
\tightlist
\item
  \textbf{Wolfram Language}: the programming language of the system
  Mathematica; on the author's computer the application is
  \texttt{/Applications/Wolfram.app}, version 15.0.1.
\item
  \textbf{Kernel}: the program that evaluates Wolfram Language
  expressions.
\item
  \textbf{Front end}: the program that shows notebooks on the screen,
  lets one type input, sends the input to the kernel and shows the
  results. The front end and the kernel are separate programs that talk
  to each other.
\item
  \textbf{Notebook} (\texttt{.nb} file): a document made of
  \textbf{cells}: input cells (code, evaluated by the kernel), output
  cells (results), text cells, section headings, printed lines and
  message cells (errors and warnings).
\item
  \textbf{Evaluate Notebook}: the menu command Evaluation \textgreater{}
  Evaluate Notebook, which evaluates every input cell of a notebook in
  order.
\item
  \textbf{Expression}, \textbf{symbol}, \textbf{atom}: every object of
  the Wolfram Language is an expression. A \textbf{symbol} is a name
  such as \texttt{x} or \texttt{Sum}. An \textbf{atom} is an expression
  without parts (a symbol, a number or a string); \texttt{AtomQ{[}e{]}}
  is \texttt{True} exactly for atoms. \texttt{Subscript{[}a,\ LL{]}} is
  not an atom: it is the function \texttt{Subscript} applied to the two
  symbols \texttt{a} and \texttt{LL}.
\item
  \textbf{Context}: the name space of a symbol, written before the name
  with the context mark \texttt{\textasciigrave{}}; symbols typed by the
  user belong to the context \texttt{Global\textasciigrave{}}; the
  symbols of this notebook's package belong to
  \texttt{SL31\textasciigrave{}}.
\item
  \textbf{Box}, \textbf{box structure}: the description of how an
  expression looks on the screen. A string such as \texttt{"a"} is a
  box; \texttt{SubscriptBox{[}"a",\ "LL"{]}} is the box of \(a\) with
  the subscript LL; \texttt{SuperscriptBox}, \texttt{SubsuperscriptBox},
  \texttt{OverscriptBox} (an accent over a letter) and
  \texttt{RowBox{[}\{...\}{]}} (several boxes in a row) are other boxes.
  \texttt{StyleBox{[}box,\ options{]}} changes the style (font, weight);
  \texttt{TagBox{[}box,\ tag{]}} carries a hidden tag and does not
  change the appearance.
\item
  \textbf{StandardForm}: the two-dimensional form in which the Wolfram
  Language reads input and shows output. \textbf{TraditionalForm}: the
  form of printed mathematics. \textbf{TeXForm}: the conversion of an
  expression into TeX, which the Wolfram Language makes from the
  TraditionalForm boxes.
\item
  \textbf{Parsing}: reading boxes as an expression.
  \texttt{MakeExpression{[}boxes,\ StandardForm{]}} parses boxes exactly
  as input is parsed and returns the result held,
  \texttt{HoldComplete{[}expression{]}} (\texttt{HoldComplete} keeps its
  contents from being evaluated).
  \texttt{MakeBoxes{[}expression,\ form{]}} makes the boxes of an
  expression; \texttt{ToBoxes} does the same after evaluating.
\item
  \textbf{The Notation package}: a package that comes with every Wolfram
  installation (loaded with
  \texttt{Needs{[}"Notation\textasciigrave{}"{]}}). Its command
  \textbf{\texttt{Symbolize{[}ParsedBoxWrapper{[}box{]},\ options{]}}}
  declares that the box, wherever it is met in input, is read as ONE new
  symbol, and that this symbol is printed as the box. The wrapper
  \texttt{ParsedBoxWrapper} marks its argument as a box structure. The
  option \textbf{\texttt{WorkingForm}} names the form (StandardForm) for
  which the reading rule is made; the option
  \textbf{\texttt{SymbolizeRootName}} gives the name of the new symbol.
\item
  \textbf{Symbolized symbol}: a symbol created by \texttt{Symbolize}. It
  behaves like any symbol: it can be added, multiplied, differentiated,
  given a value, and it is an atom.
\item
  \textbf{\texttt{Set}}, \textbf{\texttt{Clear}},
  \textbf{\texttt{Apply}}: \texttt{Set{[}x,\ v{]}} (written
  \texttt{x\ =\ v}) gives the value \texttt{v} to \texttt{x};
  \texttt{Clear{[}x{]}} removes it; \texttt{Apply{[}f,\ h{[}a,\ b{]}{]}}
  (written \texttt{f\ @@\ h{[}a,\ b{]}}) replaces the head \texttt{h} by
  \texttt{f}.
\item
  \textbf{\texttt{Block}}: \texttt{Block{[}\{x\ =\ v\},\ body{]}}
  evaluates \texttt{body} while \texttt{x} has the value \texttt{v}, and
  restores \texttt{x} afterwards.
\item
  \textbf{\texttt{Sum}}, \textbf{\texttt{FindRoot}},
  \textbf{\texttt{Solve}}, \textbf{\texttt{Simplify}},
  \textbf{\texttt{Factor}}, \textbf{\texttt{Expand}}: summation;
  numerical root finding; exact solution of equations; algebraic
  simplification; factorization; multiplication of products.
\item
  \textbf{\texttt{\$PreRead}}: a global variable; if it holds a
  function, the kernel applies it to the boxes of every input before
  parsing. It is used in section 14 to record exactly which boxes the
  front end sends.
\item
  \textbf{\texttt{SyntaxForm}}: an option of \texttt{TagBox} that tells
  the front end how the box takes part in the grammar of the surrounding
  input; \texttt{SyntaxForm\ -\textgreater{}\ "symbol"} makes the front
  end treat the whole box as one operand, like a symbol.
\item
  \textbf{\texttt{FrontEndTokenExecute{[}nb,\ "EvaluateNotebook"{]}}}: a
  command, sent from the kernel to the front end, that performs the menu
  command Evaluate Notebook for the notebook \texttt{nb}.
\item
  \texttt{FrontEnd\textasciigrave{}UndocumentedTestFEParserPacket{[}string,\ True{]}}:
  an internal command of the front end (not part of the documented
  language) that reads a string exactly as typed input is read when a
  cell is evaluated, and returns the boxes.
\item
  \textbf{wolframscript}: the program that runs a Wolfram Language
  script (a \texttt{.wls} file) in a kernel without the window of the
  application. \textbf{\texttt{UsingFrontEnd{[}expr{]}}} evaluates
  \texttt{expr} with a front end that has no window.
\item
  \textbf{Evaluator} (of a notebook): the kernel to which a notebook
  sends its input; a front end can run several kernels with different
  names.
\item
  \textbf{\texttt{.} (\texttt{Dot})}: the product of matrices and
  vectors, \texttt{a\ .\ b}. In the grammar of the Wolfram Language the
  operator \texttt{.} binds more strongly than the product written by
  placing two expressions side by side (\texttt{Times}):
  \texttt{x\ y\ .\ z} is read as \texttt{x\ (y\ .\ z)}. \texttt{Dot} of
  a number and a matrix is not defined and stays unevaluated (error
  E15).
\item
  \textbf{\texttt{KroneckerProduct}},
  \textbf{\texttt{ConjugateTranspose}}, \textbf{\texttt{ArrayFlatten}}:
  the Kronecker product of matrices; the complex conjugate of the
  transpose (the dagger); a matrix assembled from blocks.
\item
  \textbf{\texttt{BitXor}}, \textbf{\texttt{BitShiftRight}},
  \textbf{\texttt{Mod}}: the exclusive or of the bits of two integers;
  the shift of the bits to the right (division by a power of 2, rounded
  down); the remainder of a division.
\item
  \textbf{\texttt{MatrixExp}}, \textbf{\texttt{Eigenvalues}},
  \textbf{\texttt{Norm{[}m,\ "Frobenius"{]}}}: the matrix exponential;
  the eigenvalues; the square root of the sum of the squared moduli of
  all entries of a matrix (the Frobenius norm).
\item
  \textbf{\texttt{SparseArray}}: a matrix stored by its nonzero entries.
  \textbf{\texttt{Signature}}: the sign of a permutation, 0 if two
  entries are equal (\texttt{Signature{[}\{i,\ j,\ k\}{]}} is the
  Levi-Civita symbol).
\item
  \textbf{\texttt{Format}},
  \textbf{\texttt{MessageName{[}s,\ "usage"{]}}} (written
  \texttt{s::usage}): a rule for printing an expression; the usage
  message that \texttt{?s} prints.
\item
  \textbf{\texttt{ListLinePlot}}, \textbf{\texttt{ListPlot}},
  \textbf{\texttt{ListLogLogPlot}}, \textbf{\texttt{Show}},
  \textbf{\texttt{Graphics}}, \textbf{\texttt{Text}},
  \textbf{\texttt{LineLegend}}, \textbf{\texttt{Labeled}}: plots of
  lists of points (joined, as points, on logarithmic axes); the
  combination of plots; graphics and text placed at coordinates; a
  legend; a caption.
\item
  \textbf{\texttt{SingleLetterItalics}}: an option of the front end; in
  \texttt{TraditionalForm}, the form in which the labels of graphics are
  typeset, a name of one letter is printed in italics unless the option
  is \texttt{False}.
\item
  \textbf{\texttt{ImportString{[}text,\ "TeX"{]}}}: the TeX importer of
  the Wolfram Language, which turns TeX into a notebook with boxes; the
  builder uses it for the formulas in text cells (errors E4 and E17).
\end{itemize}

\hypertarget{symbols-of-the-physics-that-appear-in-the-tests}{%
\subsection{2.2 Symbols of the physics that appear in the
tests}\label{symbols-of-the-physics-that-appear-in-the-tests}}

All of them carry the label LL (the label of the super-loop revision of
the project; section 7).

\begin{longtable}[]{@{}
  >{\raggedright\arraybackslash}p{(\columnwidth - 2\tabcolsep) * \real{0.5000}}
  >{\raggedright\arraybackslash}p{(\columnwidth - 2\tabcolsep) * \real{0.5000}}@{}}
\toprule\noalign{}
\begin{minipage}[b]{\linewidth}\raggedright
symbol
\end{minipage} & \begin{minipage}[b]{\linewidth}\raggedright
meaning
\end{minipage} \\
\midrule\noalign{}
\endhead
\bottomrule\noalign{}
\endlastfoot
\(G_{\rm LL}\) & the coupling constant of the contact interaction of the
spinor charges \\
\(m_{\rm LL}\) & the mass of the spinor charges \\
\(\Lambda_{\rm LL}\) & the momentum cutoff (the unit of the numbers
below) \\
\(\Delta^{\rm vac}_{\rm LL}\) & the gap of the condensed vacuum at zero
chemical potential \\
\(M^{\rm q}_{\rm LL}\) & the mass of the quasi-particles of the
condensed vacuum,
\((M^{\rm q}_{\rm LL})^2=m_{\rm LL}^2+(\Delta^{\rm vac}_{\rm LL})^2\) \\
\(\sigma^{\rm LL}_j\), \(j=1,2,3\) & the Pauli matrices
\(\sigma^{\rm LL}_1=\begin{pmatrix}0&1\\1&0\end{pmatrix}\),
\(\sigma^{\rm LL}_2=\begin{pmatrix}0&-i\\i&0\end{pmatrix}\),
\(\sigma^{\rm LL}_3=\begin{pmatrix}1&0\\0&-1\end{pmatrix}\) \\
\(\mathsf 1_2\) & the \(2\times2\) unit matrix \\
\(\epsilon^{\rm LL}_{jj'j''}\) & the totally antisymmetric symbol:
\(\epsilon^{\rm LL}_{123}=1\), \(-1\) for an odd permutation of
\((1,2,3)\), \(0\) if two indices are equal \\
\(\delta_{jj'}\) & the Kronecker symbol: \(1\) for \(j=j'\), \(0\)
otherwise (used only in this document and in the notebook as the Wolfram
function \texttt{KroneckerDelta}) \\
\(i\) & the imaginary unit, \(i^2=-1\) \\
\(x^p_{\rm LL}\), \(p=0,1,2,3\) & the coordinates of space-time
(\(p=0\): time) \\
\(\eta^{\rm LL}_{pp'}\) & the Minkowski metric:
\(\eta^{\rm LL}_{00}=1\),
\(\eta^{\rm LL}_{11}=\eta^{\rm LL}_{22}=\eta^{\rm LL}_{33}=-1\), zero
for \(p\neq p'\) \\
\(M^{\rm W}_{\rm LL}\), \(M^{\rm F}_{\rm LL}\) & the rest masses of the
two charged vector bosons and of the vector boson of the fermion number,
in the superfluid boson singlet state \\
\(Z^{\rm s}_{\rm LL}\), \(Z^{\rm t}_{\rm LL}\) & the stiffness
coefficients of the space derivatives and of the time derivatives in the
kinetic Lagrangian of the order parameter \\
\(g^{\rm v}_{\rm LL}\), \(g^{\rm F}_{\rm LL}\) & the gauge couplings of
the isospin vector field and of the vector field of the fermion
number \\
\(k^{\rm B}_{\rm LL}\), \(\omega^{\rm B}_{\rm LL}\) & the wave number
and the frequency of a plane wave of a massive vector field \\
\(\tau^{\rm LL}_\iota\), \(\iota=1,2,3\) & the Pauli matrices acting on
the isospin index \(\iota\) (the same matrices as
\(\sigma^{\rm LL}_j\)) \\
\(T^\iota_{\rm LL}\) & the isospin generators
\(T^\iota_{\rm LL}=\tau^{\rm LL}_\iota/2\);
\((T_{\rm LL})^2=\sum_\iota T^\iota_{\rm LL}T^\iota_{\rm LL}\) is the
isospin Casimir \\
\({\rm arsinh}\) & the inverse hyperbolic sine,
\({\rm arsinh}\,y=\ln\big(y+\sqrt{y^2+1}\big)\) \\
\end{longtable}

\hypertarget{words-used-for-the-registry}{%
\subsection{2.3 Words used for the
registry}\label{words-used-for-the-registry}}

\begin{itemize}
\tightlist
\item
  \textbf{Registry}: the table of all symbols of the super-loop revision
  of the project; every row is an \textbf{entry} with a number, a
  \textbf{group} (one of nine topics), the \textbf{printed forms} of its
  symbols written in TeX between dollar signs, the \textbf{plain-text
  names} under which programs and data files write the symbols, the
  \textbf{meaning}, and the \textbf{origin}.
\item
  \textbf{TeX item}: one single symbol written in TeX, for example
  \texttt{\textbackslash{}Delta\^{}\{\textbackslash{}rm\ vac\}\_\{\textbackslash{}rm\ LL\}};
  an entry may list several TeX items.
\item
  \textbf{Kind} of a TeX item: symbol, family (a symbol with a running
  index), function (a symbol followed by its argument in parentheses),
  function family, letter, index, primed index, set, group, ket,
  operator, expression, imaginary unit (section 7).
\item
  \textbf{Label entry}: an entry whose symbol field contains no TeX
  (names of checks and hypotheses).
\end{itemize}

\hypertarget{further-symbols-named-in-this-document-as-examples}{%
\subsection{2.4 Further symbols named in this document as
examples}\label{further-symbols-named-in-this-document-as-examples}}

These symbols appear below only as examples of the forms of symbols;
their full meanings are in the registry inside the notebook.

\begin{longtable}[]{@{}
  >{\raggedright\arraybackslash}p{(\columnwidth - 2\tabcolsep) * \real{0.5000}}
  >{\raggedright\arraybackslash}p{(\columnwidth - 2\tabcolsep) * \real{0.5000}}@{}}
\toprule\noalign{}
\begin{minipage}[b]{\linewidth}\raggedright
symbol
\end{minipage} & \begin{minipage}[b]{\linewidth}\raggedright
meaning
\end{minipage} \\
\midrule\noalign{}
\endhead
\bottomrule\noalign{}
\endlastfoot
\(\mathbb R\), \(\mathbb C\), \(\mathbb H\), \(\mathbb O\) & the real
numbers, the complex numbers, the quaternions, the octonions (an
eight-dimensional real algebra) \\
\(\mathbb B_{\rm LL}\), bi-octonion & the complex octonions
\(\mathbb C\otimes\mathbb O\); an element is a bi-octonion, written in
bold, for example \(\mathbf x\), \(\mathbf y\), \(\mathbf u\),
\(\mathbf b\) \\
\(\mathbf e_0\) & the unit element 1 of \(\mathbb B_{\rm LL}\) (the
first of the eight units \(\mathbf e_0,\dots,\mathbf e_7\)) \\
\(D\) & the number of space dimensions \\
\(p\), \(p'\); \(j\); \(\iota\), \(\iota'\), \(\iota''\) & indices:
\(p=0,1,2,3\) over the directions of space-time (and over a basis of an
algebra), \(j=1,2,3\) over the directions of space, \(\iota=1,2,3\) over
the isospin; a primed letter is a second or third index of the same
kind \\
\(\partial_p\) & the partial derivative with respect to
\(x^p_{\rm LL}\) \\
\(\nabla^{\rm LL}_p\), \(\nabla^{{\rm LL}*}_p\) & the gauge-covariant
derivative of the spinor field, and the covariant derivative of a field
that transforms like the complex-conjugate spinor field \\
\(\nabla^{\rm B}_{{\rm LL},p}\) & the gauge-covariant derivative of the
field of the superfluid boson singlet state \\
\(A^\iota_{{\rm LL},p}\) & the gauge potentials of the isospin group,
three real fields for each direction \(p\) \\
\(\mathsf N^{\rm g}_{\rm LL}\) & the contact matrix, a real symmetric
\(4\times4\) matrix of the terms of the kinetic Lagrangian that are
quadratic in the gauge potentials (an example of a sans-serif symbol) \\
\(W_{\rm LL}(\mathbf u)\) & a function symbol: the linear map
\(W_{\rm LL}(\mathbf u)\mathbf b=-\mathbf b\,\mathbf u\) that depends on
the argument \(\mathbf u\) \\
\({\rm SU}(2)_{\rm v}\), \(\mathfrak{so}(3,1)\) & the isospin group; the
Lie algebra of the Lorentz group of \(3+1\) dimensions \\
\(\lvert\Phi^{\rm SF}_{\rm LL}\rangle\),
\(\lvert\Phi^{\rm N}_{\rm LL}\rangle\) & the superfluid boson singlet
state (the ground state of the spin-zero superfluid) and the normal
state from which it is built \\
\({\rm tr}\), \(\det\), \({\rm diag}\), \({\rm Re}\), \({\rm Im}\),
\({}^T\) & trace, determinant, diagonal matrix, real part, imaginary
part, transpose \\
\(\otimes\), \(\oplus\), \(\circ\) & tensor product, direct sum,
composition of maps \\
\([\mathbf x,\mathbf y]\), \([\mathbf x,\mathbf y,\mathbf z]\) & the
commutator \(\mathbf x\mathbf y-\mathbf y\mathbf x\) and the associator
\((\mathbf x\mathbf y)\mathbf z-\mathbf x(\mathbf y\mathbf z)\) \\
\([\ \cdot\ ]\), \(\langle\ \rangle\) & the mass dimension of a field or
coupling; the expectation value in the ground state \\
\(a_{\rm LL}\), \(a_{{\rm LL},p}\), \(b_{{\rm LL},p}\) & symbols of the
demonstrations of section I.2 of the notebook (not part of the
registry) \\
\end{longtable}

\hypertarget{symbols-of-part-ii-item-1}{%
\subsection{2.5 Symbols of Part II, item
{[}1{]}}\label{symbols-of-part-ii-item-1}}

Every symbol in this table is a registered symbol of the project (Part I
of the notebook symbolizes all of them; the registry is described in
section 7). The meanings are those of the registry, written out here in
full, with the formulas used by the programs. The name in the last
column is the function or variable of Part II of the notebook that
computes the object (code font). Indices: \(p,p'=0,1,2,3\) (space-time
directions), \(j,j',j''=1,2,3\) (space directions),
\(\iota,\iota',\iota''=1,2,3\) (isospin directions), \(k=1,\dots,8\)
(the components of the spinor field), \(n\) (the number of a mode).

\begin{longtable}[]{@{}
  >{\raggedright\arraybackslash}p{(\columnwidth - 4\tabcolsep) * \real{0.3333}}
  >{\raggedright\arraybackslash}p{(\columnwidth - 4\tabcolsep) * \real{0.3333}}
  >{\raggedright\arraybackslash}p{(\columnwidth - 4\tabcolsep) * \real{0.3333}}@{}}
\toprule\noalign{}
\begin{minipage}[b]{\linewidth}\raggedright
symbol
\end{minipage} & \begin{minipage}[b]{\linewidth}\raggedright
meaning
\end{minipage} & \begin{minipage}[b]{\linewidth}\raggedright
in Part II
\end{minipage} \\
\midrule\noalign{}
\endhead
\bottomrule\noalign{}
\endlastfoot
\(\sigma^{\rm LL}_0\), \(\sigma^{\rm LL}_j\) & the \(2\times2\) unit
matrix and the Pauli matrices of section 2.2 & \texttt{t2sigma{[}0{]}}
\ldots{} \texttt{t2sigma{[}3{]}} \\
\(\tau^{\rm LL}_\iota\) & the same Pauli matrices acting on the isospin
index (isospin up, down) & \texttt{t2sigma{[}i{]}} \\
\(\mathsf E_{\rm LL}\) & the real antisymmetric \(2\times2\) matrix
\(i\sigma^{\rm LL}_2=\begin{pmatrix}0&1\\-1&0\end{pmatrix}\) &
\texttt{t2eps2} \\
\(\mathsf 1_2\), \(\mathsf 1_4\), \(\mathsf 1_8\), \(\mathsf 1_{16}\) &
the unit matrices of the dimensions 2, 4, 8, 16 &
\texttt{IdentityMatrix} \\
\(\epsilon^{\rm LL}_{\iota\iota'\iota''}\) & the Levi-Civita symbol of
section 2.2, also for isospin indices & \texttt{t2eps3} \\
\(\eta^{\rm LL}_{pp'}\) & the Minkowski metric of section 2.2 &
\texttt{t2eta{[}p{]}} (its diagonal) \\
\(\mathbb B_{\rm LL}\); \(\mathbf e_0,\dots,\mathbf e_7\) & the
bi-octonions, the complex combinations of the eight units
\(\mathbf e_0=1,\mathbf e_1,\dots,\mathbf e_7\) & \\
\(\mathbf a\), \(\mathbf b\), \(\mathbf c\), \(\mathbf d\);
\(\mathbf x\), \(\mathbf y\) & bi-quaternions, the complex combinations
of \(\mathbf e_0,\dots,\mathbf e_3\) (a bi-quaternion is also a
bi-octonion, so the letters \(\mathbf x\), \(\mathbf y\) of bi-octonions
are used for bi-quaternions as well); a bi-octonion is written
\(\mathbf a+\mathbf b\,\mathbf e_4\) (Cayley-Dickson form) & lists of 4
complex numbers \\
\(\xi^{\rm LL}_p(\mathbf b)\) & the complex coefficient of
\(\mathbf e_p\) in \(\mathbf b\):
\(\mathbf b=\sum_{p=0}^3\xi^{\rm LL}_p(\mathbf b)\mathbf e_p\) & the
entries of the list \\
\(\bar{\mathbf b}\), \(\mathbf b^*\), \(\mathbf b^\dagger\) & the
conjugations: \(\bar{\mathbf b}\) changes the signs of the coefficients
of \(\mathbf e_1,\mathbf e_2,\mathbf e_3\); \(\mathbf b^*\) takes the
complex conjugate of every coefficient;
\(\mathbf b^\dagger=(\bar{\mathbf b})^*\) & \texttt{t2bqConj},
\texttt{t2bqStar}, \texttt{t2bqDag} \\
\(M_{\rm LL}(\mathbf b)\) & the \(2\times2\) matrix of a bi-quaternion,
\(M_{\rm LL}(\mathbf b)=\xi^{\rm LL}_0(\mathbf b)\sigma^{\rm LL}_0-i\sum_j\xi^{\rm LL}_j(\mathbf b)\sigma^{\rm LL}_j\),
so that \(M_{\rm LL}(\mathbf e_0)=\sigma^{\rm LL}_0\),
\(M_{\rm LL}(\mathbf e_j)=-i\sigma^{\rm LL}_j\); it is multiplicative,
\(M_{\rm LL}(\mathbf b\mathbf b')=M_{\rm LL}(\mathbf b)M_{\rm LL}(\mathbf b')\),
with \(M_{\rm LL}(\mathbf b^\dagger)=M_{\rm LL}(\mathbf b)^\dagger\) and
\(\det M_{\rm LL}(\mathbf b)=N_{\rm LL}(\mathbf b)\) & \texttt{t2mq};
inverse \texttt{t2bqFrom}; product \texttt{t2bqMul} \\
\(N_{\rm LL}(\mathbf x)\) & the complex norm,
\(N_{\rm LL}(\mathbf x)=\sum_p\xi^{\rm LL}_p(\mathbf x)^2\) & \\
\(R^{\rm LL}_{\mathbf x}\) & the right multiplication
\(\mathbf y\mapsto\mathbf y\mathbf x\) & \\
\(\mathbb R^{1,3}_{\rm LL}\); \(x^p_{\rm LL}\); \(x_{\rm LL}\);
\(\partial_p\) & Minkowski space of \(3+1\) dimensions, the Hermitian
bi-quaternions
\(x^0_{\rm LL}\mathbf e_0+i\sum_jx^j_{\rm LL}\mathbf e_j\) with real
coordinates \(x^p_{\rm LL}\); a point
\(x_{\rm LL}=(x^0_{\rm LL},\dots,x^3_{\rm LL})\); the partial derivative
with respect to \(x^p_{\rm LL}\). The coordinates of a Hermitian
bi-quaternion \(\mathbf x\) are
\(x^p_{\rm LL}=\tfrac12{\rm tr}\big(\sigma^{\rm LL}_pM_{\rm LL}(\mathbf x)\big)\)
& \texttt{t2coords} \\
\(\mathbf f_p\) & a frame: four Hermitian bi-quaternions
\(\mathbf f_0,\dots,\mathbf f_3\) (with
\(N_{\rm LL}(\mathbf f_p,\mathbf f_{p'})=\eta^{\rm LL}_{pp'}\) for an
orthonormal frame); the canonical frame is \(\mathbf f_0=\mathbf e_0\),
\(\mathbf f_j=i\,\mathbf e_j\), with
\(M_{\rm LL}(\mathbf f_p)=\sigma^{\rm LL}_p\),
\(M_{\rm LL}(\bar{\mathbf f}_0)=\sigma^{\rm LL}_0\),
\(M_{\rm LL}(\bar{\mathbf f}_j)=-\sigma^{\rm LL}_j\) &
\texttt{t2canonicalFrame} \\
\(\gamma^p_{\rm LL}[\mathbf f]\) & the Dirac matrices of a frame,
\(\gamma^p_{\rm LL}[\mathbf f]=\begin{pmatrix}0&M_{\rm LL}(\mathbf f_p)\\M_{\rm LL}(\bar{\mathbf f}_p)&0\end{pmatrix}\)
(\(4\times4\), blocks of \(2\times2\)) & \texttt{t2gammasFromFrame} \\
\(\gamma^p_{\rm LL}\), \(\gamma^{\rm ch}_{\rm LL}\), \(C_{\rm LL}\),
\(\alpha^j_{\rm LL}\), \(\beta_{\rm LL}\) & the Dirac matrices of the
canonical frame (chiral basis),
\(\gamma^0_{\rm LL}=\begin{pmatrix}0&\sigma^{\rm LL}_0\\\sigma^{\rm LL}_0&0\end{pmatrix}\),
\(\gamma^j_{\rm LL}=\begin{pmatrix}0&\sigma^{\rm LL}_j\\-\sigma^{\rm LL}_j&0\end{pmatrix}\);
the chirality matrix
\(\gamma^{\rm ch}_{\rm LL}=i\gamma^0_{\rm LL}\gamma^1_{\rm LL}\gamma^2_{\rm LL}\gamma^3_{\rm LL}=\mathrm{diag}(-1,-1,1,1)\);
the charge-conjugation matrix
\(C_{\rm LL}=i\gamma^2_{\rm LL}\gamma^0_{\rm LL}\);
\(\alpha^j_{\rm LL}=\gamma^0_{\rm LL}\gamma^j_{\rm LL}\);
\(\beta_{\rm LL}=\gamma^0_{\rm LL}\) & \texttt{t2gamma{[}p{]}},
\texttt{t2gammaCh}, \texttt{t2cconj}, \texttt{t2alpha{[}j{]}},
\texttt{t2beta} \\
\(\mathfrak l_{\rm LL}\); \(J^j_{\rm LL}\), \(K^j_{\rm LL}\);
\(\mathbf u\), \(\mathbf u_1,\dots,\mathbf u_5\) & the Lorentz part of
the Malcev algebra of the super-loop (the imaginary bi-quaternions); its
rotation and boost generators \(J^j_{\rm LL}=\mathbf e_j/2\),
\(K^j_{\rm LL}=i\,\mathbf e_j/2\); elements of \(\mathfrak l_{\rm LL}\)
(Lorentz generators; the registry lists \(\mathbf u\), \(\mathbf u'\),
\(\mathbf u_1\), \(\mathbf u_2\), and GS-15 numbers five random
generators in the same way) & \texttt{t2jgen{[}j{]}},
\texttt{t2kgen{[}j{]}} \\
\(\mathsf s_{\rm LL}(\mathbf u)\) & a Lorentz generator acting on the
Dirac spinor,
\(\mathsf s_{\rm LL}(\mathbf u)=\mathrm{diag}\big(-M_{\rm LL}(\mathbf u)^\dagger,M_{\rm LL}(\mathbf u)\big)\)
& \texttt{t2spinGen} \\
\(\mathbf g\); \(\tilde{\mathbf g}\) & a unit bi-quaternion
(\(N_{\rm LL}(\mathbf g)=1\)), an element of \(SL(2,\mathbb C)\);
\(\tilde{\mathbf g}=(\mathbf g^\dagger)^{-1}=\mathbf g^*\), with
\(M_{\rm LL}(\tilde{\mathbf g})=(M_{\rm LL}(\mathbf g)^\dagger)^{-1}\) &
\texttt{t2randSL2C} (a random one) \\
\(\mathsf S_{\rm LL}(\mathbf g)\) & the spinor matrix of the Lorentz
transformation \(\mathbf g\),
\(\mathrm{diag}\big(M_{\rm LL}(\tilde{\mathbf g}),M_{\rm LL}(\mathbf g)\big)\)
& \texttt{t2spinRep} \\
\(\mathsf L_{\rm LL}(\mathbf g)\) & the real \(4\times4\) Lorentz matrix
of \(\mathbf x\mapsto\mathbf g\mathbf x\mathbf g^\dagger\):
\([\mathsf L_{\rm LL}(\mathbf g)]_{pp'}=\tfrac12{\rm Re}\,{\rm tr}\big(\sigma^{\rm LL}_pM_{\rm LL}(\mathbf g)\sigma^{\rm LL}_{p'}M_{\rm LL}(\mathbf g)^\dagger\big)\)
& \texttt{t2lmat} \\
\(\mathcal V_{\rm LL}(\mathbf u)\) & the vector action of a Lorentz
generator,
\(\mathbf x\mapsto\mathbf u\mathbf x+\mathbf x\mathbf u^\dagger\), a
real \(4\times4\) matrix with the entries
\(\mathcal V_{\rm LL}(\mathbf u)^p{}_{p'}\) & \texttt{t2vecAction} \\
\(\Psi_{\rm LL}\); \(\Psi_{{\rm LL},k}\) & the Dirac isodoublet of the
spinor charges, 8 components: components 1 and 2 are the first Dirac
component with isospin up and down, 3 and 4 the second, 5 and 6 the
third, 7 and 8 the fourth; a \(4\times4\) Dirac matrix acts as its
Kronecker product with \(\sigma^{\rm LL}_0\) on the right, a
\(2\times2\) isospin matrix as the Kronecker product of \(\mathsf 1_4\)
with it & \texttt{t2lift4}, \texttt{t2lift2} \\
\(T^\iota_{\rm LL}\) & the isospin generators on the eight components,
\(\mathsf 1_4\otimes\tau^{\rm LL}_\iota/2\) & \texttt{t2tiso{[}i{]}} \\
\(\mathsf K^\iota_{\rm LL}\), \(\mathsf K'^\iota_{\rm LL}\) & the
\(8\times8\) pair matrices of the spin-zero channels: scalar
\(\mathsf K^\iota_{\rm LL}=C_{\rm LL}\gamma^{\rm ch}_{\rm LL}\otimes\mathsf E_{\rm LL}\tau^{\rm LL}_\iota\),
pseudoscalar
\(\mathsf K'^\iota_{\rm LL}=C_{\rm LL}\otimes\mathsf E_{\rm LL}\tau^{\rm LL}_\iota\)
& \texttt{t2kpair{[}i{]}}, \texttt{t2kpairPs{[}i{]}} \\
\(\vec d_{\rm LL}\), \(d^\iota_{\rm LL}\); \(\Delta_{\rm LL}\) & the
order parameter, a complex isospin vector
\((d^1_{\rm LL},d^2_{\rm LL},d^3_{\rm LL})\); the gap
\(\Delta_{\rm LL}=\lvert\vec d_{\rm LL}\rvert\) of a unitary state & \\
\(\mathsf P_{\rm LL}\) & the pairing block
\(\sum_\iota d^\iota_{\rm LL}(\mathsf K^\iota_{\rm LL})^\dagger\) &
\texttt{t2pairPotential} \\
\(\mathsf U_{\rm LL}\); \(\theta^\iota_{\rm LL}\);
\(\mathsf R^{\rm v}_{\rm LL}(\mathsf U_{\rm LL})\); \(\mathbf v\) & an
isospin rotation, a matrix of \(SU(2)\),
\(\mathsf U_{\rm LL}=\exp\big(i\sum_\iota\theta^\iota_{\rm LL}\tau^{\rm LL}_\iota/2\big)\)
with the three angles \(\theta^\iota_{\rm LL}\); the rotation of an
isospin vector,
\(\mathsf U_{\rm LL}\tau^{\rm LL}_\iota\mathsf U_{\rm LL}^\dagger=\sum_{\iota'}[\mathsf R^{\rm v}_{\rm LL}(\mathsf U_{\rm LL})]_{\iota'\iota}\tau^{\rm LL}_{\iota'}\);
the unit real quaternion with
\(\mathsf U_{\rm LL}=M_{\rm LL}(\mathbf v)\) & \texttt{t2su2},
\texttt{t2adjoint} \\
\(\theta^{\rm F}_{\rm LL}\) & the phase of a transformation of the
fermion number,
\(\Psi_{\rm LL}\mapsto\exp(i\theta^{\rm F}_{\rm LL})\Psi_{\rm LL}\) & \\
\(\mathcal U_{\rm LL}\) & a local gauge transformation by
\(SL(2,\mathbb C)\times SU(2)_{\rm v}\times U(1)^{\rm F}_{\rm LL}\),
\(\mathcal U_{\rm LL}=\exp(i\theta^{\rm F}_{\rm LL})\,\mathsf S_{\rm LL}(\mathbf g)\otimes\mathsf U_{\rm LL}\),
with \(\mathbf g\), \(\mathsf U_{\rm LL}\), \(\theta^{\rm F}_{\rm LL}\)
functions of \(x_{\rm LL}\) & \texttt{caluAt} \\
\(A^\iota_{{\rm LL},p}\), \(A^{\rm F}_{{\rm LL},p}\);
\(g^{\rm v}_{\rm LL}\), \(g^{\rm F}_{\rm LL}\) & the gauge potentials of
the isospin group and of the fermion number (real fields); their
couplings & \texttt{aAt}, \texttt{afAt}; \texttt{gv}, \texttt{gf} \\
\(\mathsf a_{{\rm LL},p}\); \(\vec A^{\rm v}_{{\rm LL},p}\);
\(\mathbf t^{\rm A}_{{\rm LL},p}\) & the isospin potential as a
Hermitian \(2\times2\) matrix,
\(\sum_\iota A^\iota_{{\rm LL},p}\tau^{\rm LL}_\iota/2\); as a vector
\((A^1_{{\rm LL},p},A^2_{{\rm LL},p},A^3_{{\rm LL},p})\); as a real
imaginary quaternion,
\(\sum_\iota A^\iota_{{\rm LL},p}\mathbf e_\iota/2\), with
\(M_{\rm LL}(\mathbf t^{\rm A}_{{\rm LL},p})=-i\,\mathsf a_{{\rm LL},p}\)
& \\
\(\mathbf u^\omega_{{\rm LL},p}\); \(\omega^j_{{\rm LL},p}\),
\(\varpi^j_{{\rm LL},p}\) & the Lorentz connection,
\(\mathbf u^\omega_{{\rm LL},p}=\sum_j\big(\omega^j_{{\rm LL},p}J^j_{\rm LL}+\varpi^j_{{\rm LL},p}K^j_{\rm LL}\big)\),
with real rotation and boost components;
\(\mathbf u^{\omega*}_{{\rm LL},p}\) is its complex conjugate &
\texttt{lorentzElem} \\
\(\Gamma^{\rm LL}_p\) & the spin connection,
\(\Gamma^{\rm LL}_p=\mathsf s_{\rm LL}(\mathbf u^\omega_{{\rm LL},p})\otimes\sigma^{\rm LL}_0\)
& \texttt{gammaAt} (its \(4\times4\) part) \\
\(\mathsf A_{{\rm LL},p}\) & the gauge matrix of the internal
symmetries,
\(g^{\rm v}_{\rm LL}\sum_\iota A^\iota_{{\rm LL},p}T^\iota_{\rm LL}+g^{\rm F}_{\rm LL}A^{\rm F}_{{\rm LL},p}\mathsf 1_8\)
(Hermitian) & \texttt{gauge8} \\
\(\nabla^{\rm LL}_p\), \(\nabla^{{\rm LL}*}_p\) & the gauge-covariant
derivative \(\partial_p+\Gamma^{\rm LL}_p-i\mathsf A_{{\rm LL},p}\), and
that of a field that transforms like \(\Psi^*_{\rm LL}\),
\(\partial_p+(\Gamma^{\rm LL}_p)^*+i(\mathsf A_{{\rm LL},p})^*\) &
\texttt{nabla} \\
\(\mathbf D_{\rm LL}\) & the bi-quaternion of covariant derivatives
\(\sum_p\mathbf f_p\nabla^{\rm LL}_p\);
\(\sum_p\gamma^p_{\rm LL}\nabla^{\rm LL}_p=\begin{pmatrix}0&M_{\rm LL}(\mathbf D_{\rm LL})\\M_{\rm LL}(\bar{\mathbf D}_{\rm LL})&0\end{pmatrix}\)
& \\
\(\mathsf F^\iota_{{\rm LL},pp'}\),
\(\mathsf F^{\rm F}_{{\rm LL},pp'}\), \(\mathsf R^{\rm LL}_{pp'}\) & the
field strengths and the curvature:
\(\mathsf F^\iota_{{\rm LL},pp'}=\partial_pA^\iota_{{\rm LL},p'}-\partial_{p'}A^\iota_{{\rm LL},p}+g^{\rm v}_{\rm LL}\sum_{\iota',\iota''}\epsilon^{\rm LL}_{\iota\iota'\iota''}A^{\iota'}_{{\rm LL},p}A^{\iota''}_{{\rm LL},p'}\),
\(\mathsf F^{\rm F}_{{\rm LL},pp'}=\partial_pA^{\rm F}_{{\rm LL},p'}-\partial_{p'}A^{\rm F}_{{\rm LL},p}\),
\(\mathsf R^{\rm LL}_{pp'}=\partial_p\Gamma^{\rm LL}_{p'}-\partial_{p'}\Gamma^{\rm LL}_p+\Gamma^{\rm LL}_p\Gamma^{\rm LL}_{p'}-\Gamma^{\rm LL}_{p'}\Gamma^{\rm LL}_p\)
& \\
\(\psi'_{\rm LL}(\mathbf y)\) & the Weyl isodoublet of a spinor charge
\(\mathbf y\mathbf e_4\), the \(2\times2\) matrix
\(\mathsf E_{\rm LL}^{-1}M_{\rm LL}(\mathbf y)^T\) (rows: Weyl index;
columns: isospin up, down) & \texttt{t2psiPrime} \\
\(\mathbf b^+_{\rm LL}\), \(\mathbf b^-_{\rm LL}\);
\(\Psi_{\rm LL}[\mathbf b^-_{\rm LL},\mathbf b^+_{\rm LL}]\) & the two
spinor charges \(\mathbf b^\pm_{\rm LL}\mathbf e_4\) of the Dirac
isodoublet: its upper block (components 1 to 4, chirality \(-1\)) is
\(\psi'_{\rm LL}(\mathbf b^-_{\rm LL})\), its lower block (components 5
to 8, chirality \(+1\)) is \(\psi'_{\rm LL}(\mathbf b^+_{\rm LL})\) &
\texttt{t2diracFromCharges} \\
\(\psi^{\rm p}_{\rm LL}\), \(\psi^{\rm h}_{\rm LL}\);
\(\psi^{\rm p}_{{\rm LL},k}\) & the particle and the hole component of a
mode of the mean-field equations, eight complex functions each;
\(\psi^{\rm p}_{\rm LL}\) transforms like \(\Psi_{\rm LL}\),
\(\psi^{\rm h}_{\rm LL}\) like \(\Psi^*_{\rm LL}\); a component & \\
\(\vec q_{\rm LL}\), \(q_{\rm LL}\), \(q^j_{\rm LL}\); \(m_{\rm LL}\),
\(\mu_{\rm LL}\), \(G_{\rm LL}\) & a momentum, its length and its
components; the mass of the spinor charges, the chemical potential, the
coupling of the contact interaction & \texttt{q}; \texttt{mass},
\texttt{mu} \\
\(\varepsilon_{\rm LL}\), \(E^\pm_{\rm LL}\) & the energy
\((q_{\rm LL}^2+m_{\rm LL}^2)^{1/2}\); the quasi-particle energies of
the scalar channel,
\(E^\pm_{\rm LL}=\big((\varepsilon_{\rm LL}\pm\mu_{\rm LL})^2+\Delta_{\rm LL}^2\big)^{1/2}\)
& \\
\(\mathsf h_{\rm LL}(\vec q_{\rm LL})\);
\(\mathsf H_{\rm LL}(\vec q_{\rm LL})\) & the Hamiltonian matrix of a
plane wave,
\(\sum_j\alpha^j_{\rm LL}q^j_{\rm LL}+\beta_{\rm LL}m_{\rm LL}\); the
\(16\times16\) Bogoliubov-de Gennes matrix of section 12.1 without gauge
fields & \\
\(\mathcal H^{\rm DB}_{\rm LL}\), \(h^\nabla_{\rm LL}\);
\(\mathcal H^{\rm DB}_{\rm LL}(\vec q_{\rm LL})\) & the
Dirac-Bogoliubov-de Gennes operator (section 12.1) and the covariant
Dirac Hamiltonian
\(h^\nabla_{\rm LL}=-i\sum_j\alpha^j_{\rm LL}\nabla^{\rm LL}_j+\beta_{\rm LL}m_{\rm LL}\);
its \(16\times16\) matrix on a plane wave with constant fields &
\texttt{hq} \\
\(E^{\rm qp}_{{\rm LL},n}\) & the energy of a numbered mode, an
eigenvalue of \(\mathcal H^{\rm DB}_{\rm LL}\) & \\
\(\Pi^\iota_{\rm LL}\), \(\Pi^\iota_{\rm LL}(x_{\rm LL})\) & the pair
operators, \(\Psi_{\rm LL}^T\mathsf K^\iota_{\rm LL}\Psi_{\rm LL}\) (at
a point) & \texttt{t2anomalous} (lattice) \\
\(\ell^{\rm x}_{\rm LL}\), \(\ell^{\rm t}_{\rm LL}\),
\(r^{\rm W}_{\rm LL}\) & the lattice spacing, the time step, the
coefficient of the Wilson term & \texttt{"ell"}, \texttt{dt},
\texttt{"wilson"} \\
\(n^{\rm x}_{\rm LL}\), \(n^{\rm y}_{\rm LL}\), \(n^{\rm z}_{\rm LL}\);
\(\vec s_{\rm LL}\), \(s^j_{\rm LL}\); \(\vec\delta^j_{\rm LL}\);
\(n^{\rm act}_{\rm LL}\) & the numbers of lattice sites in the three
directions; a site and its integer coordinates (counted from 0); the
step to the next site in the direction \(j\); the number of directions
with more than one site & \texttt{"n"}; \texttt{t2siteCoords} \\
\(\mathsf U^{\rm link}_{{\rm LL},j}(\vec s_{\rm LL})\) & the link matrix
from \(\vec s_{\rm LL}\) to \(\vec s_{\rm LL}+\vec\delta^j_{\rm LL}\),
\(\exp\big(-i\ell^{\rm x}_{\rm LL}(g^{\rm v}_{\rm LL}\mathsf a_{{\rm LL},j}+g^{\rm F}_{\rm LL}A^{\rm F}_{{\rm LL},j}\sigma^{\rm LL}_0)\big)\)
& \texttt{t2link} \\
\(\mathcal H^{\rm lat}_{\rm LL}\), \(\varepsilon^{\rm lat}_{\rm LL}\),
\(B^{\rm H}_{\rm LL}\) & the lattice form of
\(\mathcal H^{\rm DB}_{\rm LL}\) (16 rows per site), the energy of a
free particle on the lattice, and the bound on the largest eigenvalue
magnitude of \(\mathcal H^{\rm lat}_{\rm LL}\) & \texttt{t2apply},
\texttt{t2dense}; \texttt{t2epsLat}; \texttt{t2normBound} \\
\(\mathsf h^{\rm lat}_{\rm LL}(\vec q_{\rm LL})\) & the \(4\times4\)
free lattice operator on a plane wave,
\(\sum_j\alpha^j_{\rm LL}\sin(q^j_{\rm LL}\ell^{\rm x}_{\rm LL})/\ell^{\rm x}_{\rm LL}+\beta_{\rm LL}\big(m_{\rm LL}+(r^{\rm W}_{\rm LL}/\ell^{\rm x}_{\rm LL})\sum_j(1-\cos(q^j_{\rm LL}\ell^{\rm x}_{\rm LL}))\big)\),
sums over the active directions & \texttt{t2hLatPlane} \\
\(\mathsf G^{\rm site}_{\rm LL}\), \(\mathsf X^{\rm ph}_{\rm LL}\) & the
gauge transformation of the lattice equations,
\(\exp(i\theta^{\rm F}_{\rm LL})\mathsf U_{\rm LL}\) at each site on the
isospin index of the particle component and its complex conjugate on the
hole component; the exchange of the particle and the hole component,
\(\begin{pmatrix}0&\mathsf 1_8\\\mathsf 1_8&0\end{pmatrix}\) &
\texttt{gbig}; \texttt{swap}, \texttt{xm} \\
\(SL(2,\mathbb C)\), \(SU(2)_{\rm v}\), \(U(1)^{\rm F}_{\rm LL}\),
\(\mathfrak{so}(3,1)\) & the Lorentz group (unit bi-quaternions), the
isospin group, the group of the fermion number, the Lie algebra of the
Lorentz group & \\
\end{longtable}

\hypertarget{abbreviations-of-this-document-that-are-not-registered-symbols}{%
\subsection{2.6 Abbreviations of this document that are not registered
symbols}\label{abbreviations-of-this-document-that-are-not-registered-symbols}}

Two letters are used, in a few places only, as abbreviations; each is
defined where it is used and none of them coincides with a registered
symbol (the registry has no symbol written with these letters alone):

\begin{itemize}
\tightlist
\item
  \(z\) (section 12.2, checks GS-13 and T2-DATA-1, section 12.3, and
  Figure 2 of the notebook): the product
  \(E^{\rm qp}_{{\rm LL},n}\ell^{\rm t}_{\rm LL}\) of the energy of a
  mode and the time step.
\item
  \(\kappa\) (sections 12.2 GS-10 and GS-16, and the notebook text of
  GS-16): the integer that labels the allowed momenta of a periodic
  chain:
  \(q^1_{\rm LL}=2\pi\kappa/(n^{\rm x}_{\rm LL}\ell^{\rm x}_{\rm LL})\)
  for a chain of \(n^{\rm x}_{\rm LL}\) sites along \(x^1_{\rm LL}\),
  and
  \(q^3_{\rm LL}=2\pi\kappa/(n^{\rm z}_{\rm LL}\ell^{\rm x}_{\rm LL})\)
  for a chain along \(x^3_{\rm LL}\).
\item
  Variables of the programs are written in code font (for example
  \texttt{kv}, \texttt{waves}, \texttt{ulist}, \texttt{w}); they are
  names of the Wolfram Language code, not symbols of the physics.
\end{itemize}

\hypertarget{results-in-brief}{%
\section{3. Results in brief}\label{results-in-brief}}

\begin{longtable}[]{@{}
  >{\raggedright\arraybackslash}p{(\columnwidth - 2\tabcolsep) * \real{0.5000}}
  >{\raggedright\arraybackslash}p{(\columnwidth - 2\tabcolsep) * \real{0.5000}}@{}}
\toprule\noalign{}
\begin{minipage}[b]{\linewidth}\raggedright
quantity
\end{minipage} & \begin{minipage}[b]{\linewidth}\raggedright
value
\end{minipage} \\
\midrule\noalign{}
\endhead
\bottomrule\noalign{}
\endlastfoot
entries of the registry & 199, in 9 groups \\
TeX segments (text between two dollar signs) & 434 \\
TeX items (single symbols) & 449 \\
label entries (no TeX) & 2 (entries 86 and 87) \\
TeX items of a symbolized kind & 427 \\
different boxes, i.e.~symbolized symbols & 422 \\
tests of Part I & 15 (S-0 to S-14), all passed \\
original programs rewritten in Part II so far & the library t2core
(rng.rs, c64.rs, cmat.rs, model.rs, lattice.rs) and the program
t2\_fields (item {[}1{]} of Test 2) \\
checks of Part II so far & 20 (T2-SYM, T2-RNG, T2-LINALG, GS-1 to GS-16,
T2-DATA-1), all passed \\
components of families typed in the notebook & 10, all atomic symbols
(T2-SYM) \\
messages during the evaluation & 0 \\
evaluation time (Wolfram 15.0.1, Mac, Apple silicon) & 71 s \\
printed form of the evaluated notebook & 181 pages \\
\end{longtable}

\hypertarget{how-to-use-the-notebook}{%
\section{4. How to use the notebook}\label{how-to-use-the-notebook}}

\begin{enumerate}
\def\labelenumi{\arabic{enumi}.}
\tightlist
\item
  Open \texttt{3+1super-loop.nb} in Wolfram (version 15 or later; the
  notebook was evaluated with 15.0.1).
\item
  Choose Evaluation \textgreater{} Evaluate Notebook. Every cell is
  evaluated in order. Nothing outside the notebook is read: the
  registry, the converter and all tests are inside it.
\item
  Every test prints one line that begins with PASS (green) or FAIL
  (red), its name, the statement that it checks, and the numbers that
  decide it. Section I.9 of the notebook lists the tests of Part I in a
  table, and its last cell gives
  \texttt{\{"tests"\ -\textgreater{}\ 15,\ "passed"\ -\textgreater{}\ 15\}};
  the subsection ``Summary of item {[}1{]}'' at the end of chapter II.2
  lists the 20 checks of Part II. In Part II the numbers of a check are
  followed, in parentheses, by those of the original report.
\item
  After the evaluation every registered symbol can be typed in its
  printed form; \texttt{?symbol} prints its meaning.
\end{enumerate}

The notebook as stored in the repository is already evaluated, so its
results can be read without evaluating it. The folder
\texttt{3+1super-loop\_provenance/export} holds the printed form of the
evaluated notebook as a PDF file (181 pages; about 30 of them are the
input cell with the numbers of the data files of item {[}1{]}, decision
of section 15).

\hypertarget{the-technique-symbolized-symbols-created-by-a-program}{%
\section{5. The technique: symbolized symbols created by a
program}\label{the-technique-symbolized-symbols-created-by-a-program}}

\hypertarget{the-basic-call}{%
\subsection{5.1 The basic call}\label{the-basic-call}}

The author's note on creating a symbolized variable by a program
describes the technique that the notebook uses: the Notation package is
loaded, and \texttt{Symbolize} is called with the box of the symbol
wrapped in \texttt{ParsedBoxWrapper}:

\begin{verbatim}
Needs["Notation`"];
Symbolize[ParsedBoxWrapper[SubscriptBox["a", "LL"]], WorkingForm -> StandardForm];
\end{verbatim}

After this call the two-dimensional input \(a_{\rm LL}\) is one new
symbol (its internal name is
\texttt{a\textbackslash{}{[}UnderBracket{]}Subscript\textbackslash{}{[}UnderBracket{]}LL},
where \texttt{\textbackslash{}{[}UnderBracket{]}} is the separator
character that the package puts into names). The notebook checks:
\texttt{MakeExpression{[}SubscriptBox{[}"a",\ "LL"{]},\ StandardForm{]}}
gives \texttt{HoldComplete{[}a\_LL{]}} (displayed with the printed
form), its head is \texttt{Symbol} and it is an atom; the text input
\texttt{Subscript{[}a,\ LL{]}}, in contrast, is a compound expression
with head \texttt{Subscript}.

\hypertarget{the-working-form-must-be-given-explicitly}{%
\subsection{5.2 The working form must be given
explicitly}\label{the-working-form-must-be-given-explicitly}}

Without the option \texttt{WorkingForm\ -\textgreater{}\ StandardForm}
the package asks the front end for the format of the notebook in which
the input was typed (\texttt{InputNotebook{[}{]}}). When a notebook is
evaluated by a program, or when no notebook is selected, the answer is
\texttt{\$Failed}, and the package stores its reading rule for the
format \texttt{\$Failed}: the symbol is created but its box is never
recognized in input. This was found from the source code of the package
(the default value is taken from
\texttt{AbsoluteCurrentValue{[}InputNotebook{[}{]},\ \{CommonDefaultFormatTypes,\ "Output"\}{]}})
and is error E1 of section 14.

\hypertarget{the-name-of-every-symbol-is-given-explicitly}{%
\subsection{5.3 The name of every symbol is given
explicitly}\label{the-name-of-every-symbol-is-given-explicitly}}

The package builds the name of a new symbol from the letters and scripts
of the box, after removing \texttt{StyleBox} and \texttt{TagBox}
wrappers. Three symbols that differ only in the alphabet of their
letter, for example a bold \(\mathbf e_0\), a sans-serif \(\mathsf e_0\)
and an ordinary \(e_0\) (a letter e with the subscript 0, written in the
three alphabets), would therefore all receive one name and become ONE
symbol. The notebook gives every symbol its own name with the option
\texttt{SymbolizeRootName} (section 8). Error E5.

\hypertarget{reaching-a-symbol-from-a-program}{%
\subsection{5.4 Reaching a symbol from a
program}\label{reaching-a-symbol-from-a-program}}

A program reaches a symbolized symbol through its box:
\texttt{MakeExpression{[}box,\ StandardForm{]}} returns
\texttt{HoldComplete{[}symbol{]}}. To give it a value, the held symbol
is put into an assignment with \texttt{Apply}:

\begin{verbatim}
With[{s = MakeExpression[SubscriptBox["a", "LL"], StandardForm]},
  Set @@ Join[s, HoldComplete[42]]];        (* a_LL = 42 *)
With[{s = MakeExpression[SubscriptBox["a", "LL"], StandardForm]}, Clear @@ s]   (* Clear[a_LL] *)
\end{verbatim}

\texttt{Join{[}HoldComplete{[}a\_LL{]},\ HoldComplete{[}42{]}{]}} is
\texttt{HoldComplete{[}a\_LL,\ 42{]}}, and \texttt{Set\ @@} replaces the
head \texttt{HoldComplete} by \texttt{Set}. Because \texttt{Set} does
not evaluate its first argument, the symbol itself (and not its old
value) receives the new value. An earlier form,
\texttt{ReleaseHold{[}Hold{[}Clear{]}\ @@\ s{]}}, evaluated the symbol
before \texttt{Clear} received it and failed with the message
\texttt{Clear::ssym:\ 42\ is\ not\ a\ symbol} (error E10).

\hypertarget{a-symbolized-symbol-is-one-atom-the-authors-example}{%
\subsection{5.5 A symbolized symbol is one atom: the author's
example}\label{a-symbolized-symbol-is-one-atom-the-authors-example}}

The author showed, with the symbol \(a_{{\rm LL},p}\), that a sum whose
summation variable is called LL multiplies the whole symbol and leaves
the LL inside the subscript untouched. The notebook (section I.2)
contains this example:

\begin{verbatim}
Symbolize[ParsedBoxWrapper[SubscriptBox["a", RowBox[{"LL", ",", "p"}]]], WorkingForm -> StandardForm];
Sum[LL a_{LL,p}, {LL, 1, 2}]                          (* gives 3 a_{LL,p} *)
3 a_{LL,p} === Sum[LL a_{LL,p}, {LL, 1, 2}]           (* gives True *)
Sum[LL Subscript[b, LL, p], {LL, 1, 2}]               (* gives b_{1,p} + 2 b_{2,p} *)
\end{verbatim}

(here \texttt{a\_\{LL,p\}} stands for the two-dimensional input
\(a_{{\rm LL},p}\)). In formulas:
\(\sum_{{\rm LL}=1}^{2}{\rm LL}\,a_{{\rm LL},p}=1\cdot a_{{\rm LL},p}+2\cdot a_{{\rm LL},p}=3\,a_{{\rm LL},p}\),
while for the ordinary subscripted expression the summation variable is
replaced inside the subscript:
\(\sum_{{\rm LL}=1}^{2}{\rm LL}\,b_{{\rm LL},p}=b_{1,p}+2\,b_{2,p}\).
After the example the symbols \texttt{LL} and \texttt{b} of the context
\texttt{Global\textasciigrave{}} are removed with
\texttt{Remove{[}LL,\ b{]}}, because the name \texttt{LL} is used
afterwards for the function \texttt{LL} of the package of the notebook.

\hypertarget{the-symbolized-box-must-be-the-box-that-the-front-end-sends}{%
\subsection{5.6 The symbolized box must be the box that the front end
sends}\label{the-symbolized-box-must-be-the-box-that-the-front-end-sends}}

The Notation package recognizes an input only if its boxes are EXACTLY
the symbolized boxes. Before an input cell is evaluated, the front end
reads every string of its boxes again, as typed text, and cuts it into
tokens: names (a letter followed by letters and digits), numbers and
single other characters. The subscript LL,p of a typed
\(a_{{\rm LL},p}\) therefore arrives as
\texttt{RowBox{[}\{"LL",\ ",",\ "p"\}{]}}. A box written by a program as
\texttt{SubscriptBox{[}"a",\ "LL,p"{]}} (the whole subscript as one
string) is never matched by typed input. A test notebook, evaluated in
the front end, showed exactly this (folder \texttt{experiments}: the sum
gave \(a_{1,p}+2\,a_{2,p}\), \texttt{AtomQ} gave \texttt{False}); for
the registry, 9 symbols were affected (error E7). In the same way, a box
that the front end reads as an operator must be protected (error E9).
Section 6 describes how the converter builds boxes that the front end
sends unchanged, and the tests S-0 and S-14 verify it.

\hypertarget{the-converter-from-tex-to-boxes}{%
\section{6. The converter from TeX to
boxes}\label{the-converter-from-tex-to-boxes}}

The registry stores every symbol in TeX. The notebook contains a reader
of exactly this kind of TeX (function \texttt{SL31TeXToBoxes} of the
package \texttt{SL31\textasciigrave{}}), written in the Wolfram
Language. The Wolfram TeX importer
(\texttt{ImportString{[}...,\ "TeX"{]}}) was not used for symbols
because it loses the alphabets of styled letters (error E4). The rules
of the reader:

\begin{longtable}[]{@{}
  >{\raggedright\arraybackslash}p{(\columnwidth - 2\tabcolsep) * \real{0.5000}}
  >{\raggedright\arraybackslash}p{(\columnwidth - 2\tabcolsep) * \real{0.5000}}@{}}
\toprule\noalign{}
\begin{minipage}[b]{\linewidth}\raggedright
TeX
\end{minipage} & \begin{minipage}[b]{\linewidth}\raggedright
box
\end{minipage} \\
\midrule\noalign{}
\endhead
\bottomrule\noalign{}
\endlastfoot
a letter or digit & the string itself; neighbouring digits are joined
into one number \\
\texttt{\textbackslash{}Delta}, \texttt{\textbackslash{}alpha}, \ldots{}
(Greek letters and other commands) & the Wolfram character, for example
\texttt{"\textbackslash{}{[}CapitalDelta{]}"} \\
\texttt{x\^{}\{...\}}, \texttt{x\_\{...\}}, both &
\texttt{SuperscriptBox}, \texttt{SubscriptBox},
\texttt{SubsuperscriptBox} (a prime \texttt{\textquotesingle{}} is a
superscript \texttt{"\textbackslash{}{[}Prime{]}"}) \\
\texttt{\{\textbackslash{}rm\ text\}} & the text, cut into tokens as the
front end cuts typed text: \texttt{\{\textbackslash{}rm\ vac\}} gives
\texttt{"vac"}, \texttt{\{\textbackslash{}rm\ int,F\}} gives
\texttt{RowBox{[}\{"int",\ ",",\ "F"\}{]}},
\texttt{\{\textbackslash{}rm\ s(2)\}} gives
\texttt{RowBox{[}\{"s",\ RowBox{[}\{"(",\ "2",\ ")"\}{]}\}{]}} \\
\texttt{\textbackslash{}mathbb\ R}, \texttt{\textbackslash{}mathcal\ L},
\texttt{\textbackslash{}mathfrak\ m} & the double-struck, script and
Gothic characters (separate characters of the Wolfram Language) \\
\texttt{\textbackslash{}mathsf\ N}, \texttt{\textbackslash{}mathbf\ x} &
\texttt{TagBox{[}StyleBox{[}"N",\ FontFamily\ -\textgreater{}\ "Helvetica",\ FontSlant\ -\textgreater{}\ "Plain"{]},\ "SL31SansSerif"{]}},
\texttt{TagBox{[}StyleBox{[}"x",\ FontWeight\ -\textgreater{}\ "Bold",\ FontSlant\ -\textgreater{}\ "Plain"{]},\ "SL31Bold"{]}} \\
\texttt{\textbackslash{}vec\ d}, \texttt{\textbackslash{}hat\ x},
\texttt{\textbackslash{}tilde\ x}, \texttt{\textbackslash{}bar\ x} &
\texttt{OverscriptBox} with
\texttt{"\textbackslash{}{[}RightVector{]}"}, \texttt{"\^{}"},
\texttt{"\textasciitilde{}"}, \texttt{"\_"} \\
\texttt{(...)\^{}2} & a script after a closing parenthesis belongs to
the whole group:
\texttt{SuperscriptBox{[}RowBox{[}\{"(",\ ...,\ ")"\}{]},\ "2"{]}}, the
box that the front end makes of typed input \\
\end{longtable}

The \texttt{TagBox} of the sans-serif and bold letters is kept by the
Notation package when it compares boxes (it removes only
\texttt{StyleBox}, \texttt{AdjustmentBox} and \texttt{FrameBox}), so
that a bold, a sans-serif and an ordinary letter with the same scripts
remain three different symbols.

The boxes that are symbolized depend on the kind of the TeX item
(function \texttt{SL31SymbolBoxes}):

\begin{itemize}
\tightlist
\item
  a single letter (\(D\), the number of space dimensions) is wrapped in
  \texttt{TagBox{[}..,\ "SL31Plain"{]}}, so that the ordinary letter
  \texttt{D} keeps its meaning in the Wolfram Language (it is the
  derivative function there);
\item
  a function symbol \(W_{\rm LL}(\mathbf u)\) gives the box of its head
  \(W_{\rm LL}\);
\item
  a group name, a ket, an expression, and a symbol whose base is an
  operator sign (\(\partial_p\), \(\nabla^{\rm LL}_p\),
  \(\nabla^{{\rm LL}*}_p\), \(\nabla^{\rm B}_{{\rm LL},p}\)) is wrapped
  in
  \texttt{TagBox{[}..,\ "SL31Name",\ SyntaxForm\ -\textgreater{}\ "symbol"{]}},
  so that the front end treats it as ONE operand wherever it stands (37
  symbols);
\item
  every other symbol gives its box unchanged.
\end{itemize}

\hypertarget{the-registry-inside-the-notebook}{%
\section{7. The registry inside the
notebook}\label{the-registry-inside-the-notebook}}

The registry of the symbols of the super-loop revision has 199 entries
in nine groups:

\begin{longtable}[]{@{}
  >{\raggedright\arraybackslash}p{(\columnwidth - 2\tabcolsep) * \real{0.5000}}
  >{\raggedright\arraybackslash}p{(\columnwidth - 2\tabcolsep) * \real{0.5000}}@{}}
\toprule\noalign{}
\begin{minipage}[b]{\linewidth}\raggedright
group
\end{minipage} & \begin{minipage}[b]{\linewidth}\raggedright
entries
\end{minipage} \\
\midrule\noalign{}
\endhead
\bottomrule\noalign{}
\endlastfoot
1 Indices, conventions and number sets & 11 \\
2 The bi-octonions and the super-loop & 26 \\
3 Operators on the super-loop and their eigenvalues & 3 \\
4 The covariant families and their classification & 5 \\
5 The manifolds & 8 \\
6 The spin-zero superfluid & 20 \\
7 Power counting and renormalizability & 8 \\
8 Groups, Lie algebras, operations and labels & 6 \\
9 Test 2: the gauge-covariant field equations of the spinor field and
their solutions & 112 \\
\end{longtable}

The registry is stored in the repository as the file
\texttt{Revision\_SuperLoop/provenance/symbols\_ll.csv} (columns group,
symbol, names\_in\_files, meaning, origin). The script
\texttt{tools/registry\_to\_wl.wls} converts it into the Wolfram
Language file \texttt{notebook\_source/registry\_data.wl} (a list
\texttt{SL31Registry} of 199 associations with the keys ``entry'',
``group'', ``symbol'', ``names'', ``meaning'', ``meaningEdited'',
``origin''), and the notebook contains this file in full in one input
cell. In 16 entries (2, 36, 58, 86, 115, 123, 139, 148, 169, 174, 176,
177, 178, 179, 192, 199) the meaning was restated (file
\texttt{notebook\_source/meaning\_edits.wl}, key ``meaningEdited''
-\textgreater{} True), because its original wording pointed to
equations, theorems, sections or files of other texts; for example, one
restated meaning writes out a transformation of the order parameter of
the superfluid in full, where the original text only gave the number of
the equation that defines it.

Every TeX segment (the text between two dollar signs) is split at the
commas that stand outside braces, parentheses and brackets; a list
\texttt{a,\ \textbackslash{}dots,\ b} of numbered symbols is expanded.
This gives 449 TeX items in 434 segments; the 2 label entries (86: the
names of the machine checks and theorems; 87: the hypothesis labels
H-grp, H-op, H-Lie, H-gr, H-dyn) are names, not symbols, and are kept as
text. The kinds of the 449 items:

\begin{longtable}[]{@{}
  >{\raggedright\arraybackslash}p{(\columnwidth - 6\tabcolsep) * \real{0.2500}}
  >{\raggedright\arraybackslash}p{(\columnwidth - 6\tabcolsep) * \real{0.2500}}
  >{\raggedright\arraybackslash}p{(\columnwidth - 6\tabcolsep) * \real{0.2500}}
  >{\raggedright\arraybackslash}p{(\columnwidth - 6\tabcolsep) * \real{0.2500}}@{}}
\toprule\noalign{}
\begin{minipage}[b]{\linewidth}\raggedright
kind
\end{minipage} & \begin{minipage}[b]{\linewidth}\raggedright
items
\end{minipage} & \begin{minipage}[b]{\linewidth}\raggedright
symbolized
\end{minipage} & \begin{minipage}[b]{\linewidth}\raggedright
examples
\end{minipage} \\
\midrule\noalign{}
\endhead
\bottomrule\noalign{}
\endlastfoot
symbol & 289 & yes & \(\Delta^{\rm vac}_{\rm LL}\), \(G_{\rm LL}\) \\
family (running index) & 56 & yes & \(\tau^{\rm LL}_\iota\),
\(x^p_{\rm LL}\) \\
function & 37 & yes (its head) & \(W_{\rm LL}(\mathbf u)\) \\
group & 30 & yes & \({\rm SU}(2)_{\rm v}\), \(\mathfrak{so}(3,1)\) \\
operator & 13 & no (Wolfram functions) & tr, det, diag, Re, Im,
\({}^T\), \(\otimes\), \(\oplus\), \(\circ\), \([\mathbf x,\mathbf y]\),
\([\mathbf x,\mathbf y,\mathbf z]\), \([\ \cdot\ ]\),
\(\langle\ \rangle\) \\
function family & 7 & yes (its head) & \\
set & 4 & no (Wolfram domains) & \(\mathbb R\), \(\mathbb C\),
\(\mathbb H\), \(\mathbb O\) \\
index & 3 & no & \(j\), \(p\), \(\iota\) \\
primed index & 3 & yes & \(p'\), \(\iota'\), \(\iota''\) \\
imaginary unit & 2 & no (the Wolfram symbol \texttt{I}) & \(i\) \\
expression & 2 & yes & \((T_{\rm LL})^2\) \\
ket & 2 & yes & \(\lvert\Phi^{\rm SF}_{\rm LL}\rangle\),
\(\lvert\Phi^{\rm N}_{\rm LL}\rangle\) \\
letter & 1 & yes & \(D\) \\
\end{longtable}

The 427 TeX items of the symbolized kinds give 422 different boxes, so
422 symbolized symbols. Four boxes are shared by several TeX items,
always a function symbol written with different arguments; they are ONE
symbol, as the head \(f\) of \(f(x)\) and \(f(x,y)\) is one symbol:

\begin{longtable}[]{@{}
  >{\raggedright\arraybackslash}p{(\columnwidth - 4\tabcolsep) * \real{0.3333}}
  >{\raggedright\arraybackslash}p{(\columnwidth - 4\tabcolsep) * \real{0.3333}}
  >{\raggedright\arraybackslash}p{(\columnwidth - 4\tabcolsep) * \real{0.3333}}@{}}
\toprule\noalign{}
\begin{minipage}[b]{\linewidth}\raggedright
symbol
\end{minipage} & \begin{minipage}[b]{\linewidth}\raggedright
TeX items
\end{minipage} & \begin{minipage}[b]{\linewidth}\raggedright
entries
\end{minipage} \\
\midrule\noalign{}
\endhead
\bottomrule\noalign{}
\endlastfoot
\(N_{\rm LL}\) &
\texttt{N\_\{\textbackslash{}rm\ LL\}(\textbackslash{}mathbf\ x)},
\texttt{N\_\{\textbackslash{}rm\ LL\}(\textbackslash{}mathbf\ x,\textbackslash{}mathbf\ y)}
& 19 \\
\(\Pi^\iota_{\rm LL}\) &
\texttt{\textbackslash{}Pi\^{}\textbackslash{}iota\_\{\textbackslash{}rm\ LL\}},
\texttt{\textbackslash{}Pi\^{}\textbackslash{}iota\_\{\textbackslash{}rm\ LL\}(x\_\{\textbackslash{}rm\ LL\})}
& 59, 113 \\
\(\varepsilon_{\rm LL}\) &
\texttt{\textbackslash{}varepsilon\_\{\textbackslash{}rm\ LL\}},
\texttt{\textbackslash{}varepsilon\_\{\textbackslash{}rm\ LL\}(+)},
\texttt{\textbackslash{}varepsilon\_\{\textbackslash{}rm\ LL\}(-)} & 62,
156 \\
\(\mathfrak g_2\) & \texttt{\textbackslash{}mathfrak\ g\_2},
\texttt{\textbackslash{}mathfrak\ g\_2(\textbackslash{}mathbb\ C)} &
83 \\
\end{longtable}

The meanings show that each group is one object:
\(N_{\rm LL}(\mathbf x)\) is the complex norm and
\(N_{\rm LL}(\mathbf x,\mathbf y)\) its polarization;
\(\Pi^\iota_{\rm LL}=\sum_{\vec q_{\rm LL}}\Psi_{\rm LL}(\vec q_{\rm LL})^T\mathsf K^\iota_{\rm LL}\Psi_{\rm LL}(-\vec q_{\rm LL})\)
is the pair operator and \(\Pi^\iota_{\rm LL}(x_{\rm LL})\) the same
operator at the point \(x_{\rm LL}\);
\(\varepsilon_{\rm LL}=(q_{\rm LL}^2+m_{\rm LL}^2)^{1/2}\) is the energy
of a spinor charge and
\(\varepsilon_{\rm LL}(\pm)=(\lvert\vec q_{\rm LL}\pm\vec k^{\rm B}_{\rm LL}/2\rvert^2+m_{\rm LL}^2)^{1/2}\)
the same energy at the two momenta of a pair; \(\mathfrak g_2\) is the
Lie algebra of the exceptional group \(G_2\) and
\(\mathfrak g_2(\mathbb C)\) its complexification. (In these meanings
\(\Psi_{\rm LL}\) is the spinor field, the isospin doublet of Dirac
spinors of the spinor charges; \(\mathsf K^\iota_{\rm LL}\) are the
\(8\times8\) pair matrices of the scalar channel; \({}^T\) is the
transpose; \(\vec q_{\rm LL}\) is a momentum and \(q_{\rm LL}\) its
length; \(\vec k^{\rm B}_{\rm LL}\) is the total momentum of a pair,
whose length is the wave number \(k^{\rm B}_{\rm LL}\) of section 2.2;
\(x_{\rm LL}\) is a point of space-time, with the coordinates
\(x^p_{\rm LL}\); the complex norm \(N_{\rm LL}(\mathbf x)\) of a
bi-octonion \(\mathbf x\) is the sum of the squares of its eight complex
coefficients; \(G_2\) is the exceptional Lie group of the automorphisms
of the octonions.) Symbolizing one box twice would only replace the
first symbolization, with the warning \texttt{Symbolize::bsymbexs}
(error E11).

Components of a family, for example \(\sigma^{\rm LL}_1\) of
\(\sigma^{\rm LL}_j\), are separate symbols. A symbolized symbol is an
atom, so the index letter inside it cannot be replaced by a number in a
sum; the function
\texttt{SL31Instance{[}box,\ \{index\ -\textgreater{}\ value\}{]}}
replaces the index in the scripts of the box, symbolizes the new box and
returns its symbol. The tests use the components
\(\sigma^{\rm LL}_{1,2,3}\), the 27 components of
\(\epsilon^{\rm LL}_{jj'j''}\), \(x^{0,1,2,3}_{\rm LL}\), the 16
components of \(\eta^{\rm LL}_{pp'}\), \(\tau^{\rm LL}_{1,2,3}\) and
\(T^{1,2,3}_{\rm LL}\).

\hypertarget{the-names-of-the-symbols}{%
\section{8. The names of the symbols}\label{the-names-of-the-symbols}}

The name of a symbol (function \texttt{SL31NameOf}) is built from its
box: the letters are kept, every character that may not stand in a
symbol name is replaced by its character name, and the separator
\texttt{\textbackslash{}{[}UnderBracket{]}} marks the parts:
\texttt{Sub} and \texttt{Sup} before a subscript and a superscript,
\texttt{Over} before an accent, \texttt{Bold}, \texttt{SansSerif},
\texttt{Plain} and \texttt{Name} for the four kinds of \texttt{TagBox};
a name that would not begin with a letter is prefixed with \texttt{Sym}.
Examples (TeX on the left, the name of the symbol, as the Wolfram
Language writes it, on the right):

\begin{verbatim}
\Delta^{\rm vac}_{\rm LL}    \[CapitalDelta]\[UnderBracket]Sub\[UnderBracket]LL\[UnderBracket]Sup\[UnderBracket]vac
A^\iota_{{\rm LL},p}         A\[UnderBracket]Sub\[UnderBracket]LL\[UnderBracket]RawComma\[UnderBracket]p\[UnderBracket]Sup\[UnderBracket]\[Iota]
\mathsf N^{\rm g}_{\rm LL}    SansSerif\[UnderBracket]N\[UnderBracket]Sub\[UnderBracket]LL\[UnderBracket]Sup\[UnderBracket]g
\mathbf e_0                  Bold\[UnderBracket]e\[UnderBracket]Sub\[UnderBracket]0
D                            Plain\[UnderBracket]D
\partial_p                   Name\[UnderBracket]\[UnderBracket]PartialD\[UnderBracket]\[UnderBracket]Sub\[UnderBracket]p
\end{verbatim}

The internal names never appear in the notebook's output: every symbol
is printed with its box, and (with a rule for \texttt{TraditionalForm})
written in TeX from its box.

\hypertarget{the-tests-of-part-i}{%
\section{9. The tests of Part I}\label{the-tests-of-part-i}}

Every test is recorded by the function
\texttt{SL31Check{[}name,\ statement,\ condition,\ details{]}} in the
list \texttt{SL31\textasciigrave{}Test\textasciigrave{}\$Results} and
prints one line. All 15 tests passed in the evaluation of 17 November
2026; the numbers below are the numbers of that evaluation.

\textbf{S-0 (the strings as the front end reads them).} Every string
inside every symbolized box (the tags of \texttt{TagBox} and the options
of \texttt{StyleBox} excepted) must be one token for the front end; the
test reads each with
\texttt{FrontEnd\textasciigrave{}UndocumentedTestFEParserPacket{[}string,\ True{]}}
and requires that the result is the string itself. Result: 201 different
strings, none cut.

\textbf{S-1.} Every one of the 449 TeX items is converted into a box
structure (no \texttt{\$Failed}, no \texttt{Null}). Result: 449.

\textbf{S-2.} Every one of the 422 symbolized boxes is read by
\texttt{MakeExpression} as one atomic symbol. Result: 422.

\textbf{S-14 (the typed symbols).} One input cell contains all 422
symbols, each in its printed form, as a reader would type them. When the
cell is evaluated, the front end sends the boxes to the kernel; every
one must be one atomic symbol, and it must be the very symbol created
from its registry TeX. Result: 422 of 422.

\textbf{S-3 (uniqueness).} Different boxes give different symbols, and
the boxes stay different after the removal of \texttt{StyleBox},
\texttt{AdjustmentBox} and \texttt{FrameBox} (the comparison that the
Notation package makes). Result: 422 boxes, 422 symbols, 422 boxes after
the removal.

\textbf{S-4 (round trip).} Every symbol printed again gives exactly its
box: \texttt{ToBoxes{[}symbol,\ StandardForm{]}\ ===\ box}. Result: 422
of 422.

\textbf{S-5 (TeX).} \texttt{TeXForm} of every symbol is written from its
printed form; the internal name with its separator does not appear.
Result: 422 of 422.

\textbf{S-6 (arithmetic: the gap equation).} The gap of the condensed
vacuum solves
\[\frac{1}{G_{\rm LL}}=\frac{1}{2\pi^2}\Big[\Lambda_{\rm LL}\big(\Lambda_{\rm LL}^2+(M^{\rm q}_{\rm LL})^2\big)^{1/2}-(M^{\rm q}_{\rm LL})^2\,{\rm arsinh}\frac{\Lambda_{\rm LL}}{M^{\rm q}_{\rm LL}}\Big],\qquad (M^{\rm q}_{\rm LL})^2=m_{\rm LL}^2+(\Delta^{\rm vac}_{\rm LL})^2 .\]
The notebook writes this equation with the symbolized symbols, gives
them the values \(G_{\rm LL}=45/2\), \(m_{\rm LL}=1/5\),
\(\Lambda_{\rm LL}=1\) only inside \texttt{Block}, and solves for
\(\Delta^{\rm vac}_{\rm LL}\) with \texttt{FindRoot} (30 digits, start
value 0.2). Result: \(\Delta^{\rm vac}_{\rm LL}=0.209493303270\); the
conditions
\(\lvert\Delta^{\rm vac}_{\rm LL}-0.2094933033\rvert<10^{-10}\) (the
value found by the Rust program of Test 2),
\(\lvert2\Delta^{\rm vac}_{\rm LL}-0.4189866\rvert<10^{-7}\) and
\(\lvert2M^{\rm q}_{\rm LL}-0.5792666\rvert<10^{-7}\) hold.

\textbf{S-7 (sums over a family: the Pauli matrices).} With the three
symbolized components \(\sigma^{\rm LL}_{1,2,3}\) and the 27 symbolized
components of \(\epsilon^{\rm LL}_{jj'j''}\), given their values inside
\texttt{Block}:
\[\sum_{j=1}^3\sigma^{\rm LL}_j\sigma^{\rm LL}_j-3\,\mathsf 1_2=0,\qquad \sigma^{\rm LL}_j\sigma^{\rm LL}_{j'}-\Big(\delta_{jj'}\mathsf 1_2+i\sum_{j''=1}^3\epsilon^{\rm LL}_{jj'j''}\sigma^{\rm LL}_{j''}\Big)=0\quad\text{for all nine pairs }(j,j').\]
Result: the largest absolute value of all matrix elements of both
left-hand sides is 0 and 0 (exact).

\textbf{S-8 (a double sum with a metric).} With the 4 components
\(x^p_{\rm LL}\) and the 16 components \(\eta^{\rm LL}_{pp'}\) as free
symbols,
\(Q=\sum_{p,p'=0}^{3}\eta^{\rm LL}_{pp'}x^p_{\rm LL}x^{p'}_{\rm LL}\)
has 16 terms; after the metric receives its values, \texttt{Expand}
gives
\[Q=(x^0_{\rm LL})^2-(x^1_{\rm LL})^2-(x^2_{\rm LL})^2-(x^3_{\rm LL})^2,\]
with the coordinate components still free symbols. Result: equal (the
difference expands to 0).

\textbf{S-9 (algebra: the Higgs mechanism).} With
\((M^{\rm W}_{\rm LL})^2=2Z^{\rm s}_{\rm LL}(g^{\rm v}_{\rm LL})^2(\Delta^{\rm vac}_{\rm LL})^2\)
and
\((M^{\rm F}_{\rm LL})^2=8Z^{\rm s}_{\rm LL}(g^{\rm F}_{\rm LL})^2(\Delta^{\rm vac}_{\rm LL})^2\),
\texttt{Simplify} gives
\[\Big(\frac{M^{\rm F}_{\rm LL}}{M^{\rm W}_{\rm LL}}\Big)^2=4\Big(\frac{g^{\rm F}_{\rm LL}}{g^{\rm v}_{\rm LL}}\Big)^2,\]
independent of the stiffness and of the gap, and \texttt{Solve} recovers
the gap from \(M^{\rm W}_{\rm LL}\):
\(\Delta^{\rm vac}_{\rm LL}=\pm M^{\rm W}_{\rm LL}/\big(\sqrt2\,g^{\rm v}_{\rm LL}(Z^{\rm s}_{\rm LL})^{1/2}\big)\)
(two solutions). Result: as stated.

\textbf{S-10 (algebra: the longitudinal vector boson).} For a plane wave
the time component and the longitudinal component of a massive vector
field obey two linear equations with the determinant
\[\det=\Big((k^{\rm B}_{\rm LL})^2+\frac{Z^{\rm t}_{\rm LL}}{Z^{\rm s}_{\rm LL}}M^2\Big)\big((\omega^{\rm B}_{\rm LL})^2-M^2\big)-(k^{\rm B}_{\rm LL}\omega^{\rm B}_{\rm LL})^2,\qquad M=M^{\rm W}_{\rm LL}.\]
\texttt{Factor} gives
\(\det=-\dfrac{M^2}{Z^{\rm s}_{\rm LL}}\Big(Z^{\rm s}_{\rm LL}(k^{\rm B}_{\rm LL})^2+Z^{\rm t}_{\rm LL}M^2-Z^{\rm t}_{\rm LL}(\omega^{\rm B}_{\rm LL})^2\Big)\),
and \texttt{Solve} (for the unknown \(w=(\omega^{\rm B}_{\rm LL})^2\))
gives the single root
\(w=M^2+(Z^{\rm s}_{\rm LL}/Z^{\rm t}_{\rm LL})(k^{\rm B}_{\rm LL})^2\):
the absorbed Goldstone boson moves with the velocity
\((Z^{\rm s}_{\rm LL}/Z^{\rm t}_{\rm LL})^{1/2}\). Result: one root,
equal to the stated one.

\textbf{S-11 (a sum over the isospin index: the Casimir).} With the
symbolized components \(T^\iota_{\rm LL}\) and \(\tau^{\rm LL}_\iota\)
and the rule \(T^\iota_{\rm LL}=\tau^{\rm LL}_\iota/2\):
\[(T_{\rm LL})^2=\sum_{\iota=1}^3T^\iota_{\rm LL}T^\iota_{\rm LL}=\tfrac34\,\mathsf 1_2 .\]
Result: \(\{\{3/4,0\},\{0,3/4\}\}\).

\textbf{S-12 (plain-text names).} No plain-text name of the registry
belongs to two different entries. Result: 592 different names, none
shared.

\textbf{S-13 (the kinds that are not symbolized).} The imaginary unit
\(i\) is the Wolfram symbol \texttt{I}; \(\mathbb R\) and \(\mathbb C\)
are the domains \texttt{Reals} and \texttt{Complexes}; tr, det, diag,
Re, Im, \(\otimes\) and \({}^T\) are \texttt{Tr}, \texttt{Det},
\texttt{DiagonalMatrix}, \texttt{Re}, \texttt{Im},
\texttt{KroneckerProduct} and \texttt{Transpose}. The test checks eleven
statements on the Pauli matrices: \(i^2=-1\); \(\sqrt2\in\mathbb R\);
\(i\in\mathbb C\); \(i\notin\mathbb R\);
\({\rm tr}\,\sigma^{\rm LL}_3=0\); \(\det\sigma^{\rm LL}_1=-1\);
\({\rm diag}(1,-1)=\sigma^{\rm LL}_3\);
\({\rm Re}\,\sigma^{\rm LL}_2=0\);
\({\rm Im}\,\sigma^{\rm LL}_2=\begin{pmatrix}0&-1\\1&0\end{pmatrix}\);
\((\sigma^{\rm LL}_1\otimes\sigma^{\rm LL}_3)^2=\mathsf 1_4\)
(\(\mathsf 1_4\): the \(4\times4\) unit matrix);
\((\sigma^{\rm LL}_2)^T=-\sigma^{\rm LL}_2\). Result: all eleven
\texttt{True}.

\hypertarget{part-ii-the-encyclopedia-of-the-programs-of-test-2}{%
\section{10. Part II: the encyclopedia of the programs of Test
2}\label{part-ii-the-encyclopedia-of-the-programs-of-test-2}}

\hypertarget{what-part-ii-contains}{%
\subsection{10.1 What Part II contains}\label{what-part-ii-contains}}

Test 2 of the super-loop revision of the project asks eleven questions,
called items {[}1{]} to {[}11{]}: the gauge-covariant field equations of
the self-interacting spinor field \(\Psi_{\rm LL}\) on
\(\mathbb R^{1,3}_{\rm LL}\) and their checks (item {[}1{]}); their
analysis and analytic solutions (item {[}2{]}); a solver of the field
equations and a plotter of its solutions, written twice, in Python and
in Rust (items {[}3{]} to {[}6{]}); the construction and the physics of
the superfluid boson singlet state, SFBS (items {[}7{]} to {[}11{]}).
The original programs are in the folder
\texttt{Revision\_SuperLoop/test2} of the repository:

\begin{itemize}
\tightlist
\item
  the Rust library \texttt{t2core} (files \texttt{rng.rs},
  \texttt{c64.rs}, \texttt{cmat.rs}, \texttt{model.rs},
  \texttt{lattice.rs}, \texttt{sym.rs}, \texttt{lib.rs}) and the Rust
  programs \texttt{t2\_fields} (item {[}1{]}), \texttt{t2\_analytic}
  (item {[}2{]}), \texttt{t2\_solve} (items {[}3{]}, {[}5{]}),
  \texttt{t2\_plot} (items {[}4{]}, {[}6{]}) and \texttt{t2\_sfbs}
  (items {[}7{]} to {[}11{]}), 12 615 lines in all;
\item
  the Python programs \texttt{t2\_solve.py}, \texttt{t2\_plot.py},
  \texttt{t2\_examples\_check.py}, \texttt{make\_t2\_figures.py} and
  \texttt{make\_t2\_sfbs\_figures.py}.
\end{itemize}

Part II of the notebook rewrites all of them, item by item, in the
Wolfram Language, with every check of the original programs.
\textbf{This stage} contains chapter II.1 (the library \texttt{t2core})
and chapter II.2 (item {[}1{]}, the program \texttt{t2\_fields}: the
checks GS-1 to GS-16, the two figures and the data files of the
figures). The items {[}2{]} to {[}11{]} follow in the next stages; each
will be added to this document when it is added to the notebook.

\hypertarget{names-indices-and-numbers}{%
\subsection{10.2 Names, indices and
numbers}\label{names-indices-and-numbers}}

\begin{itemize}
\tightlist
\item
  \textbf{Names.} Every function of Part II begins with \texttt{t2} and
  keeps the name of the original function in the style of the Wolfram
  Language (for example the Rust function \texttt{model::gamma\_ch} is
  \texttt{t2gammaCh}). Every name stands for one object only (section
  15).
\item
  \textbf{Indices.} The indices keep the values of the physics,
  \(p=0,\dots,3\) and \(j,\iota=1,2,3\); because the lists of the
  Wolfram Language begin at position 1, the code writes the entry of
  \(p\) at position \(p+1\).
\item
  \textbf{Numbers.} The original programs compute in double precision
  (about 16 decimal digits). Part II does the same. Exact objects (the
  Dirac matrices, the generators, the pair matrices) are computed
  exactly, with integers and the imaginary unit.
\item
  \textbf{Results compared with the original report.} Every printed
  check line ends with the numbers of the Wolfram Language and, in
  parentheses, those of the original report (the file
  \texttt{Revision\_SuperLoop/test2/data/t2\_fields\_checks.txt},
  written by the Rust program). Results that are fixed by the random
  numbers alone (the spectral radius 8.4349 and its bound 17.5474 of
  GS-10, the factors of GS-13, the smallest Gram eigenvalue 8.000 of
  GS-15) agree to all printed digits. Results that measure rounding
  errors (the largest deviation of an exact identity, of the size
  \(10^{-16}\) to \(10^{-13}\)) agree in their order of magnitude but
  not in their digits: a sum of floating-point numbers depends in its
  last bits on the order of the terms, and the matrix products of the
  Wolfram Language add the terms in a different order than the loops of
  the Rust code. Such a deviation is the accumulated rounding error of
  the computation and has no further meaning; what matters is that it
  stays far below the tolerance of the check.
\end{itemize}

\hypertarget{random-numbers}{%
\subsection{10.3 Random numbers}\label{random-numbers}}

The original programs draw their random matrices, fields and momenta
from the generator SplitMix64, whose state is one unsigned integer with
64 bits. One step, written with the names of the program (all arithmetic
modulo \(2^{64}\); \texttt{XOR} is the exclusive or of the bits,
\texttt{\textgreater{}\textgreater{}} the shift of the bits to the
right; the constants are hexadecimal):

\begin{verbatim}
state  = state + 0x9E3779B97F4A7C15
t      = state
t      = (t XOR (t >> 30)) * 0xBF58476D1CE4E5B9
t      = (t XOR (t >> 27)) * 0x94D049BB133111EB
output = t XOR (t >> 31)
\end{verbatim}

A uniform number in \([-1,1)\) is
\texttt{2\ (output\ \textgreater{}\textgreater{}\ 11)/2\^{}53\ -\ 1};
the division and the subtraction are exact in double precision, so the
Wolfram Language (\texttt{t2Seed}, \texttt{t2NextU64}, \texttt{t2Unit};
the integer operations \texttt{Mod}, \texttt{BitXor},
\texttt{BitShiftRight}) produces exactly the numbers of the Rust code. A
random complex number is drawn as real part, then imaginary part
(\texttt{t2UnitC}). Part II uses the seeds of the original programs
(2026111707 for \texttt{t2\_fields}) and draws the numbers in the same
order, so that every random object is the same as in the original
program. The check T2-RNG compares the first three integers and the
first four uniform numbers for the seed 2026111707 with the values of
the same algorithm computed independently in Python 3 (script
\texttt{tools/splitmix64\_reference.py}, output
\texttt{logs/T2-RNG\_reference\_values.txt}): 16253586484067153034,
11771902715344777314, 8742176756923514478 and 0.7622173776706638,
0.2763122498264816, \(-0.05217129678912458\), 0.4038650534416697; all
are equal.

\hypertarget{components-of-families}{%
\subsection{10.4 Components of families}\label{components-of-families}}

A component of a registered family, for example \(\gamma^0_{\rm LL}\) of
\(\gamma^p_{\rm LL}\), is a symbol of its own, made from the box of the
family with the index replaced by its value (\texttt{SL31Instance},
section 7). When an input cell is evaluated, the front end sends the
whole cell to the kernel, which reads all of it before it evaluates its
first line; a component typed in the cell must therefore be a symbol
before the cell is evaluated. The builder collects every component that
is typed anywhere in the notebook (every reference
\texttt{LL{[}"tex",\ "index"\ -\textgreater{}\ "value",\ ...{]}} of the
source) and writes the list into the first input cell of Part II, which
makes the components that Part I has not made, gives them their
\texttt{TraditionalForm} boxes and usage messages, and checks them
(T2-SYM): 10 components are typed in the notebook
(\(\sigma^{\rm LL}_0,\dots,\sigma^{\rm LL}_3\),
\(\gamma^0_{\rm LL},\dots,\gamma^3_{\rm LL}\),
\(\mathsf A_{{\rm LL},j}\), \(q^1_{\rm LL}\)), 6 of them are made in
this cell, and all 10 are read as atomic symbols whose TeX contains no
internal name. Error E18 (section 14) describes what happened before
this cell existed.

\hypertarget{part-ii-chapter-ii.1-the-library-t2core-in-the-wolfram-language}{%
\section{11. Part II, chapter II.1: the library t2core in the Wolfram
Language}\label{part-ii-chapter-ii.1-the-library-t2core-in-the-wolfram-language}}

\hypertarget{complex-numbers-and-matrices-c64.rs-cmat.rs}{%
\subsection{11.1 Complex numbers and matrices (c64.rs,
cmat.rs)}\label{complex-numbers-and-matrices-c64.rs-cmat.rs}}

The Rust library defines its own complex numbers and dense complex
matrices; the Wolfram Language has both built in (section 2.1). Three
routines of \texttt{cmat.rs} are rewritten because the programs and
their checks use them:

\begin{itemize}
\tightlist
\item
  \texttt{t2det2}, \texttt{t2inv2}: the determinant and the inverse of a
  \(2\times2\) matrix by the explicit formulas: for the matrix
  \texttt{\{\{m11,\ m12\},\ \{m21,\ m22\}\}} the determinant is
  \texttt{m11\ m22\ -\ m12\ m21} and the inverse is
  \texttt{\{\{m22,\ -m12\},\ \{-m21,\ m11\}\}} divided by the
  determinant. \texttt{t2csqrt}: the principal square root of a complex
  number, the square root of its modulus times the phase factor of half
  its argument.
\item
  \texttt{t2expm}: the matrix exponential by scaling and squaring. The
  matrix is halved \texttt{s} times, until its Frobenius norm is at most
  \(1/4\); the Taylor series of the exponential of the halved matrix
  \texttt{x} (the sum of its powers divided by the factorials) is summed
  until a term has a Frobenius norm below \(10^{-18}\) (at most 29
  terms); the sum is then squared \texttt{s} times, because the
  exponential of twice a matrix is the square of the exponential of the
  matrix.
\item
  \texttt{t2eigh}: all eigenvalues and eigenvectors of a complex
  Hermitian matrix, in two stages. Stage 1, Householder reduction: for
  every column \texttt{k} a reflection (the unit matrix minus twice the
  projector on a complex unit vector \texttt{vv}) is chosen that sets to
  zero the entries of the column below the subdiagonal; applied from
  both sides it keeps the eigenvalues, and after all columns the matrix
  is tridiagonal; the product of the reflections is accumulated in the
  matrix \texttt{q}. The complex subdiagonal entries are then made real
  and non-negative by a diagonal unitary matrix of phases. Stage 2, the
  implicit QL algorithm with Wilkinson shifts (\texttt{tqli}): plane
  rotations drive the off-diagonal entries of the real tridiagonal
  matrix to zero; the diagonal then holds the eigenvalues, and the
  accumulated rotations, multiplied by \texttt{q}, give the
  eigenvectors. The eigenvalues are sorted in ascending order.
\end{itemize}

Check \textbf{T2-LINALG} (seed 12345): for a random \(9\times9\)
Hermitian matrix (\texttt{z\ +\ ConjugateTranspose{[}z{]}} with a random
complex matrix \texttt{z}) the eigenvalues of \texttt{t2eigh} differ
from those of the built-in \texttt{Eigenvalues} by at most
\(5.33\times10^{-15}\) (tolerance \(10^{-12}\)), the eigenvectors
reproduce the matrix (Frobenius norm of the difference below
\(10^{-12}\)), and for a random \(3\times3\) matrix \texttt{t2expm}
differs from the built-in \texttt{MatrixExp} by \(1.24\times10^{-14}\)
in the Frobenius norm (tolerance \(10^{-13}\)). The full test of the
eigensolver is GS-11.

\hypertarget{the-super-loop-operators-as-matrices-model.rs}{%
\subsection{11.2 The super-loop operators as matrices
(model.rs)}\label{the-super-loop-operators-as-matrices-model.rs}}

Every function of \texttt{model.rs} is rewritten, in the same order; the
formulas are those of section 2.5. In particular: \texttt{t2mq} is
\(M_{\rm LL}\); \texttt{t2bqFrom} is its inverse, which gives for a
\(2\times2\) matrix \texttt{m} the coefficients
\(\xi^{\rm LL}_0=(\)\texttt{m11}\(+\)\texttt{m22}\()/2\),
\(\xi^{\rm LL}_1=i(\)\texttt{m12}\(+\)\texttt{m21}\()/2\),
\(\xi^{\rm LL}_2=(\)\texttt{m21}\(-\)\texttt{m12}\()/2\),
\(\xi^{\rm LL}_3=i(\)\texttt{m11}\(-\)\texttt{m22}\()/2\) (these follow
from
\(M_{\rm LL}(\mathbf b)=\begin{pmatrix}\xi^{\rm LL}_0-i\xi^{\rm LL}_3&-i\xi^{\rm LL}_1-\xi^{\rm LL}_2\\-i\xi^{\rm LL}_1+\xi^{\rm LL}_2&\xi^{\rm LL}_0+i\xi^{\rm LL}_3\end{pmatrix}\));
\texttt{t2bqMul} multiplies two bi-quaternions through their matrices;
\texttt{t2bioctMul} multiplies two bi-octonions in Cayley-Dickson form,

\[(\mathbf a+\mathbf b\,\mathbf e_4)(\mathbf c+\mathbf d\,\mathbf e_4)=(\mathbf a\mathbf c-\bar{\mathbf d}\mathbf b)+(\mathbf d\mathbf a+\mathbf b\bar{\mathbf c})\mathbf e_4 ;\]

\texttt{t2randSL2C} draws a random element of \(SL(2,\mathbb C)\): four
random complex entries (in the order \texttt{m11}, \texttt{m12},
\texttt{m21}, \texttt{m22}), accepted if the modulus of the determinant
exceeds 0.1, then divided by the principal square root of the
determinant, so that the determinant is 1; \texttt{t2su2} is
\(\mathsf U_{\rm LL}=\exp\big(i\sum_\iota\theta^\iota_{\rm LL}\tau^{\rm LL}_\iota/2\big)\)
(with \texttt{t2expm}); \texttt{t2adjoint} is
\(\mathsf R^{\rm v}_{\rm LL}\), with
\([\mathsf R^{\rm v}_{\rm LL}(\mathsf U_{\rm LL})]_{\iota'\iota}=\tfrac12{\rm Re}\,{\rm tr}\big(\tau^{\rm LL}_{\iota'}\mathsf U_{\rm LL}\tau^{\rm LL}_\iota\mathsf U_{\rm LL}^\dagger\big)\);
\texttt{t2eps3} is the Levi-Civita symbol (\texttt{Signature}). The
notebook shows the matrices
\(\sigma^{\rm LL}_0,\dots,\sigma^{\rm LL}_3\),
\(\gamma^0_{\rm LL},\dots,\gamma^3_{\rm LL}\),
\(\gamma^{\rm ch}_{\rm LL}\), \(C_{\rm LL}\), \(\beta_{\rm LL}\) and
\(\mathsf E_{\rm LL}\), each labelled with its registered symbol.

\hypertarget{the-lattice-operator-lattice.rs}{%
\subsection{11.3 The lattice operator
(lattice.rs)}\label{the-lattice-operator-lattice.rs}}

The solvers discretize space on a lattice of
\(n^{\rm x}_{\rm LL}\times n^{\rm y}_{\rm LL}\times n^{\rm z}_{\rm LL}\)
sites with the spacing \(\ell^{\rm x}_{\rm LL}\). A site
\(\vec s_{\rm LL}\) has the number
\(s^1_{\rm LL}+n^{\rm x}_{\rm LL}(s^2_{\rm LL}+n^{\rm y}_{\rm LL}s^3_{\rm LL})\)
in the code (\texttt{t2siteOf}, \texttt{t2siteCoords}); a lattice is an
association with the keys \texttt{"n"}, \texttt{"ell"},
\texttt{"periodic"} (periodic or Dirichlet boundary in each direction),
\texttt{"wilson"}, \texttt{"mass"}, \texttt{"mu"}. On every site a mode
has 16 components, \(\psi^{\rm p}_{\rm LL}\) and
\(\psi^{\rm h}_{\rm LL}\). The link matrix (\texttt{t2link}) has the
closed form

\[\mathsf U^{\rm link}_{{\rm LL},j}=e^{-i\ell^{\rm x}_{\rm LL}g^{\rm F}_{\rm LL}A^{\rm F}_{{\rm LL},j}}\Big(\cos\big(\tfrac12\ell^{\rm x}_{\rm LL}g^{\rm v}_{\rm LL}\lvert\vec A^{\rm v}_{{\rm LL},j}\rvert\big)\sigma^{\rm LL}_0-i\sin\big(\tfrac12\ell^{\rm x}_{\rm LL}g^{\rm v}_{\rm LL}\lvert\vec A^{\rm v}_{{\rm LL},j}\rvert\big)\frac{\sum_\iota A^\iota_{{\rm LL},j}\tau^{\rm LL}_\iota}{\lvert\vec A^{\rm v}_{{\rm LL},j}\rvert}\Big),\]

because
\(\big(\sum_\iota A^\iota_{{\rm LL},j}\tau^{\rm LL}_\iota\big)^2=\lvert\vec A^{\rm v}_{{\rm LL},j}\rvert^2\sigma^{\rm LL}_0\).
The potential of the time direction at a site (\texttt{t2potential}) is
the \(2\times2\) matrix
\(g^{\rm v}_{\rm LL}\mathsf a_{{\rm LL},0}+g^{\rm F}_{\rm LL}A^{\rm F}_{{\rm LL},0}\sigma^{\rm LL}_0\),
which acts on the eight components as \(\mathsf A_{{\rm LL},0}\). The
particle row of \(\mathcal H^{\rm lat}_{\rm LL}\) is

\[\sum_j\Big[-\frac{i\alpha^j_{\rm LL}}{2\ell^{\rm x}_{\rm LL}}\big(\mathsf U^{\rm link}_{{\rm LL},j}(\vec s_{\rm LL})\psi^{\rm p}_{\rm LL}(\vec s_{\rm LL}+\vec\delta^j_{\rm LL})-\mathsf U^{\rm link}_{{\rm LL},j}(\vec s_{\rm LL}-\vec\delta^j_{\rm LL})^\dagger\psi^{\rm p}_{\rm LL}(\vec s_{\rm LL}-\vec\delta^j_{\rm LL})\big)\]
\[-\frac{r^{\rm W}_{\rm LL}}{2\ell^{\rm x}_{\rm LL}}\beta_{\rm LL}\big(\mathsf U^{\rm link}_{{\rm LL},j}(\vec s_{\rm LL})\psi^{\rm p}_{\rm LL}(\vec s_{\rm LL}+\vec\delta^j_{\rm LL})+\mathsf U^{\rm link}_{{\rm LL},j}(\vec s_{\rm LL}-\vec\delta^j_{\rm LL})^\dagger\psi^{\rm p}_{\rm LL}(\vec s_{\rm LL}-\vec\delta^j_{\rm LL})-2\psi^{\rm p}_{\rm LL}(\vec s_{\rm LL})\big)\Big]\]
\[+m_{\rm LL}\beta_{\rm LL}\psi^{\rm p}_{\rm LL}(\vec s_{\rm LL})-\big(\mathsf A_{{\rm LL},0}(\vec s_{\rm LL})+\mu_{\rm LL}\big)\psi^{\rm p}_{\rm LL}(\vec s_{\rm LL})+\mathsf P_{\rm LL}(\vec s_{\rm LL})\psi^{\rm h}_{\rm LL}(\vec s_{\rm LL}),\]

where \(\alpha^j_{\rm LL}\) and \(\beta_{\rm LL}\) act on the Dirac
index and the link matrices on the isospin index; a direction without a
neighbour (a Dirichlet boundary, or a direction with one site)
contributes zero (\texttt{t2neighbour} returns \texttt{None} there). The
hole row is
\(\mathsf P_{\rm LL}(\vec s_{\rm LL})^\dagger\psi^{\rm p}_{\rm LL}(\vec s_{\rm LL})\)
minus the complex conjugate of the same expression without the pair
term, applied to \((\psi^{\rm h}_{\rm LL})^*\) (\texttt{t2hMinusV},
\texttt{t2apply}). \texttt{t2apply} applies
\(\mathcal H^{\rm lat}_{\rm LL}\) to a vector or, column by column, to a
matrix; \texttt{t2dense} applies it to the unit matrix and returns the
whole matrix. Further functions: \texttt{t2normBound}, the bound

\[B^{\rm H}_{\rm LL}=n^{\rm act}_{\rm LL}\frac{1+2r^{\rm W}_{\rm LL}}{\ell^{\rm x}_{\rm LL}}+\lvert m_{\rm LL}\rvert+\lvert\mu_{\rm LL}\rvert+\max_{\vec s_{\rm LL}}\lVert g^{\rm v}_{\rm LL}\mathsf a_{{\rm LL},0}+g^{\rm F}_{\rm LL}A^{\rm F}_{{\rm LL},0}\sigma^{\rm LL}_0\rVert_{\rm F}+\max_{\vec s_{\rm LL}}\sum_\iota\lvert d^\iota_{\rm LL}(\vec s_{\rm LL})\rvert\]

(\(\lVert\cdot\rVert_{\rm F}\): the Frobenius norm); \texttt{t2epsLat},
the energy \(\varepsilon^{\rm lat}_{\rm LL}\) of section 12.2 (GS-16);
\texttt{t2hLatPlane}, \(\mathsf h^{\rm lat}_{\rm LL}(\vec q_{\rm LL})\);
\texttt{t2anomalous}, the lattice pair densities
\(\Pi^\iota_{\rm LL}(\vec s_{\rm LL})=\sum_n\psi^{\rm p}_{{\rm LL},n}(\vec s_{\rm LL})^T\mathsf K^\iota_{\rm LL}\psi^{\rm h}_{{\rm LL},n}(\vec s_{\rm LL})^*\)
of a set of modes; \texttt{t2rk4Step}, one step of the Runge-Kutta
method of fourth order for
\(i\partial_0\)\texttt{w}\(=\mathcal H^{\rm lat}_{\rm LL}\)\texttt{w}
with the time step \(\ell^{\rm t}_{\rm LL}\) (\texttt{k1\ =\ -i\ H\ w},
\texttt{k2\ =\ -i\ H\ (w\ +\ dt\ k1/2)},
\texttt{k3\ =\ -i\ H\ (w\ +\ dt\ k2/2)},
\texttt{k4\ =\ -i\ H\ (w\ +\ dt\ k3)}, new
\texttt{w\ =\ w\ +\ dt\ (k1\ +\ 2\ k2\ +\ 2\ k3\ +\ k4)/6}, with
\texttt{dt} \(=\ell^{\rm t}_{\rm LL}\)).

\hypertarget{the-symbolic-algebra-sym.rs}{%
\subsection{11.4 The symbolic algebra
(sym.rs)}\label{the-symbolic-algebra-sym.rs}}

\texttt{sym.rs} is a small symbolic algebra of the Rust library (linear
combinations of monomials written in TeX, with exact coefficients
\(a+bi\), \(a\) and \(b\) rational), used to print matrices and
component equations in TeX and to evaluate every printed formula
numerically. The Wolfram Language is itself a symbolic algebra with
exact Gaussian rationals; the module is therefore not rewritten, and
check GS-12 builds and evaluates the same formulas with symbols of the
Wolfram Language.

\hypertarget{part-ii-item-1-the-field-equations-on-the-31-manifold-program-t2_fields}{%
\section{12. Part II, item {[}1{]}: the field equations on the 3+1
manifold (program
t2\_fields)}\label{part-ii-item-1-the-field-equations-on-the-31-manifold-program-t2_fields}}

\hypertarget{the-physics-of-item-1}{%
\subsection{12.1 The physics of item
{[}1{]}}\label{the-physics-of-item-1}}

The covariant derivative of the spinor field is
\(\nabla^{\rm LL}_p=\partial_p+\Gamma^{\rm LL}_p-i\mathsf A_{{\rm LL},p}\),
with the spin connection
\(\Gamma^{\rm LL}_p=\mathsf s_{\rm LL}(\mathbf u^\omega_{{\rm LL},p})\otimes\sigma^{\rm LL}_0\)
and the gauge matrix \(\mathsf A_{{\rm LL},p}\) (section 2.5). The
mean-field equations of a mode
\((\psi^{\rm p}_{\rm LL},\psi^{\rm h}_{\rm LL})\) are
\(i\partial_0(\psi^{\rm p}_{\rm LL},\psi^{\rm h}_{\rm LL})=\mathcal H^{\rm DB}_{\rm LL}(\psi^{\rm p}_{\rm LL},\psi^{\rm h}_{\rm LL})\)
with the Dirac-Bogoliubov-de Gennes operator

\[\mathcal H^{\rm DB}_{\rm LL}=\begin{pmatrix}h^\nabla_{\rm LL}-i\Gamma^{\rm LL}_0-\mathsf A_{{\rm LL},0}-\mu_{\rm LL}&\mathsf P_{\rm LL}\\\mathsf P_{\rm LL}^\dagger&-\big(h^\nabla_{\rm LL}-i\Gamma^{\rm LL}_0-\mathsf A_{{\rm LL},0}-\mu_{\rm LL}\big)^*\end{pmatrix},\]

with the covariant Dirac Hamiltonian

\[h^\nabla_{\rm LL}=-i\sum_j\alpha^j_{\rm LL}\nabla^{\rm LL}_j+\beta_{\rm LL}m_{\rm LL},\]

and the order parameter is fixed by the self-consistency
\(d^\iota_{\rm LL}=-\tfrac12G_{\rm LL}\langle\Pi^\iota_{\rm LL}\rangle\).
On a plane wave \(\exp(i\sum_jq^j_{\rm LL}x^j_{\rm LL})\) with constant
fields and without the spin connection, \(\mathcal H^{\rm DB}_{\rm LL}\)
becomes the \(16\times16\) matrix
\(\mathcal H^{\rm DB}_{\rm LL}(\vec q_{\rm LL})\) (code: \texttt{hq})
with the upper left block

\[\sum_j\alpha^j_{\rm LL}\big(q^j_{\rm LL}\mathsf 1_8-\mathsf A_{{\rm LL},j}\big)+\beta_{\rm LL}m_{\rm LL}-\mathsf A_{{\rm LL},0}-\mu_{\rm LL},\]

the off-diagonal blocks \(\mathsf P_{\rm LL}\) and
\(\mathsf P_{\rm LL}^\dagger\), and the lower right block

\[-\Big(-\sum_j\alpha^j_{\rm LL}\big(q^j_{\rm LL}\mathsf 1_8+\mathsf A_{{\rm LL},j}\big)+\beta_{\rm LL}m_{\rm LL}-\mathsf A_{{\rm LL},0}-\mu_{\rm LL}\Big)^*\]

(the hole component carries the momentum \(-\vec q_{\rm LL}\) and is
complex conjugated); without gauge fields it is the matrix
\(\mathsf H_{\rm LL}(\vec q_{\rm LL})\) of the registry.

\hypertarget{the-checks-gs-1-to-gs-16}{%
\subsection{12.2 The checks GS-1 to
GS-16}\label{the-checks-gs-1-to-gs-16}}

The constants of the program: \(g^{\rm v}_{\rm LL}=0.7\),
\(g^{\rm F}_{\rm LL}=0.45\), \(m_{\rm LL}=0.4\), \(\mu_{\rm LL}=0.3\)
(from GS-8 on). ``Random'' means drawn from the generator of section
10.3 in the order of the program, uniform in \([-1,1)\) unless another
range is given. The tolerances are those of the original program; the
tolerance of an exact identity is \(10^{-12}\).

\textbf{GS-1 (the Dirac matrices of the frame).} For the canonical
frame: \(M_{\rm LL}(\mathbf f_p)=\sigma^{\rm LL}_p\),
\(M_{\rm LL}(\bar{\mathbf f}_0)=\sigma^{\rm LL}_0\),
\(M_{\rm LL}(\bar{\mathbf f}_j)=-\sigma^{\rm LL}_j\),
\(\mathsf E_{\rm LL}=M_{\rm LL}(-\mathbf e_2)\); the Clifford relations
\(\gamma^p_{\rm LL}\gamma^{p'}_{\rm LL}+\gamma^{p'}_{\rm LL}\gamma^p_{\rm LL}=2\eta^{\rm LL}_{pp'}\mathsf 1_4\);
\(\gamma^{\rm ch}_{\rm LL}=\mathrm{diag}(-1,-1,1,1)\);
\(C_{\rm LL}=\mathrm{diag}(\mathsf E_{\rm LL},-\mathsf E_{\rm LL})\);
\(\alpha^j_{\rm LL}=\mathrm{diag}(-\sigma^{\rm LL}_j,\sigma^{\rm LL}_j)\)
(\(2\times2\) blocks). Result: the largest deviation of the Clifford
relations is 0 (exact integers; original 0).

\textbf{GS-2 (the Lorentz generators on the Dirac spinor).}
\(\mathsf s_{\rm LL}(J^j_{\rm LL})=\mathrm{diag}(-\tfrac i2\sigma^{\rm LL}_j,-\tfrac i2\sigma^{\rm LL}_j)=\tfrac12\gamma^{j'}_{\rm LL}\gamma^{j''}_{\rm LL}\)
for cyclic \((j,j',j'')\),
\(\mathsf s_{\rm LL}(K^j_{\rm LL})=\mathrm{diag}(-\tfrac12\sigma^{\rm LL}_j,\tfrac12\sigma^{\rm LL}_j)=\tfrac12\gamma^0_{\rm LL}\gamma^j_{\rm LL}\);
the 27 relations of \(\mathfrak{so}(3,1)\),
\[[\mathsf s(J^j),\mathsf s(J^{j'})]=\sum_{j''}\epsilon^{\rm LL}_{jj'j''}\mathsf s(J^{j''}),\]
\[[\mathsf s(J^j),\mathsf s(K^{j'})]=\sum_{j''}\epsilon^{\rm LL}_{jj'j''}\mathsf s(K^{j''}),\]
\[[\mathsf s(K^j),\mathsf s(K^{j'})]=-\sum_{j''}\epsilon^{\rm LL}_{jj'j''}\mathsf s(J^{j''})\]
(\(\mathsf s\) for \(\mathsf s_{\rm LL}\), the subscripts LL of the
generators omitted; 9 pairs \((j,j')\), three relations each); and the
24 intertwining relations
\([\gamma^p_{\rm LL},\mathsf s_{\rm LL}(\mathbf u)]=\sum_{p'}\mathcal V_{\rm LL}(\mathbf u)^p{}_{p'}\gamma^{p'}_{\rm LL}\)
for the six generators and the four values of \(p\). Result: all 27 and
all 24 hold within \(10^{-12}\) (original: 27 and 24).

\textbf{GS-3 (the gauge group as a change of representation).} For 200
random \(\mathbf g\):
\(\mathsf S_{\rm LL}(\mathbf g)\mathsf S_{\rm LL}(\mathbf g)^{-1}=\mathsf 1_4\)
with
\(\mathsf S_{\rm LL}(\mathbf g)^{-1}=\mathrm{diag}(M_{\rm LL}(\mathbf g)^\dagger,M_{\rm LL}(\mathbf g)^{-1})\);
\(\mathsf S_{\rm LL}(\mathbf g)^{-1}\gamma^p_{\rm LL}\mathsf S_{\rm LL}(\mathbf g)=\sum_{p'}[\mathsf L_{\rm LL}(\mathbf g)]_{pp'}\gamma^{p'}_{\rm LL}\);
and, for four random Hermitian bi-quaternions
\(\mathbf f_0,\dots,\mathbf f_3\), each with four random real
coordinates (the relation holds for any four Hermitian bi-quaternions,
not only for orthonormal frames),
\(\mathsf S_{\rm LL}(\mathbf g)\gamma^p_{\rm LL}[\mathbf f]\mathsf S_{\rm LL}(\mathbf g)^{-1}=\gamma^p_{\rm LL}[\tilde{\mathbf g}\mathbf f\tilde{\mathbf g}^\dagger]\),
where the transformed bi-quaternions are computed from the coordinates
of
\(M_{\rm LL}(\tilde{\mathbf g})M_{\rm LL}(\mathbf f_p)M_{\rm LL}(\tilde{\mathbf g})^\dagger\).
Tolerance \(10^{-10}\). Result: \(5.13\times10^{-14}\) and
\(7.21\times10^{-14}\) (original \(5.8\times10^{-14}\) and
\(6.6\times10^{-14}\)).

\textbf{GS-4 (isospin and the pair matrices).}
\([T^\iota_{\rm LL},T^{\iota'}_{\rm LL}]=i\sum_{\iota''}\epsilon^{\rm LL}_{\iota\iota'\iota''}T^{\iota''}_{\rm LL}\);
\(T^\iota_{\rm LL}\) commutes with
\(\mathsf s_{\rm LL}(J^j_{\rm LL})\otimes\sigma^{\rm LL}_0\) and
\(\mathsf s_{\rm LL}(K^j_{\rm LL})\otimes\sigma^{\rm LL}_0\);
\((\mathsf K^\iota_{\rm LL})^T=-\mathsf K^\iota_{\rm LL}\); for 100
random \(\mathbf g\) and \(\mathsf U_{\rm LL}\) (angles
\(\theta^\iota_{\rm LL}\) uniform in \([-3,3)\)):
\((\mathsf S_{\rm LL}(\mathbf g)\otimes\sigma^{\rm LL}_0)^T\mathsf K^\iota_{\rm LL}(\mathsf S_{\rm LL}(\mathbf g)\otimes\sigma^{\rm LL}_0)=\mathsf K^\iota_{\rm LL}\)
(Lorentz invariance; the transpose appears because the pair operator is
\(\Psi^T_{\rm LL}\mathsf K^\iota_{\rm LL}\Psi_{\rm LL}\)) and
\((\mathsf 1_4\otimes\mathsf U_{\rm LL})^T\mathsf K^\iota_{\rm LL}(\mathsf 1_4\otimes\mathsf U_{\rm LL})=\sum_{\iota'}[\mathsf R^{\rm v}_{\rm LL}(\mathsf U_{\rm LL})]_{\iota\iota'}\mathsf K^{\iota'}_{\rm LL}\)
(an isospin vector); and
\(\Psi^T_{\rm LL}\mathsf K^\iota_{\rm LL}\Psi_{\rm LL}=0\) for a random
commuting (c-number) \(\Psi_{\rm LL}\), because the matrix is
antisymmetric. Tolerances \(10^{-10}\) and \(10^{-12}\). Result:
\(5.59\times10^{-15}\) and \(2.99\times10^{-16}\) (original
\(5.2\times10^{-15}\) and \(3.0\times10^{-16}\)).

\textbf{GS-5 (gauge covariance of the covariant derivative).} Random
fields linear in \(x_{\rm LL}\): \(A^\iota_{{\rm LL},p}\) and
\(A^{\rm F}_{{\rm LL},p}\) (constant parts in \([-1,1)\), slopes in
\([-0.5,0.5)\)), \(\omega^j_{{\rm LL},p}\) and \(\varpi^j_{{\rm LL},p}\)
(constant parts in \([-0.5,0.5)\), slopes in \([-0.3,0.3)\)); a spinor
field made of three plane waves (\texttt{waves}: four random numbers and
a random 8-component amplitude each); a gauge transformation
\(\mathcal U_{\rm LL}\) whose \(\theta^{\rm F}_{\rm LL}\) and
\(\theta^\iota_{\rm LL}\) are linear in \(x_{\rm LL}\) (coefficients in
\([-1,1)\)) and whose \(\mathbf g\) has
\(M_{\rm LL}(\mathbf g)=\exp M_{\rm LL}(\mathbf u)\) with a Lorentz
generator \(\mathbf u\) linear in \(x_{\rm LL}\) (coefficients in
\([-0.5,0.5)\)). The transformed fields are
\(A^{\rm F}_{{\rm LL},p}+\partial_p\theta^{\rm F}_{\rm LL}/g^{\rm F}_{\rm LL}\),
\(\mathsf U_{\rm LL}\mathsf a_{{\rm LL},p}\mathsf U_{\rm LL}^{-1}-(i/g^{\rm v}_{\rm LL})(\partial_p\mathsf U_{\rm LL})\mathsf U_{\rm LL}^{-1}\)
and
\(\mathsf S_{\rm LL}(\mathbf g)\mathsf s_{\rm LL}(\mathbf u^\omega_{{\rm LL},p})\mathsf S_{\rm LL}(\mathbf g)^{-1}-(\partial_p\mathsf S_{\rm LL}(\mathbf g))\mathsf S_{\rm LL}(\mathbf g)^{-1}\).
At 20 random points and for every \(p\), the covariant derivative made
of the transformed fields applied to
\(\mathcal U_{\rm LL}\Psi_{\rm LL}\) is compared with
\(\mathcal U_{\rm LL}\nabla^{\rm LL}_p\Psi_{\rm LL}\); the derivatives
of the gauge functions and of \(\mathcal U_{\rm LL}\Psi_{\rm LL}\) are
central differences with the step \(10^{-5}\), so the comparison is
limited by the error of the differences. The relative deviation
\(\lvert{\rm lhs}-{\rm rhs}\rvert/(1+\lvert{\rm rhs}\rvert)\) must stay
below \(10^{-7}\), and the transformed isospin potential must be
Hermitian within \(10^{-9}\). Result: \(1.45\times10^{-10}\) (original
\(1.4\times10^{-10}\)).

\textbf{GS-6 (the curvature).} At 20 random points and for all 12
ordered pairs \(p\neq p'\):
\([\nabla^{\rm LL}_p,\nabla^{\rm LL}_{p'}]\Psi_{\rm LL}=\big(\mathsf R^{\rm LL}_{pp'}-ig^{\rm v}_{\rm LL}\sum_\iota\mathsf F^\iota_{{\rm LL},pp'}T^\iota_{\rm LL}-ig^{\rm F}_{\rm LL}\mathsf F^{\rm F}_{{\rm LL},pp'}\big)\Psi_{\rm LL}\),
with the fields of GS-5; the outer derivative of the left side is a
central difference with the step \(10^{-4}\), the derivatives of the
linear fields on the right side are exact. Relative tolerance
\(10^{-6}\). Result: \(2.14\times10^{-9}\) (original
\(2.1\times10^{-9}\)).

\textbf{GS-7 (the super-loop form of the Dirac operator).} For 50 random
plane waves with four random numbers \texttt{kv} and random constant
fields, \(\nabla^{\rm LL}_p\) acts as the matrix
\(i\)\texttt{kv}\(_p\mathsf 1_8-i\mathsf A_{{\rm LL},p}\). Checked:
\(\sum_p(\gamma^p_{\rm LL}\otimes\sigma^{\rm LL}_0)\nabla^{\rm LL}_p\)
has no diagonal blocks and its off-diagonal blocks are
\(M_{\rm LL}(\mathbf D_{\rm LL})\) and
\(M_{\rm LL}(\bar{\mathbf D}_{\rm LL})\);
\(\big(\sum_p\gamma^p_{\rm LL}\nabla^{\rm LL}_p\big)^2=\sum_p\eta^{\rm LL}_{pp}\nabla^{\rm LL}_p\nabla^{\rm LL}_p+\tfrac14\sum_{p,p'}[\gamma^p_{\rm LL},\gamma^{p'}_{\rm LL}][\nabla^{\rm LL}_p,\nabla^{\rm LL}_{p'}]\);
and for a free plane wave
\(M_{\rm LL}(\mathbf D_{\rm LL})M_{\rm LL}(\bar{\mathbf D}_{\rm LL})=-N_{\rm LL}\,\sigma^{\rm LL}_0\),
where \(N_{\rm LL}\) is the norm of the Hermitian bi-quaternion with the
coordinates \texttt{kv}, the first of the four numbers squared minus the
squares of the other three (the product of the two \(2\times2\) matrices
is this difference of squares times \(\sigma^{\rm LL}_0\), because the
Pauli matrices anticommute with each other and square to
\(\sigma^{\rm LL}_0\)). Result: \(1.52\times10^{-15}\) (original
\(1.6\times10^{-15}\)).

\textbf{GS-8 (the Dirac-Bogoliubov-de Gennes matrix).} For 50 random
momenta \(\vec q_{\rm LL}\in[-2,2)^3\), random constant fields and
random complex \(\vec d_{\rm LL}\):
\(\mathcal H^{\rm DB}_{\rm LL}(\vec q_{\rm LL})\) is Hermitian, and it
is particle-hole symmetric,
\(\mathsf X^{\rm ph}_{\rm LL}\mathcal H^{\rm DB}_{\rm LL}(\vec q_{\rm LL})^*\mathsf X^{\rm ph}_{\rm LL}=-\mathcal H^{\rm DB}_{\rm LL}(-\vec q_{\rm LL})\).
For the polar order parameter \(\vec d_{\rm LL}=(0,0,\Delta_{\rm LL})\)
with \(\Delta_{\rm LL}=0.25+0.2\times{}\)(a random number in
\([-1,1)\)), so \(\Delta_{\rm LL}\in[0.05,0.45)\), and without gauge
fields, the 16 eigenvalues (from \texttt{t2eigh}) are
\(\pm E^+_{\rm LL}\) and \(\pm E^-_{\rm LL}\), each four times (30
random momenta). Tolerance \(10^{-10}\). Result: 0 and
\(7.99\times10^{-15}\) (original 0 and \(6.7\times10^{-15}\)).

\textbf{GS-9 (the self-consistency).} For the polar state of GS-8 and 30
random momenta, the sum over the eight eigenvectors with positive
energy,
\(\sum_{n:E^{\rm qp}_{{\rm LL},n}>0}(\psi^{\rm p}_{{\rm LL},n})^T\mathsf K^3_{\rm LL}(\psi^{\rm h}_{{\rm LL},n})^*=-2\Delta_{\rm LL}\big(1/E^+_{\rm LL}+1/E^-_{\rm LL}\big)\).
With \(d^3_{\rm LL}=-\tfrac12G_{\rm LL}\langle\Pi^3_{\rm LL}\rangle\)
and \(\langle\Pi^3_{\rm LL}\rangle\) the sum of this expression over the
momenta,
\(\Delta_{\rm LL}=G_{\rm LL}\Delta_{\rm LL}\sum_{\vec q_{\rm LL}}(1/E^+_{\rm LL}+1/E^-_{\rm LL})\),
which per volume is the gap equation
\(1/G_{\rm LL}=\int d^3q_{\rm LL}(2\pi)^{-3}(1/E^+_{\rm LL}+1/E^-_{\rm LL})\).
Tolerance \(10^{-10}\). Result: \(1.34\times10^{-15}\) (original
\(1.7\times10^{-15}\)).

\textbf{GS-10 (the lattice operator).} A lattice of \(3\times2\times4\)
sites, \(\ell^{\rm x}_{\rm LL}=0.7\), periodic in the directions 1 and
3, Dirichlet in the direction 2, \(r^{\rm W}_{\rm LL}=1\), with random
links, potentials and order parameters on every site; the dense matrix
of \(\mathcal H^{\rm lat}_{\rm LL}\) has \(16\cdot24=384\) rows.
Checked: Hermiticity; particle-hole symmetry,
\(\mathsf X(\mathcal H^{\rm lat}_{\rm LL})^*\mathsf X=-\mathcal H^{\rm lat}_{\rm LL}\)
with \(\mathsf X\) the exchange \(\mathsf X^{\rm ph}_{\rm LL}\) at every
site; exact gauge covariance under random site transformations
\(\mathsf G^{\rm site}_{\rm LL}\) with
\(\exp(i\theta^{\rm F}_{\rm LL})\mathsf U_{\rm LL}\) at each site
(angles in \([-3,3)\)): links
\(\mathsf G\,\mathsf U^{\rm link}\,\mathsf G(\vec s_{\rm LL}+\vec\delta^j_{\rm LL})^\dagger\)
(the unit matrix in place of the second factor at a Dirichlet boundary),
potentials \(\mathsf G\,(\cdot)\,\mathsf G^\dagger\), order parameters
\(\exp(2i\theta^{\rm F}_{\rm LL})\mathsf R^{\rm v}_{\rm LL}(\mathsf U_{\rm LL})\vec d_{\rm LL}\),
and the operator transformed by the block-diagonal matrix of
\(\mathsf 1_4\otimes\mathsf G\) (particles) and its complex conjugate
(holes); on a free periodic chain of 12 sites along \(x^3_{\rm LL}\)
(\(\ell^{\rm x}_{\rm LL}=0.5\), \(\mu_{\rm LL}=0\)) the spectrum is
\(\pm\varepsilon^{\rm lat}_{\rm LL}\), each eight times, at the momenta
\(q^3_{\rm LL}=2\pi\kappa/(12\cdot0.5)\), \(\kappa=0,\dots,11\); the
Wilson term lifts the doubler,
\(\varepsilon^{\rm lat}_{\rm LL}=m_{\rm LL}+2r^{\rm W}_{\rm LL}/\ell^{\rm x}_{\rm LL}\)
at \(q^3_{\rm LL}\ell^{\rm x}_{\rm LL}=\pi\); and the spectral radius
stays below \(B^{\rm H}_{\rm LL}\). Tolerance \(10^{-10}\). Result: 0,
0, \(1.75\times10^{-13}\), \(2.84\times10^{-14}\); spectral radius
\(8.4349\le B^{\rm H}_{\rm LL}=17.5474\) (original 0, 0,
\(1.8\times10^{-13}\), \(1.4\times10^{-14}\); \(8.4349\le17.5474\)).

\textbf{GS-11 (the eigensolver).} For random Hermitian matrices of the
dimensions 7, 64 and 200 (real diagonal and complex off-diagonal
entries, drawn row by row), the larger of the relative residual of the
reconstruction and the deviation of the eigenvectors from orthonormality
stays below \(10^{-11}\). Result: \(6.49\times10^{-14}\) (original
\(6.0\times10^{-14}\)).

\textbf{GS-12 (the printed formulas).} The formulas that the original
program prints are built from symbols of the Wolfram Language (ordinary
symbols such as \texttt{dS{[}i{]}} for \(d^\iota_{\rm LL}\), with
formatting rules so that they print as the registered symbols), shown,
and evaluated for numbers assigned to the symbols; each value is
compared with the matrix computation: (a) \(\mathsf P_{\rm LL}\) with
symbolic \(d^\iota_{\rm LL}\) (block-diagonal with two equal
\(4\times4\) blocks; values
\(\vec d_{\rm LL}=(0.31-0.2i,-0.17+0.4i,0.52+0.11i)\)); (b)
\(\mathsf s_{\rm LL}(\mathbf u^\omega_{{\rm LL},p})\) with symbolic
\(\omega^j_{{\rm LL},p}\), \(\varpi^j_{{\rm LL},p}\), and its blocks
\(-M_{\rm LL}(\mathbf u^\omega_{{\rm LL},p})^\dagger\) and
\(M_{\rm LL}(\mathbf u^\omega_{{\rm LL},p})\); (c) the \(8\times8\)
matrix \(\Gamma^{\rm LL}_p\) and the gauge matrix
\(\mathsf A_{{\rm LL},p}\); (d) the component expansions of
\(\nabla^{\rm LL}_j\psi^{\rm p}_{\rm LL}\),
\(\nabla^{{\rm LL}*}_j\psi^{\rm h}_{\rm LL}\) and
\(\nabla^{\rm LL}_j\Psi_{\rm LL}\) for a plane wave; (e) the sixteen
component equations \(i\partial_0\psi^{\rm p}_{{\rm LL},k}=\dots\),
\(i\partial_0\psi^{\rm h}_{{\rm LL},k}=\dots\) of
\(\mathcal H^{\rm DB}_{\rm LL}\), compared with the \(16\times16\)
matrix applied to a random vector; (f)
\(\Pi^\iota_{\rm LL}=\sum_{k<k'}\big((\mathsf K^\iota_{\rm LL})_{kk'}-(\mathsf K^\iota_{\rm LL})_{k'k}\big)\Psi_{{\rm LL},k}\Psi_{{\rm LL},k'}\)
(only \(k<k'\), because the components anticommute); (g) the eight
component equations of the operator equation
\(i\partial_0\Psi_{\rm LL}=\big(m_{\rm LL}\beta_{\rm LL}-i\Gamma^{\rm LL}_0-\mathsf A_{{\rm LL},0}-\mu_{\rm LL}\big)\Psi_{\rm LL}-i\sum_j\alpha^j_{\rm LL}\nabla^{\rm LL}_j\Psi_{\rm LL}-\tfrac12G_{\rm LL}\sum_\iota(\mathsf K^\iota_{\rm LL})^\dagger\Psi^*_{\rm LL}\Pi^\iota_{\rm LL}\),
with independent random values for \(\Psi_{\rm LL}\),
\(\Psi^*_{\rm LL}\), \(\Pi^\iota_{\rm LL}\),
\(\nabla^{\rm LL}_j\Psi_{\rm LL}\) and \(G_{\rm LL}=1.3\). The notebook
also shows the Weyl blocks \(M_{\rm LL}(\mathbf D_{\rm LL})\) and
\(M_{\rm LL}(\bar{\mathbf D}_{\rm LL})\) as \(2\times2\) matrices of
operators. Tolerance \(10^{-12}\). Result: 0, \(3.6\times10^{-16}\),
\(1.6\times10^{-16}\), \(1.1\times10^{-16}\), \(2.3\times10^{-16}\) for
(a)-(c), (e), (d), (f), (g) (original 0, \(4.0\times10^{-16}\),
\(3.4\times10^{-16}\), 0, \(5.0\times10^{-16}\)).

\textbf{GS-13 (the time step).} One RK4 step multiplies the component of
an eigenvector with the energy \(E^{\rm qp}_{{\rm LL},n}\) by
\(1-iz-z^2/2+iz^3/6+z^4/24\),
\(z=E^{\rm qp}_{{\rm LL},n}\ell^{\rm t}_{\rm LL}\) (the Taylor
polynomial of degree 4 of \(e^{-iz}\)). Its squared modulus is
\((1-z^2/2+z^4/24)^2+(z-z^3/6)^2=1-z^6/72+z^8/576\), so the modulus is
at most 1 exactly for \(z^2\le8\), that is \(\lvert z\rvert\le2\sqrt2\),
and \(1-\lvert\cdot\rvert=z^6/144+O(z^8)\) for small \(z\). The solvers
use \(\ell^{\rm t}_{\rm LL}\le2.5/B^{\rm H}_{\rm LL}\). The printed
check line keeps the wording of the original report, in which this
abbreviation is called x. Result: modulus 0.508186 at
\(\lvert z\rvert=2.5\) and 1.193063 at 2.9; \(1-{}\)modulus at \(z=0.1\)
is \(6.93576\times10^{-9}\), within \(10^{-11}\) of \(0.1^6/144\)
(original 0.508186; 1.193063; \(6.9\times10^{-9}\)).

\textbf{GS-14 (the super-loop form of the field equations).} With random
bi-quaternions: (a)
\(M_{\rm LL}(\mathbf x)\psi'_{\rm LL}(\mathbf y)=\psi'_{\rm LL}(\mathbf y\bar{\mathbf x})\),
and the Cayley-Dickson product gives
\((\mathbf y\mathbf e_4)\mathbf x=(\mathbf y\bar{\mathbf x})\mathbf e_4\)
(50 trials); (b)
\(\gamma^p_{\rm LL}[\mathbf f]\Psi_{\rm LL}[\mathbf b^-_{\rm LL},\mathbf b^+_{\rm LL}]=\Psi_{\rm LL}[\mathbf b^+_{\rm LL}\bar{\mathbf f}_p,\mathbf b^-_{\rm LL}\mathbf f_p]\)
for the canonical frame and 19 sets of random Hermitian bi-quaternions;
(c) \(\mathcal U_{\rm LL}\) with
\(\mathsf U_{\rm LL}=M_{\rm LL}(\mathbf v)\) replaces
\(\mathbf b^+_{\rm LL}\) by
\(e^{i\theta^{\rm F}_{\rm LL}}\mathbf v\mathbf b^+_{\rm LL}\bar{\mathbf g}\)
and \(\mathbf b^-_{\rm LL}\) by
\(e^{i\theta^{\rm F}_{\rm LL}}\mathbf v\mathbf b^-_{\rm LL}\mathbf g^\dagger\);
\(\mathbf v\) is a real quaternion and \(\tilde{\mathbf g}=\mathbf g^*\)
(50 trials); (d) \(\nabla^{\rm LL}_p\) acts on \(\mathbf b^+_{\rm LL}\)
as
\(\partial_p\mathbf b^+_{\rm LL}-\mathbf b^+_{\rm LL}\mathbf u^\omega_{{\rm LL},p}+g^{\rm v}_{\rm LL}\mathbf t^{\rm A}_{{\rm LL},p}\mathbf b^+_{\rm LL}-ig^{\rm F}_{\rm LL}A^{\rm F}_{{\rm LL},p}\mathbf b^+_{\rm LL}\)
and on \(\mathbf b^-_{\rm LL}\) in the same way with
\(\mathbf u^{\omega*}_{{\rm LL},p}\) (30 trials, four directions each);
(e)
\(\sum_p\gamma^p_{\rm LL}\nabla^{\rm LL}_p\Psi_{\rm LL}=\Psi_{\rm LL}\big[\sum_p(\nabla^{\rm LL}_p\mathbf b^+_{\rm LL})\bar{\mathbf f}_p,\sum_p(\nabla^{\rm LL}_p\mathbf b^-_{\rm LL})\mathbf f_p\big]\),
the right multiplications \(R^{\rm LL}_{\bar{\mathbf f}_p}\),
\(R^{\rm LL}_{\mathbf f_p}\) applied to the covariant derivatives of the
charges. Tolerances \(10^{-12}\) (\(10^{-10}\) for (c)). Result:
\(6.0\times10^{-16}\), \(5.9\times10^{-16}\), \(3.5\times10^{-15}\),
\(9.7\times10^{-16}\), \(2.0\times10^{-15}\) (original
\(3.2\times10^{-16}\), \(4.1\times10^{-16}\), \(3.1\times10^{-15}\),
\(8.8\times10^{-16}\), \(2.1\times10^{-15}\)).

\textbf{GS-15 (the Lorentz connection is fixed by the frame).} Five
random Lorentz generators \(\mathbf u_1,\dots,\mathbf u_5\) define
\(M_{\rm LL}(\mathbf g(x_{\rm LL}))=\exp\big(M_{\rm LL}(\mathbf u_5)+\sum_{p=0}^3x^p_{\rm LL}M_{\rm LL}(\mathbf u_{p+1})\big)\)
and the connection
\(\Gamma_p=-(\partial_p\mathsf S_{\rm LL}(\mathbf g))\mathsf S_{\rm LL}(\mathbf g)^{-1}\)
(a \(4\times4\) matrix; central differences with the step \(10^{-4}\)).
At 20 random points: \(\Gamma_p\) has the form
\(\mathsf s_{\rm LL}(\mathbf u)\) (no off-diagonal blocks, upper block
\(=-\)(lower block)\(^\dagger\), traceless lower block); the Dirac
matrices of the transformed frame,
\(\mathsf S_{\rm LL}(\mathbf g)\gamma^{p'}_{\rm LL}\mathsf S_{\rm LL}(\mathbf g)^{-1}\),
are covariantly constant, \(\partial_p(\cdot)+[\Gamma_p,\cdot]=0\); the
curvature
\(\partial_p\Gamma_{p'}-\partial_{p'}\Gamma_p+[\Gamma_p,\Gamma_{p'}]\)
vanishes. Uniqueness: the real-linear map
\(\mathbf u\mapsto\big([\mathsf s_{\rm LL}(\mathbf u),\gamma^{p'}_{\rm LL}]\big)_{p'=0,\dots,3}\)
has the rank 6 on the six generators \(J^j_{\rm LL}\), \(K^j_{\rm LL}\):
the smallest eigenvalue of the \(6\times6\) Gram matrix of the real and
imaginary parts of all entries of the commutators exceeds \(10^{-6}\).
Tolerance \(10^{-6}\) (limited by the finite differences). Result:
\(1.03\times10^{-9}\), \(3.55\times10^{-8}\), \(7.98\times10^{-9}\);
smallest Gram eigenvalue 8.000 (original \(1.0\times10^{-9}\),
\(3.6\times10^{-8}\), \(8.0\times10^{-9}\); 8.000). (The first
evaluation failed this check because of error E15.)

\textbf{GS-16 (the lattice operator on plane waves).} On a periodic
chain of 24 sites along \(x^1_{\rm LL}\) (\(\ell^{\rm x}_{\rm LL}=0.5\),
\(m_{\rm LL}=0.4\), \(\mu_{\rm LL}=0\), no gauge fields, no pairing),
\(\mathcal H^{\rm lat}_{\rm LL}\), applied by \texttt{t2apply} to
\(\exp(iq^1_{\rm LL}x^1_{\rm LL})\) times each of the 16 unit vectors of
a site, gives a plane wave of the same momentum, and the \(16\times16\)
matrix of the images has the eigenvalues
\(\pm\varepsilon^{\rm lat}_{\rm LL}\), each eight times, at all 24
momenta \(q^1_{\rm LL}=2\pi\kappa/(24\ell^{\rm x}_{\rm LL})\),
\(\kappa=-11,\dots,12\), for \(r^{\rm W}_{\rm LL}=0\) and \(1\), where

\[\varepsilon^{\rm lat}_{\rm LL}(\vec q_{\rm LL})=\Big(\sum_j\frac{\sin^2(q^j_{\rm LL}\ell^{\rm x}_{\rm LL})}{(\ell^{\rm x}_{\rm LL})^2}+\Big(m_{\rm LL}+\frac{r^{\rm W}_{\rm LL}}{\ell^{\rm x}_{\rm LL}}\sum_j\big(1-\cos(q^j_{\rm LL}\ell^{\rm x}_{\rm LL})\big)\Big)^2\Big)^{1/2}\]

(sums over the directions with more than one site). Tolerance
\(10^{-12}\). Result: \(2.73\times10^{-14}\) (original
\(1.4\times10^{-14}\)).

\textbf{T2-DATA-1 (the data of the figures).} The original program wrote
the numbers of its two figures into three files of
\texttt{Revision\_SuperLoop/test2/data}:
\texttt{t2\_dispersion\_lattice.csv} (48 rows: \(r^{\rm W}_{\rm LL}\),
\(q^1_{\rm LL}\), the largest eigenvalue of GS-16),
\texttt{t2\_dispersion\_curves.csv} (802 rows: \(r^{\rm W}_{\rm LL}\),
\(q^1_{\rm LL}\), \(\varepsilon_{\rm LL}\),
\(\varepsilon^{\rm lat}_{\rm LL}\) on 401 momenta) and
\texttt{t2\_rk4\_stability.csv} (641 rows:
\(\lvert z\rvert=0,0.005,\dots,3.2\) and the modulus of the RK4 factor),
printed with 10 decimals (12 for the RK4 factor). Every number of the
three files is copied into the notebook (file
\texttt{notebook\_source/t2\_fields\_figure\_data.wl}, written by
\texttt{tools/make\_fig\_data\_cell.py} from the files, included in one
input cell), and the check compares them, entry by entry, with the
numbers computed in the notebook. Tolerances: half a unit of the last
printed decimal, \(5.1\times10^{-11}\) and \(5.1\times10^{-13}\).
Result: \(4.55\times10^{-11}\), \(5.00\times10^{-11}\),
\(4.99\times10^{-13}\), as expected from the rounding of the files; so
the rewritten programs compute the numbers of the original program.

\hypertarget{the-figures-of-item-1}{%
\subsection{12.3 The figures of item
{[}1{]}}\label{the-figures-of-item-1}}

\textbf{Figure 1} (data: \texttt{points} and \texttt{curves} of GS-16).
Panel (a): the energy of a free particle on the lattice as a function of
\(q^1_{\rm LL}\ell^{\rm x}_{\rm LL}\in[-\pi,\pi]\) for
\(m_{\rm LL}=0.4\), \(\ell^{\rm x}_{\rm LL}=0.5\): the continuum energy
\(\varepsilon_{\rm LL}\) (grey, dashed),
\(\varepsilon^{\rm lat}_{\rm LL}\) for \(r^{\rm W}_{\rm LL}=0\) (orange)
and for \(r^{\rm W}_{\rm LL}=1\) (blue), and, as open circles, the
largest eigenvalues of the lattice operator at the 24 allowed momenta,
which lie on the curves. Without the Wilson term the energy returns to
\(m_{\rm LL}\) at \(q^1_{\rm LL}\ell^{\rm x}_{\rm LL}=\pm\pi\) (a
second, spurious particle, the doubler); with \(r^{\rm W}_{\rm LL}=1\)
the doubler is lifted to
\(m_{\rm LL}+2r^{\rm W}_{\rm LL}/\ell^{\rm x}_{\rm LL}=4.4\). Panel (b):
the relative deviation
\((\varepsilon^{\rm lat}_{\rm LL}-\varepsilon_{\rm LL})/\varepsilon_{\rm LL}\)
for \(\lvert q^1_{\rm LL}\ell^{\rm x}_{\rm LL}\rvert\le\pi/2\): it is of
second order in \(q^1_{\rm LL}\ell^{\rm x}_{\rm LL}\), negative without
the Wilson term and positive (at most 0.07) with it.

\textbf{Figure 2} (data: \texttt{rkData} of GS-13). Panel (a): the
modulus of the RK4 factor against \(\lvert z\rvert\), with the line of
modulus 1 (grey, dashed), the stability limit
\(\lvert z\rvert=2\sqrt2=2.83\) (red, dotted) and the largest value used
by the solvers, \(\lvert z\rvert=2.5\) (green, dash-dotted); the modulus
has its minimum 0.5 near \(\lvert z\rvert=2.45\) and exceeds 1 beyond
the stability limit. Panel (b): on logarithmic axes, \(1-{}\)modulus
(blue) and \(z^6/144\) (orange, dashed) for \(0.02<\lvert z\rvert<1.5\);
the two agree, which shows the small loss of norm per step.

\textbf{How the figures are drawn.} The figures are made with
\texttt{ListLinePlot}, \texttt{ListPlot}, \texttt{ListLogLogPlot} and
\texttt{Show} from the computed lists, so every point of a figure is a
number of the notebook (and, by T2-DATA-1, of the original data files).
The colors are those of the project's figures: blue \texttt{\#2a78d6}
and orange \texttt{\#eb6834} for the two series, grey \texttt{\#8a8984}
for the continuum and the line of modulus 1, a muted grey
\texttt{\#5f5e5a} for the text of labels. The two series colors were
checked for readers with color-vision deficiencies: their color
difference in the perceptual color space OKLab (times 100) is 33.6 for
normal vision and at least 24.7 under simulated protanopia, deuteranopia
and tritanopia (8 is the usual threshold of distinctness), and both have
a contrast of at least 3:1 against the white page; in addition, every
series is named in a legend or in a direct label, so that no information
depends on color alone. Labels are placed in empty regions of the plots
(checked on the exported page, error E16), the panel with a single curve
has no legend, and all labels use the registered symbols, printed with
upright letters as in the cells of the notebook (error E18).

\hypertarget{how-the-notebook-was-built-and-evaluated}{%
\section{13. How the notebook was built and
evaluated}\label{how-the-notebook-was-built-and-evaluated}}

\textbf{The source.} The notebook is written as a text file in the
Wolfram package format,
\texttt{notebook\_source/3+1super-loop.source.wl}: every cell is
introduced by a marker line such as \texttt{(*\ ::Section::\ *)},
\texttt{(*\ ::Text::\ *)} or \texttt{(*\ ::Input::\ *)}. Inside an input
cell a symbol is written as \texttt{LL{[}"tex"{]}} (for example
\texttt{LL{[}"G\_\{\textbackslash{}\textbackslash{}rm\ LL\}"{]}}), and a
line \texttt{(*\ @@INCLUDE\ file@@\ *)} inserts a file: the package
\texttt{sl31\_core.wl}, the registry \texttt{registry\_data.wl}, and the
list \texttt{all\_symbols.wl} of all 422 symbols for test S-14 (written
by \texttt{tools/gen\_all\_symbols.wls}).

\textbf{The source of Part II} is kept in separate files, one or two per
chapter (\texttt{notebook\_source/part2\_A\_core.source.wl},
\texttt{part2\_B\_item1.source.wl},
\texttt{part2\_C\_item1b.source.wl}); a line
\texttt{(*\ @@INCLUDECELLS\ file@@\ *)} at the end of the main source
inserts such a file with its own cell markers. A line
\texttt{(*\ @@INCLUDE\ t2\_fields\_figure\_data.wl@@\ *)} inserts the
numbers of the data files of item {[}1{]} (made by
\texttt{tools/make\_fig\_data\_cell.py}), and the line
\texttt{(*\ @@COMPONENTS@@\ *)} of the first cell of Part II is replaced
by the builder with the list of all components of families typed in the
notebook (section 10.4).

\textbf{The builder} \texttt{tools/build\_notebook.wls} (run by
wolframscript) turns the source into the notebook
\texttt{notebook\_source/3+1super-loop.unevaluated.nb}: the code of
every input cell is parsed into boxes by the front end itself (so
comments, line breaks and indentation are kept), every
\texttt{LL{[}"tex"{]}} is replaced by the two-dimensional box of the
symbol, and text cells receive inline formulas (TeX between dollar
signs, converted by the Wolfram TeX importer, with the alphabets of the
letters kept, error E17), code font and bold. The notebook is written by
the kernel, not through the front end (error E9).

\textbf{The evaluation} is done by the front end of Wolfram.app, exactly
as Evaluation \textgreater{} Evaluate Notebook does it (file
\texttt{tools/guirun.wl}): the notebook is opened (visible), its
evaluator is set to a kernel named \texttt{SL31Run} that is declared for
the running session only
(\texttt{SetOptions{[}\$FrontEndSession,\ EvaluatorNames\ -\textgreater{}\ ...{]}},
so the preferences of Wolfram.app are not changed), and
\texttt{FrontEndTokenExecute{[}nb,\ "EvaluateNotebook"{]}} starts the
evaluation. A hidden one-cell notebook, queued after it in the same
kernel, writes the test list of the kernel to a file and then quits the
kernel, so that every run starts with a fresh kernel. Afterwards the
evaluated notebook is saved (with its evaluator reset to the default),
exported to PDF, and a log lists every message and every printed line
with the label of its input cell
(\texttt{logs/3+1super-loop.evaluation.log.txt}). The commands of
\texttt{guirun.wl} are sent to the kernel of the running Wolfram.app
session through a small file-request mechanism
(\texttt{tools/watcher.wl}, \texttt{tools/handler.wl}): a scheduled task
looks every 3 seconds for request files in a scratch folder.

\textbf{To reproduce} without any of these tools: open
\texttt{notebook\_source/3+1super-loop.unevaluated.nb} (or the evaluated
\texttt{3+1super-loop.nb}) in Wolfram and choose Evaluation
\textgreater{} Evaluate Notebook; the result must be the same 35 PASS
lines and no message. To rebuild the unevaluated notebook from the
source:

\begin{verbatim}
wolframscript -file tools/build_notebook.wls notebook_source/3+1super-loop.source.wl \
    notebook_source/3+1super-loop.unevaluated.nb notebook_source/sl31_core.wl
\end{verbatim}

(the source includes its files relative to its own folder;
\texttt{registry\_data.wl} is made by
\texttt{tools/registry\_to\_wl.wls} from the CSV file of the registry,
and \texttt{all\_symbols.wl} by \texttt{tools/gen\_all\_symbols.wls};
the absolute paths inside these two scripts name the scratch folder of
the author's Mac and must be adapted elsewhere).

(the files of Part II and \texttt{t2\_fields\_figure\_data.wl} stand in
the same folder as the source). The data cell is remade from the data
files of the original program with

\begin{verbatim}
python3 tools/make_fig_data_cell.py Revision_SuperLoop/test2/data \
    3+1super-loop_provenance/notebook_source/t2_fields_figure_data.wl
\end{verbatim}

\textbf{Copies to the author's computer.} The files are written in a
workspace and copied into the scratch folder of the author's Mac before
each evaluation. A copy can deliver an earlier version of a file whose
name was used for a copy shortly before; this happened twice during this
stage (a test script and the builder arrived in their previous
versions). Every copied file is therefore compared with its MD5 digest
on both sides before an evaluation; for the final evaluation of this
stage the digests of all sources, of the package and of the builder
agreed.

\hypertarget{errors-found-their-causes-and-their-fixes}{%
\section{14. Errors found, their causes and their
fixes}\label{errors-found-their-causes-and-their-fixes}}

Every error below was met while the notebook was prepared and evaluated;
each is fixed in the final notebook. E1 to E14 belong to Part I, E15 to
E19 to Part II.

\textbf{E1. The symbolized box was never recognized when the notebook
was evaluated by a program.} Cause: the Notation package takes the
default of the option \texttt{WorkingForm} from
\texttt{AbsoluteCurrentValue{[}InputNotebook{[}{]},\ \{CommonDefaultFormatTypes,\ "Output"\}{]}};
without a selected input notebook this is \texttt{\$Failed}, and the
reading rule was stored for the form \texttt{\$Failed} (found by reading
the source code of the package, \texttt{Notation.m}, in the
installation). Fix: \texttt{WorkingForm\ -\textgreater{}\ StandardForm}
is given explicitly in every call of \texttt{Symbolize}. Evidence: the
tests S-2 and S-14.

\textbf{E2. Results in plain text.}
\texttt{NotebookEvaluate{[}nb,\ InsertResults\ -\textgreater{}\ True{]}},
run by wolframscript, inserted the results as plain-text output form
instead of typeset boxes. A first remedy was a custom cell-by-cell
evaluator (\texttt{tools/run\_notebook.wls}); it was replaced by the
evaluation in the front end of Wolfram.app (E6).

\textbf{E3. Code shown in linear syntax.} Input cells written as plain
strings showed the symbols as linear syntax (for example
\texttt{\textbackslash{}!\textbackslash{}(\textbackslash{}*SubscriptBox{[}...{]}\textbackslash{})}).
Fix: the builder parses the code into boxes with the front end and
inserts the boxes of the symbols.

\textbf{E4. The TeX importer loses alphabets.}
\texttt{ImportString{[}"\$...\$",\ "TeX"{]}} turns
\texttt{\textbackslash{}mathbb}, \texttt{\textbackslash{}mathsf},
\texttt{\textbackslash{}mathfrak} and \texttt{\textbackslash{}mathbf}
letters into plain letters, so different symbols would coincide. Fix:
the notebook's own converter (section 6); the TeX importer is used only
for formulas in text.

\textbf{E5. Equal names for styled letters.} The Notation package builds
names after removing \texttt{StyleBox} and \texttt{TagBox}, so
\(\mathbf e_0\), \(\mathsf e_0\) and \(e_0\) would become one symbol.
Fix: \texttt{SymbolizeRootName} with names that record the alphabet
(section 8). Evidence: S-3.

\textbf{E6. The evaluator run by wolframscript failed.} Two full runs
stopped without a result; a diagnostic run showed that a cell with
several statements on separate lines was read by \texttt{ToExpression}
as ONE product of the statements (the output of a 52-statement cell was
\texttt{Null\^{}52}), and that cells stored as lists of boxes returned
unevaluated \texttt{HoldComplete{[}...{]}} results. In the front end the
statements of a cell are separated at its line breaks before they are
sent to the kernel. Fix: the evaluation was moved into the front end of
Wolfram.app (section 13), which also answered the author's request to
execute the notebook in his Wolfram.app;
\texttt{tools/run\_notebook.wls} is kept, marked as superseded.

\textbf{E7. Nine symbols were not recognized in typed input.} A test
notebook in the front end showed that
\texttt{Symbolize{[}ParsedBoxWrapper{[}SubscriptBox{[}"a",\ "LL,p"{]}{]}{]}}
is not matched by the typed \(a_{{\rm LL},p}\): the sum
\(\sum_{{\rm LL}=1}^2{\rm LL}\,a_{{\rm LL},p}\) gave
\(a_{1,p}+2\,a_{2,p}\) and \texttt{AtomQ} gave \texttt{False}, and the
saved input cell held
\texttt{SubscriptBox{[}"a",\ RowBox{[}\{"LL",\ ",",\ "p"\}{]}{]}}
(folder \texttt{experiments}). To find every affected symbol, a notebook
with one cell for each of the 430 boxes then symbolized was evaluated in
the front end with \texttt{\$PreRead} recording the boxes that the
kernel received (\texttt{tools/canon\_build.wls},
\texttt{tools/canon\_compare.wls}): 9 differed, all strings of roman
text with several tokens (\texttt{"int,F"}, \texttt{"int,B"},
\texttt{"ext,F"}, \texttt{"ext,B"}, \texttt{"v,Y"}, \texttt{"v,4"},
\texttt{"s(2)"}, \texttt{"s(3)"}, \texttt{"s(4)"}), which the front end
had cut into \texttt{RowBox}es
(\texttt{logs/E07\_front\_end\_boxes\_before\_fix.txt}). The front end's
own parser confirmed exactly these 9 among the 213 strings then present
(\texttt{tools/leaf\_check.wls},
\texttt{logs/E07\_leaf\_check\_front\_end\_parser.txt}). Cause: before
evaluation the front end reads every string of the boxes again as typed
text (the reference page of \texttt{StripBoxes} states that the front
end strips boxes ``using a more extensive process'' when it sends them
to the kernel). Fix: the converter cuts roman text into tokens as the
front end does (section 6); after the fix the same comparison found 0
differences (\texttt{logs/E07\_front\_end\_boxes\_after\_fix.txt}). The
test S-0 repeats the check in every evaluation.

\textbf{E8. Three operators were symbolized as letters.} The tensor
product \(\otimes\), the direct sum \(\oplus\) and the composition
\(\circ\) are single characters and were classified as letters before
the test for operators was reached. In an input cell the front end read
them as operators and regrouped the whole line around them. Fix: the
test for operators comes first; these three are operators and correspond
to Wolfram functions (S-13). The number of symbolized TeX items fell
from 430 to 427.

\textbf{E9. The cell of test S-14 was not evaluated.} In the first full
evaluation the input cell that lists all symbols received an input label
but no evaluation: no result, no message, and the next cell received the
same input number. The stored cell showed a regrouped line: the
partial-derivative sign with a subscript, \(\partial_p\), had been read
as an operator applied to everything after it (a \texttt{RowBox} that
begins with a comma followed it). An experiment that split the list into
cells of ten symbols (\texttt{tools/s14\_build.wls},
\texttt{logs/E09\_s14\_chunk\_experiment.log.txt}) located the problem
in the part with \(\otimes\), \(\oplus\), \(\circ\) and \(\partial_p\).
Evaluated alone in cells of the form \texttt{\{b,\ 1\}}, every symbol
box reached the kernel unchanged (\texttt{tools/canon2\_build.wls},
\texttt{tools/canon2\_compare.wls},
\texttt{logs/E09\_operand\_check\_in\_kernel\_input.txt}), and writing
and reading the cells through the front end with \texttt{NotebookPut}
and \texttt{NotebookGet} changed none of them either
(\texttt{tools/reparse\_check.wls},
\texttt{logs/E09\_reparse\_check.txt}); the regrouping had been made
when the builder wrote the notebook file through the front end
(\texttt{Export{[}file,\ notebook,\ "NB"{]}} inside
\texttt{UsingFrontEnd}). Fix (three changes, made together, so that the
evaluation does not separate their effects): (a) the builder writes the
notebook expression with the kernel itself; (b) every symbol whose base
is an operator sign, and every group name, ket and expression, is
wrapped in
\texttt{TagBox{[}box,\ "SL31Name",\ SyntaxForm\ -\textgreater{}\ "symbol"{]}},
which tells the front end to treat the box as one operand wherever it
stands (the reference page of \texttt{SyntaxForm} describes the option;
the value \texttt{"symbol"} for a box that should act as a standalone
thing was recommended by the front-end developer J. Fultz on the
MathGroup mailing list in January 2014); (c) a script on a closing
parenthesis is attached to the whole group, the form that the front end
makes of typed input. After the three changes the built notebook
contains no regrouped line, and the test S-14 passes for all 422
symbols.

\textbf{E10.
\texttt{Clear::ssym:\ 42\ is\ not\ a\ symbol\ or\ a\ valid\ string\ pattern.}}
The demonstration of section I.2 cleared the value of \(a_{\rm LL}\)
with \texttt{ReleaseHold{[}Hold{[}Clear{]}\ @@\ s{]}};
\texttt{Hold{[}Clear{]}{[}a\_LL{]}} evaluates its argument before
\texttt{ReleaseHold} acts, so \texttt{Clear} received 42. Fix:
\texttt{Clear\ @@\ s} and
\texttt{Set\ @@\ Join{[}s,\ HoldComplete{[}42{]}{]}} (section 5.4).

\textbf{E11. Warnings \texttt{Symbolize::bsymbexs}} (``the box structure
attempting to be symbolized has a similar or identical symbol already
defined''). Cause: TeX items with the same box (the four groups of
section 7) were symbolized once each. Fix: the items are gathered by
their boxes (\texttt{GatherBy}), and each box is symbolized once. The
test S-3 was rewritten to compare boxes and symbols (it had compared TeX
items with symbols and reported 430 against 425 as a failure).

\textbf{E12. TeX of the symbols showed the internal names.} A rule
\texttt{Format{[}s,\ TeXForm{]}\ =\ ...} had no effect, because
\texttt{TeXForm} converts the \texttt{TraditionalForm} boxes of an
expression, and the Notation package defines only the
\texttt{StandardForm} boxes. Fix:
\texttt{MakeBoxes{[}s,\ TraditionalForm{]}\ =\ box} for every symbol.
Evidence: S-5 (422 of 422).

\textbf{E13. Plain-text names ``Missing{[}NotAvailable{]}''.} The CSV
import of Wolfram 15 gives \texttt{Missing{[}"NotAvailable"{]}} for an
empty field; \texttt{ToString} made the name
``Missing{[}NotAvailable{]}'' for 14 entries, and the test S-12 reported
it as a name shared by 14 entries. Fix:
\texttt{tools/registry\_to\_wl.wls} turns every field into a string and
an empty field into ``\,'' (no name).

\textbf{E14. Code spans with a context mark.} In text cells the builder
marks code with backquotes; code that itself contains the context mark
\texttt{\textasciigrave{}} (for example
\texttt{SL31\textasciigrave{}Test\textasciigrave{}\$Results}) broke the
code span. Fix: such code is marked with double backquotes.

\textbf{E15. GS-15 failed in the first evaluation of Part II (32 of 33
checks passed).} The connection was written
\texttt{conn15{[}x\_,\ p\_{]}\ :=\ -(Sp\ -\ Sm)\ (0.5/h15)\ .\ Sinv},
where \texttt{Sp} and \texttt{Sm} stand here for the spinor matrices at
the two shifted points and \texttt{Sinv} for the inverse spinor matrix
at the point. In the grammar of the Wolfram Language the matrix product
\texttt{.} (\texttt{Dot}) binds more strongly than the product written
by placing two expressions side by side (\texttt{Times}): in the table
of the operators in order of decreasing precedence (tutorial ``Operator
Input Forms'') \texttt{.} stands above the implicit product. The
expression was therefore read as
\texttt{-(Sp\ -\ Sm)\ *\ ((0.5/h15)\ .\ Sinv)}. \texttt{Dot} of a number
and a matrix is not defined; it stays unevaluated without a message, so
the ``connection'' was a matrix whose entries contained the unevaluated
product, and the comparisons of the check could not give \texttt{True}.
The evaluation gave no message, so only the check itself revealed the
error. Evidence: \texttt{tools/e15\_dot\_precedence.wls} with its output
\texttt{logs/E15\_dot\_precedence.txt}: the parsed form of
\texttt{-(a\ -\ b)\ (s)\ .\ c} is
\texttt{Times{[}Times{[}-1,\ Plus{[}a,\ Times{[}-1,\ b{]}{]}{]},\ Dot{[}s,\ c{]}{]}};
\texttt{0.5\ .\ \{\{1.,\ 2.\},\ \{3.,\ 4.\}\}} stays unevaluated (head
\texttt{Dot}); the faulty form applied to numbers gives a matrix of
unevaluated products; the corrected form
\texttt{(-(0.5)\ (A\ -\ B))\ .\ m} gives the expected matrix. Fix: the
number multiplies the difference first, inside parentheses,
\texttt{conn15{[}x\_,\ p\_{]}\ :=\ (-(0.5/h15)\ (Sp\ -\ Sm))\ .\ Sinv},
with a comment in the code. After the fix GS-15 passes with the
deviations of the original program (section 12.2).

\textbf{E16. Labels of the figures covered curves; a label and the
summary table were cut at the edge of the page.} Found by looking at the
exported PDF (\texttt{logs/E16\_figures\_and\_table\_before\_fix.png}):
in Figure 1(a) a label crossed the vertical axis line at
\(q^1_{\rm LL}=0\), the label of the doubler for
\(r^{\rm W}_{\rm LL}=0\) lay on the orange curve, and the label
\(m_{\rm LL}=0.4\) touched the same curve; in Figure 1(b) the vertical
axis line crossed the legend; in Figure 2(a) the label of the stability
limit was cut at the right edge of the plot and the label of the largest
step lay on the curve; and the summary table of item {[}1{]} had fixed
column widths (81 character widths in all), wider than the page, so it
was cut at the right edge. Fix: no axis lines inside the framed plots
(only the horizontal line of zero deviation in Figure 1(b)); the labels
moved into empty regions and written in the muted grey of the text; the
horizontal range of Figure 2(a) extended to 4, so that the label fits to
the right of its line; no legend for the single curve of Figure 2(a),
which its title names; the summary table without fixed widths, so that
the front end wraps the statements at the width of the page, as in the
table of section I.9. Evidence after the fix:
\texttt{logs/E16\_figures\_and\_table\_after\_fix.png}.

\textbf{E17. Formulas in text cells lost the alphabets of their
letters.} The Wolfram TeX importer, which the builder uses for formulas
in text cells (error E4), turns \texttt{\textbackslash{}mathbf\ b} into
\texttt{StyleBox{[}"b",\ FontSlant\ -\textgreater{}\ "Bold"{]}}, an
invalid option value that the front end ignores, so that the letter is
printed plain; it turns \texttt{\textbackslash{}mathsf\ N} and
\texttt{\textbackslash{}mathbb\ R} into plain letters,
\texttt{\textbackslash{}mathcal\ L} and
\texttt{\textbackslash{}mathfrak\ g} into italic letters, and
\texttt{\textbackslash{}varepsilon} into the lunate epsilon of
\texttt{\textbackslash{}epsilon}. In the texts of Part II the bold
bi-quaternions \(\mathbf b\), \(\mathbf e_p\), \(\mathbf f_p\) were
therefore printed as plain letters and \(\mathbb R^{1,3}_{\rm LL}\) with
a plain R
(\texttt{logs/E17\_text\_formulas\_before\_and\_after\_fix.png}), and
the energy \(\varepsilon_{\rm LL}\) was printed with the character of
the Levi-Civita symbol \(\epsilon^{\rm LL}_{jj'j''}\); the texts of Part
I were affected in the same way. Evidence:
\texttt{tools/e17\_tex\_importer.wls} with its output
\texttt{logs/E17\_tex\_importer.txt}. Fix (in the builder): before the
import every styled letter is replaced by a placeholder word inside
\texttt{\textbackslash{}mathrm\{...\}} (for example
\texttt{\textbackslash{}mathbb\ R} by
\texttt{\textbackslash{}mathrm\{QZbbQRZQ\}}), which the importer keeps
as one string, and \texttt{\textbackslash{}varepsilon} by the
placeholder of the curly epsilon; after the import each placeholder is
replaced by the box that the converter of the notebook uses for the same
letter (section 6): double-struck, script and Gothic characters, bold
and sans-serif upright letters, and the curly epsilon. In the built
notebook the 588 formulas of the text cells then contain 270 bold
letters, 116 sans-serif letters, 76 double-struck, script and Gothic
characters, 12 curly and 8 lunate epsilons, no invalid bold option and
no placeholder (\texttt{logs/E17\_tex\_importer.txt}).

\textbf{E18. Components of families typed in input cells were not
symbols.} In the table of the basic matrices (section II.1.3 of the
notebook) the label \(\gamma^0_{\rm LL}\) was printed as 1 and
\(\gamma^1_{\rm LL}\) as \(\gamma_{\rm LL}\), and the axis label
\(q^1_{\rm LL}\) of Figure 1 lost its superscript
(\texttt{logs/E18\_components\_before\_fix.png}). Cause: the builder
writes the box of a component into the input cell, but the component had
never been symbolized, so the kernel read the box \(\gamma^0_{\rm LL}\)
as the power 0 of the expression \(\gamma_{\rm LL}\), which is 1, and
\(q^1_{\rm LL}\) as the first power of \(q_{\rm LL}\). No check had
caught this, because these results were only displayed. Fix (section
10.4): the builder collects all components typed in the notebook, the
first cell of Part II makes those that do not yet exist, and the check
T2-SYM verifies all of them. An intermediate state showed the internal
name of \(q^1_{\rm LL}\) in the label of the graphics, because labels of
graphics are typeset in \texttt{TraditionalForm}, for which the new
component had no rule (the same cause as E12); the cell now gives every
component its \texttt{TraditionalForm} box and a usage message, and
T2-SYM also checks that the TeX of every component contains no internal
name. Evidence after the fix:
\texttt{logs/E18\_components\_after\_fix.png}.

\textbf{E19. Notations of Part II that were not registered symbols.} The
first version of the texts and the labels of Part II wrote some objects
with the short notations of the original program instead of the
registered symbols: \(\Delta x\) for the lattice spacing
\(\ell^{\rm x}_{\rm LL}\), \(\Delta t\) for the time step
\(\ell^{\rm t}_{\rm LL}\), \(r_W\) for \(r^{\rm W}_{\rm LL}\),
\(U_j(s)\) for \(\mathsf U^{\rm link}_{{\rm LL},j}(\vec s_{\rm LL})\),
\(V_0\) for \(\mathsf A_{{\rm LL},0}\), \(P\) for
\(\mathsf P_{\rm LL}\), \(\xi_p\) for \(\xi^{\rm LL}_p(\mathbf b)\),
\(S\) and \(L\) for \(\mathsf S_{\rm LL}(\mathbf g)\) and
\(\mathsf L_{\rm LL}(\mathbf g)\), \(E\) for
\(E^{\rm qp}_{{\rm LL},n}\), a plain \(x\) for the product of an energy
and the time step, and \(m\), \(\varepsilon\), \(\varepsilon_{\rm lat}\)
without the label LL. A plain \(x\), for example, could be taken for the
coordinates \(x^p_{\rm LL}\). This contradicts the rule of the project
that every object has one unique, registered symbol. Fix: all texts and
all labels of Part II use the registered symbols; in the labels of
graphics they are the symbolized symbols themselves, printed with
upright letters, as in the cells of the notebook (option
\texttt{SingleLetterItalics\ -\textgreater{}\ False}; the reference page
of the option states that it is typically \texttt{True} in cells of
\texttt{TraditionalForm}); the one abbreviation that is not registered,
\(z=E^{\rm qp}_{{\rm LL},n}\ell^{\rm t}_{\rm LL}\), is defined where it
is used (section 2.6). The printed lines of the checks keep the
plain-text wording of the original report, whose names (for example
\texttt{dx}, \texttt{r\_W}, \texttt{Sspin}, \texttt{Kch}) are the
plain-text names of the registry; only the abbreviation x of GS-13 is
not a registered name, and section 12.2 states its meaning.

\hypertarget{decisions}{%
\section{15. Decisions}\label{decisions}}

\begin{itemize}
\tightlist
\item
  ``./'' of the task is the top folder of the repository, so the
  notebook is \texttt{3+1super-loop.nb} there; ``./test'' is the
  author's folder of the Test 2 documents, whose symbols are the
  complete registry (section 7). The author's folders \texttt{test} and
  \texttt{mathematica\_symbols} are private and are not part of the
  repository.
\item
  Function symbols are symbolized by their heads:
  \(W_{\rm LL}(\mathbf u)\) is the head \(W_{\rm LL}\) applied to
  \(\mathbf u\), written \texttt{W\_LL{[}u{]}} in the Wolfram Language.
\item
  Indices (\(p\), \(j\), \(\iota\)) are not symbolized: they are the
  running letters inside the scripts; their components are symbolized on
  demand (\texttt{SL31Instance}). Primed indices are symbolized, because
  they also appear alone.
\item
  Operators, number sets and the imaginary unit are not symbolized: they
  are Wolfram functions, domains and the symbol \texttt{I} (S-13).
\item
  The single letter \(D\) is symbolized in the box
  \texttt{TagBox{[}"D",\ "SL31Plain"{]}}, so that the letter \texttt{D},
  typed plainly, keeps its meaning (the derivative) in the Wolfram
  Language.
\item
  Values are given to symbols only inside \texttt{Block}, so that after
  every test the symbols are free again.
\item
  The registry is contained in the notebook as data (one input cell of
  about 133 000 characters) and also shown as formatted tables, group by
  group (section I.8).
\item
  Part II rewrites the original programs in the order of their code and
  draws the random numbers in the same order with the same seeds, so
  that every random object of a check equals the one of the original
  program; results that are pure rounding errors are reported next to
  the values of the original report (section 10.2).
\item
  The numbers of the three data files of the figures of item {[}1{]} are
  copied into the notebook (about 60 000 characters in one input cell)
  instead of being read from the files, so that the notebook is complete
  and stand-alone; they are used only by the check T2-DATA-1. The cell
  makes the printed notebook about 30 pages longer.
\item
  The symbolic algebra \texttt{sym.rs} of the original library is not
  rewritten, because the Wolfram Language is itself a symbolic algebra
  (section 11.4); GS-12 uses symbols of the Wolfram Language instead.
\item
  Every name of Part II stands for one object only. For example, the
  variable of the sum over the plane waves of GS-5 is called
  \texttt{wave}, because the name \texttt{wv} belongs to the random
  vector of GS-12.
\item
  Labels of graphics are built from expressions (\texttt{Row},
  \texttt{Subscript}, \texttt{Superscript} and the registered symbols),
  not from strings in linear syntax, so that the source stays readable
  and the labels are typeset in the same way in the notebook and in the
  PDF.
\item
  The checks T2-SYM, T2-RNG, T2-LINALG and T2-DATA-1 are checks of the
  rewriting itself and are not in the original report; GS-1 to GS-16
  keep the names, the statements and the tolerances of the original
  report.
\item
  The components of families are collected by the builder from the
  source rather than listed by hand, so that a component typed in a
  later item of Part II is made automatically (section 10.4).
\end{itemize}

\hypertarget{files}{%
\section{16. Files}\label{files}}

All paths are relative to the top folder of the repository.

\begin{longtable}[]{@{}
  >{\raggedright\arraybackslash}p{(\columnwidth - 2\tabcolsep) * \real{0.5000}}
  >{\raggedright\arraybackslash}p{(\columnwidth - 2\tabcolsep) * \real{0.5000}}@{}}
\toprule\noalign{}
\begin{minipage}[b]{\linewidth}\raggedright
file
\end{minipage} & \begin{minipage}[b]{\linewidth}\raggedright
content
\end{minipage} \\
\midrule\noalign{}
\endhead
\bottomrule\noalign{}
\endlastfoot
\texttt{3+1super-loop.nb} & the notebook, evaluated (Part I; Part II:
the library and item {[}1{]}) \\
\texttt{3+1super-loop\_provenance/notebook\_source/3+1super-loop.source.wl}
& the source of the notebook (cells marked in the Wolfram package
format) \\
\texttt{3+1super-loop\_provenance/notebook\_source/sl31\_core.wl} & the
package
SL31\texttt{(TeX\ reader,\ kinds,\ names,\ symbolization),\ included\ in\ the\ notebook\ \textbar{}\ \textbar{}}3+1super-loop\_provenance/notebook\_source/registry\_data.wl\texttt{\textbar{}\ the\ registry\ as\ Wolfram\ Language\ data\ (199\ entries),\ included\ in\ the\ notebook\ \textbar{}\ \textbar{}}3+1super-loop\_provenance/notebook\_source/meaning\_edits.wl\texttt{\textbar{}\ the\ 16\ restated\ meanings\ applied\ by\ registry\_to\_wl.wls\ \textbar{}\ \textbar{}}3+1super-loop\_provenance/notebook\_source/all\_symbols.wl\texttt{\textbar{}\ the\ 422\ symbols\ as\ LL{[}tex{]}\ references,\ included\ in\ the\ cell\ of\ test\ S-14\ \textbar{}\ \textbar{}}3+1super-loop\_provenance/notebook\_source/3+1super-loop.unevaluated.nb\texttt{\textbar{}\ the\ notebook\ as\ built,\ before\ evaluation\ \textbar{}\ \textbar{}}3+1super-loop\_provenance/export/3+1super-loop.pdf\texttt{\textbar{}\ the\ evaluated\ notebook\ printed\ to\ PDF\ by\ the\ front\ end\ \textbar{}\ \textbar{}}3+1super-loop\_provenance/logs/3+1super-loop.evaluation.log.txt\texttt{\textbar{}\ the\ log\ of\ the\ final\ evaluation:\ messages,\ printed\ lines,\ input\ labels,\ tests\ \textbar{}\ \textbar{}}3+1super-loop\_provenance/logs/symbolization\_one\_by\_one.txt\texttt{\textbar{}\ diagnostic:\ every\ item\ symbolized\ alone,\ with\ its\ result\ (no\ item\ failed\ or\ hung)\ \textbar{}\ \textbar{}}3+1super-loop\_provenance/logs/E07\_front\_end\_boxes\_before\_fix.txt\texttt{\textbar{}\ E7:\ the\ 9\ boxes\ changed\ by\ the\ front\ end\ \textbar{}\ \textbar{}}3+1super-loop\_provenance/logs/E07\_leaf\_check\_front\_end\_parser.txt\texttt{\textbar{}\ E7:\ the\ front\ end\textquotesingle{}s\ parser\ on\ all\ 213\ strings\ \textbar{}\ \textbar{}}3+1super-loop\_provenance/logs/E07\_front\_end\_boxes\_after\_fix.txt\texttt{\textbar{}\ E7:\ 0\ boxes\ changed\ after\ the\ fix\ \textbar{}\ \textbar{}}3+1super-loop\_provenance/logs/E09\_kinds\_and\_shared\_boxes.txt\texttt{\textbar{}\ E9,\ E11:\ kinds\ and\ the\ TeX\ items\ with\ shared\ boxes\ \textbar{}\ \textbar{}}3+1super-loop\_provenance/logs/E09\_operand\_check\_in\_kernel\_input.txt\texttt{\textbar{}\ E9:\ every\ box\ reaches\ the\ kernel\ unchanged\ in\ the\ cell\ \{b,\ 1\}\ \textbar{}\ \textbar{}}3+1super-loop\_provenance/logs/E09\_reparse\_check.txt\texttt{\textbar{}\ E9:\ NotebookPut\ and\ NotebookGet\ change\ no\ box\ \textbar{}\ \textbar{}}3+1super-loop\_provenance/logs/E09\_s14\_chunk\_experiment.log.txt\texttt{\textbar{}\ E9:\ the\ list\ of\ symbols\ evaluated\ in\ cells\ of\ ten\ \textbar{}\ \textbar{}}3+1super-loop\_provenance/experiments/gui\_test.unevaluated.nb\texttt{\textbar{}\ E7:\ the\ test\ notebook\ with\ SubscriptBox{[}"a",\ "LL,p"{]}\ \textbar{}\ \textbar{}}3+1super-loop\_provenance/experiments/gui\_test.evaluated.nb\texttt{\textbar{}\ E7:\ the\ same\ notebook\ evaluated\ in\ the\ front\ end\ \textbar{}\ \textbar{}}3+1super-loop\_provenance/experiments/gui\_test.evaluation.log.txt\texttt{\textbar{}\ E7:\ its\ log\ \textbar{}\ \textbar{}}3+1super-loop\_provenance/tools/build\_notebook.wls\texttt{\textbar{}\ the\ builder:\ source\ to\ notebook\ (includes,\ components\ of\ families,\ alphabets\ of\ text\ formulas)\ \textbar{}\ \textbar{}}3+1super-loop\_provenance/tools/registry\_to\_wl.wls\texttt{\textbar{}\ the\ registry\ CSV\ file\ to\ registry\_data.wl\ \textbar{}\ \textbar{}}3+1super-loop\_provenance/tools/gen\_all\_symbols.wls\texttt{\textbar{}\ writes\ all\_symbols.wl\ \textbar{}\ \textbar{}}3+1super-loop\_provenance/tools/guirun.wl\texttt{\textbar{}\ evaluation\ of\ a\ notebook\ in\ the\ front\ end\ of\ Wolfram.app\ in\ a\ fresh\ kernel;\ saving,\ PDF\ export\ and\ log\ \textbar{}\ \textbar{}}3+1super-loop\_provenance/tools/start\_main.wl\texttt{\textbar{}\ starts\ the\ evaluation\ of\ the\ notebook\ with\ guirun.wl\ \textbar{}\ \textbar{}}3+1super-loop\_provenance/tools/status.wl\texttt{\textbar{}\ asks\ guirun.wl\ for\ the\ state\ of\ the\ evaluation\ \textbar{}\ \textbar{}}3+1super-loop\_provenance/tools/watcher.wl\texttt{\textbar{}\ the\ scheduled\ task\ in\ the\ Wolfram.app\ session\ that\ reads\ request\ files\ \textbar{}\ \textbar{}}3+1super-loop\_provenance/tools/handler.wl\texttt{\textbar{}\ the\ handler\ of\ the\ request\ files\ (scripts,\ and\ files\ evaluated\ in\ the\ session)\ \textbar{}\ \textbar{}}3+1super-loop\_provenance/tools/run\_notebook.wls\texttt{\textbar{}\ the\ superseded\ evaluator\ run\ by\ wolframscript\ (E2,\ E6)\ \textbar{}\ \textbar{}}3+1super-loop\_provenance/tools/canon\_build.wls\texttt{\textbar{}\ E7:\ notebook\ with\ one\ cell\ per\ symbol\ box\ and\ a\ \$PreRead\ recorder\ \textbar{}\ \textbar{}}3+1super-loop\_provenance/tools/canon\_compare.wls\texttt{\textbar{}\ E7:\ compares\ recorded\ and\ built\ boxes\ \textbar{}\ \textbar{}}3+1super-loop\_provenance/tools/leaf\_check.wls\texttt{\textbar{}\ E7:\ the\ front\ end\textquotesingle{}s\ parser\ on\ every\ string\ of\ the\ boxes\ \textbar{}\ \textbar{}}3+1super-loop\_provenance/tools/canon2\_build.wls\texttt{\textbar{}\ E9,\ E11:\ cells\ \{b,\ 1\}\ with\ and\ without\ the\ SyntaxForm\ wrapper;\ shared\ boxes\ \textbar{}\ \textbar{}}3+1super-loop\_provenance/tools/canon2\_compare.wls\texttt{\textbar{}\ E9:\ compares\ the\ recorded\ cells\ \textbar{}\ \textbar{}}3+1super-loop\_provenance/tools/reparse\_check.wls\texttt{\textbar{}\ E9:\ writes\ and\ reads\ every\ cell\ \{x,\ b,\ y\}\ through\ the\ front\ end\ \textbar{}\ \textbar{}}3+1super-loop\_provenance/tools/s14\_build.wls\texttt{\textbar{}\ E9:\ the\ list\ of\ all\ symbols\ split\ into\ cells\ of\ ten\ \textbar{}\ \textbar{}}3+1super-loop\_provenance/tools/build\_provenance\_docs.sh\texttt{\textbar{}\ makes\ the\ .tex\ and\ .pdf\ versions\ of\ this\ document\ from\ the\ .md\ \textbar{}\ \textbar{}}3+1super-loop\_provenance/tools/doc\_header.tex\texttt{\textbar{}\ the\ LaTeX\ header\ used\ by\ build\_provenance\_docs.sh\ \textbar{}\ \textbar{}}3+1super-loop\_provenance/notebook\_source/part2\_A\_core.source.wl\texttt{\textbar{}\ the\ source\ of\ the\ introduction\ of\ Part\ II,\ of\ its\ first\ cell\ (the\ components\ of\ families)\ and\ of\ chapter\ II.1\ (the\ library\ t2core)\ \textbar{}\ \textbar{}}3+1super-loop\_provenance/notebook\_source/part2\_B\_item1.source.wl\texttt{\textbar{}\ the\ source\ of\ item\ {[}1{]},\ checks\ GS-1\ to\ GS-11\ \textbar{}\ \textbar{}}3+1super-loop\_provenance/notebook\_source/part2\_C\_item1b.source.wl\texttt{\textbar{}\ the\ source\ of\ item\ {[}1{]},\ checks\ GS-12\ to\ GS-16,\ the\ figures,\ T2-DATA-1\ and\ the\ summary\ of\ the\ item\ \textbar{}\ \textbar{}}3+1super-loop\_provenance/notebook\_source/t2\_fields\_figure\_data.wl\texttt{\textbar{}\ the\ numbers\ of\ the\ three\ data\ files\ of\ the\ figures\ of\ t2\_fields,\ included\ in\ the\ notebook\ (T2-DATA-1)\ \textbar{}\ \textbar{}}3+1super-loop\_provenance/tools/make\_fig\_data\_cell.py\texttt{\textbar{}\ writes\ t2\_fields\_figure\_data.wl\ from\ the\ data\ files\ of\ the\ original\ program\ (Python\ 3)\ \textbar{}\ \textbar{}}3+1super-loop\_provenance/tools/splitmix64\_reference.py\texttt{\textbar{}\ the\ reference\ values\ of\ the\ check\ T2-RNG,\ computed\ in\ Python\ 3\ \textbar{}\ \textbar{}}3+1super-loop\_provenance/tools/e15\_dot\_precedence.wls\texttt{\textbar{}\ E15:\ the\ parsed\ form\ of\ the\ faulty\ expression\ and\ the\ behaviour\ of}Dot\texttt{with\ a\ number\ \textbar{}\ \textbar{}}3+1super-loop\_provenance/tools/e17\_tex\_importer.wls\texttt{\textbar{}\ E17:\ what\ the\ TeX\ importer\ makes\ of\ styled\ letters,\ and\ the\ boxes\ of\ the\ text\ formulas\ of\ the\ built\ notebook\ \textbar{}\ \textbar{}}3+1super-loop\_provenance/logs/T2-RNG\_reference\_values.txt\texttt{\textbar{}\ the\ output\ of\ splitmix64\_reference.py\ \textbar{}\ \textbar{}}3+1super-loop\_provenance/logs/E15\_dot\_precedence.txt\texttt{\textbar{}\ E15:\ the\ output\ of\ e15\_dot\_precedence.wls\ \textbar{}\ \textbar{}}3+1super-loop\_provenance/logs/E17\_tex\_importer.txt\texttt{\textbar{}\ E17:\ the\ output\ of\ e17\_tex\_importer.wls\ \textbar{}\ \textbar{}}3+1super-loop\_provenance/logs/E16\_figures\_and\_table\_before\_fix.png\texttt{\textbar{}\ E16:\ the\ pages\ of\ the\ figures\ and\ of\ the\ summary\ table\ before\ the\ fix\ \textbar{}\ \textbar{}}3+1super-loop\_provenance/logs/E16\_figures\_and\_table\_after\_fix.png\texttt{\textbar{}\ E16:\ the\ same\ pages\ after\ the\ fix\ \textbar{}\ \textbar{}}3+1super-loop\_provenance/logs/E17\_text\_formulas\_before\_and\_after\_fix.png\texttt{\textbar{}\ E17:\ a\ text\ of\ section\ II.1.3\ before\ and\ after\ the\ fix\ \textbar{}\ \textbar{}}3+1super-loop\_provenance/logs/E18\_components\_before\_fix.png\texttt{\textbar{}\ E18:\ the\ table\ of\ the\ basic\ matrices\ and\ the\ axis\ label\ of\ Figure\ 1\ before\ the\ fix\ \textbar{}\ \textbar{}}3+1super-loop\_provenance/logs/E18\_components\_after\_fix.png` \\
\end{longtable}

The SHA-256 digests of these files, as committed (output of
\texttt{sha256sum}; the three versions of this document cannot list
their own digests):

\begin{verbatim}
1bf6e439ad07e37444696b5f661059398a5450f0410ae0cdc96925b5e1cdf3ae  3+1super-loop.nb
ee5ad6db199bed6dcc39c84bb5cfe5755e55efdfc0b4bc2c4ab02e62b688edb5  3+1super-loop_provenance/notebook_source/3+1super-loop.source.wl
36cf048620323e4ac04d5ebf285576e2e191404791ac666a780e1165791357ef  3+1super-loop_provenance/notebook_source/sl31_core.wl
e3c8cbaa0005577b0d984daaea939c8ceffabc7479dde57e32e11714f7be48f1  3+1super-loop_provenance/notebook_source/registry_data.wl
5b5f96ef59e2057b284b362e7e49e67423b3d122bf1e0db37abd3df81a550fde  3+1super-loop_provenance/notebook_source/meaning_edits.wl
d35c3bf0c6d80a1c8d95b71e6602871755caed36e0551ec9639694fa81a83f87  3+1super-loop_provenance/notebook_source/all_symbols.wl
f56b422c9514bf956cc3f4d72bf51206effa8373d3a78ebc048a89a189a0f026  3+1super-loop_provenance/notebook_source/3+1super-loop.unevaluated.nb
8e0f6f1e85e7898904c121e2da3fe690e723c5a9b68573f64a0d973a4b1b800f  3+1super-loop_provenance/export/3+1super-loop.pdf
38defe45b116d23adc9577fe2d869d0652569d0554bfbb1fd63a7ff4e3877677  3+1super-loop_provenance/logs/3+1super-loop.evaluation.log.txt
664d1e799a0aa5114b03f9b10db89b7ec45f66897d40601b7970fa5db426dae6  3+1super-loop_provenance/logs/symbolization_one_by_one.txt
fe4de274da007c4ae06a5ede22c05776f1279f1515225b904d693b1f02ea7ed8  3+1super-loop_provenance/logs/E07_front_end_boxes_before_fix.txt
94a1d676722de723e6ec58c2c2b1489df566a457955d33b46b594802ca7554cc  3+1super-loop_provenance/logs/E07_leaf_check_front_end_parser.txt
a5361540c3674a70cc39b2bb16019591ca32a75a0ba94e90a94b45c90ea99907  3+1super-loop_provenance/logs/E07_front_end_boxes_after_fix.txt
c5b7e2c3f40e10c02df896742ac1e1eecc5d2589d68bf29bcc88ef2046a3bac9  3+1super-loop_provenance/logs/E09_kinds_and_shared_boxes.txt
6cc8a6be8f70c3910140d5ed17872d240b7d44e98709a23e5508b7b52a33469e  3+1super-loop_provenance/logs/E09_operand_check_in_kernel_input.txt
557a7c3dbb539ac1f7607509942374dc4373bc41657c3eb25d8d9c41ccd27d3c  3+1super-loop_provenance/logs/E09_reparse_check.txt
747754b20292e9f1427b0249f86964ac9dd74380da26ec2f0b13ecac05925341  3+1super-loop_provenance/logs/E09_s14_chunk_experiment.log.txt
7aee2310fa163ab02afe2881a46ba51f08590e3e4ced458c8986fdee17cdd955  3+1super-loop_provenance/experiments/gui_test.unevaluated.nb
cb0bc845dd86a23fd7bbd1fd5b18d6d9a2a2caf4b26a0ab58d1c136f5ccd42bc  3+1super-loop_provenance/experiments/gui_test.evaluated.nb
de33e81f087369495076880b88f0ad3e2eb83cdeacd9f346b04180e21fa8e66d  3+1super-loop_provenance/experiments/gui_test.evaluation.log.txt
5f659852270513081501df3aa23d66699e702364cc27bd242006c3b37256f7dc  3+1super-loop_provenance/tools/build_notebook.wls
42767cd1ebc43492c95d19c5d67507a6a2450da9f0893f41f367f681a9ce5bd5  3+1super-loop_provenance/tools/registry_to_wl.wls
46211d0b974221fb90e02b5459c3fc84577b35fb7ce732b565aa428ae1c5c604  3+1super-loop_provenance/tools/gen_all_symbols.wls
6a7d810350cfcd4f69cc12b7d4bbc92ca1d57dcafd4a6ce6a1ecef4b37beef8d  3+1super-loop_provenance/tools/guirun.wl
9083afbd62a52895d166dbb1e7efcae9820834ca503cd5f21e6863336785a37c  3+1super-loop_provenance/tools/start_main.wl
d08e464e906ef6b296b3c6b27a89cded6a447551df3faeb77f96d34ecbe2029c  3+1super-loop_provenance/tools/status.wl
a20ca599120a23107d27a7c4f024b3a137303d68868a11293032f4f76f3f5436  3+1super-loop_provenance/tools/watcher.wl
3511e534722ca460f99dfc5e6479ae51eeb33a1f89fb9a41f440689a26c492bc  3+1super-loop_provenance/tools/handler.wl
24ca3f2b9246b58759f4c1f7a1d447ad9bf6003d5a838e88ede0b503dc1de871  3+1super-loop_provenance/tools/run_notebook.wls
33da3e6e322c6c184900c25167bddba499c5a9c2dbeee0d65ec760bc7aee7d92  3+1super-loop_provenance/tools/canon_build.wls
049fd2e57dc8a5551fc90f12ca64e0bdfa0d12e64926a8b4e7584b6dadbfcc11  3+1super-loop_provenance/tools/canon_compare.wls
7e1e688b9d26606103d84ff246aaa8834a55b53cbdbae1d0fc0bb71480f54d6a  3+1super-loop_provenance/tools/leaf_check.wls
a79915c5253a7cdd9b162eda7b9e466ec3b57a63d670f614f1890c983bdc28e7  3+1super-loop_provenance/tools/canon2_build.wls
05faf4a855ce53e514a34a2966ae43345df7e613a5e10fe459a99292437011ed  3+1super-loop_provenance/tools/canon2_compare.wls
7803f4410097cbf0af6a7cd557bdc54cb0c75c9fe26411dc20eb51a3fa6640f3  3+1super-loop_provenance/tools/reparse_check.wls
ddd7824acd0e8c1746ec4548d8912f9cf1d179438f18b69e47a90978687881a1  3+1super-loop_provenance/tools/s14_build.wls
518fc0768f4da0aa5528e11f5207be50dbe453051b96dc03917352296cb0670b  3+1super-loop_provenance/tools/build_provenance_docs.sh
00d3e3c225b60acac822b408deed88744968de7f8ef78467f6799ea1b68a3d30  3+1super-loop_provenance/tools/doc_header.tex
131f786bf2b89cf4ba2f219320dd75ed5202dd6481500a21545e0dc9844fae26  3+1super-loop_provenance/notebook_source/part2_A_core.source.wl
3c323752b0e34c0775f6041104bc8479b196647e4ea0ae098bb545e93a04c833  3+1super-loop_provenance/notebook_source/part2_B_item1.source.wl
97ea97f6302be0b554d3fb9c4ce285047f3f287c14f53b7fe5a6695c3cebdac0  3+1super-loop_provenance/notebook_source/part2_C_item1b.source.wl
277b1e4817587d9a5f5aa80620d94910edc32b423c23a5888d79150f80923334  3+1super-loop_provenance/notebook_source/t2_fields_figure_data.wl
7a46cd1d8efcd05bc2839e9abd62dce628753a7a75269f44a65b2976407a0a18  3+1super-loop_provenance/tools/make_fig_data_cell.py
06be9891ca81c88602943cabdca85f971193481d5f669207352b1a18bc53dca5  3+1super-loop_provenance/tools/splitmix64_reference.py
f28f35adb6abc25e59f0e6f389e4566632631d238d20e044c27ba23247ff437b  3+1super-loop_provenance/tools/e15_dot_precedence.wls
f34de46eaa0bbb6a7261868c62dc076e3927185bfcc0fdea91c0403770d172cb  3+1super-loop_provenance/tools/e17_tex_importer.wls
ceafb181fbec77e8efb288b8f2c25ee73aba60c5121779d12da90840b1e2ad19  3+1super-loop_provenance/logs/T2-RNG_reference_values.txt
c64f1a909b9d008365fd5707547547cceb6f0df78b314099c676ba7d1e711f16  3+1super-loop_provenance/logs/E15_dot_precedence.txt
0b051cae79f44a26285ece8d21c4c46994342dd7a0f529d2358c35b621358647  3+1super-loop_provenance/logs/E17_tex_importer.txt
c719036debc776aab97c49fe2cb07f56e0e5bc7f14c0b4ee334f6e429265521b  3+1super-loop_provenance/logs/E16_figures_and_table_before_fix.png
dfe1af9e06b087239047f4f4ed05fea6bb9d27b0e14f4b738296b64df81e9298  3+1super-loop_provenance/logs/E16_figures_and_table_after_fix.png
8d1981333fc22357b76b79d9b6d30ef86579a24425067578d8fde2b48c67399a  3+1super-loop_provenance/logs/E17_text_formulas_before_and_after_fix.png
95f9e2e5adde8e72e471cd7a0bb0b81cf491c25a80e7b3857ca949c25e6eb82a  3+1super-loop_provenance/logs/E18_components_before_fix.png
87c29f8d4a0fedc884dce45f905201f9ac1a95c1b2b2052c35d02896d442c6f7  3+1super-loop_provenance/logs/E18_components_after_fix.png
\end{verbatim}

\hypertarget{sources-consulted}{%
\section{17. Sources consulted}\label{sources-consulted}}

\begin{itemize}
\tightlist
\item
  The source code of the Notation package in the Wolfram 15.0.1
  installation (\texttt{Notation.m}): the default of
  \texttt{WorkingForm}, the removal of \texttt{StyleBox},
  \texttt{AdjustmentBox} and \texttt{FrameBox} before comparison, the
  conversion of strings that end in an underscore into patterns, the
  warning \texttt{Symbolize::bsymbexs} (E1, E5, E11).
\item
  The Wolfram Language reference page of \texttt{SyntaxForm}: an option
  of \texttt{TagBox}, \texttt{InterpretationBox} and
  \texttt{TemplateBox} that specifies how a box takes part in the
  grammar of the surrounding input; by default the front end treats a
  box as an operand (E9).
  \url{https://reference.wolfram.com/language/ref/SyntaxForm.html}
\item
  The reference page of \texttt{TagBox}: its options
  \texttt{StripWrapperBoxes} and \texttt{SyntaxForm} (E9).
  \url{https://reference.wolfram.com/language/ref/TagBox.html}
\item
  The reference page of \texttt{StripBoxes}: when boxes are sent by the
  front end to the kernel for evaluation, the front end strips them with
  a more extensive process (E7).
  \url{https://reference.wolfram.com/language/ref/StripBoxes.html}
\item
  The Notation package tutorial ``Complex Patterns and Advanced
  Features'': tag boxes and the option \texttt{SyntaxForm} in notations
  (E9).
  \url{https://reference.wolfram.com/language/Notation/tutorial/ComplexPatternsAndAdvancedFeatures.html}
\item
  J. Fultz, answers on Mathematica Stack Exchange about the internal
  command of the front end that reads a string as typed input (its
  second argument \texttt{True} makes the stripping that is made at
  evaluation time) (E7).
\item
  J. Fultz, MathGroup mailing list, 21 January 2014, ``Re: Unnecessary
  parentheses around TagBox for MakeBoxes?'': the value
  \texttt{SyntaxForm\ -\textgreater{}\ "symbol"} for a box that should
  act as a standalone thing (E9).
\item
  The Wolfram Language tutorial ``Operator Input Forms'': the table of
  the operators in order of decreasing precedence, in which
  \texttt{expr1\ .\ expr2} (\texttt{Dot}) stands above the implicit
  product \texttt{expr1\ expr2} (\texttt{Times}) (E15).
  \url{https://reference.wolfram.com/language/tutorial/OperatorInputForms.html}
\item
  The reference page of \texttt{SingleLetterItalics}: a cell option that
  controls whether names of one letter are printed in italics, typically
  \texttt{True} for cells of \texttt{TraditionalForm} (E18, E19).
  \url{https://reference.wolfram.com/language/ref/SingleLetterItalics.html}
\item
  The original programs of Test 2 in the repository
  (\texttt{Revision\_SuperLoop/test2/rust}: the library \texttt{t2core}
  and the program \texttt{t2fields}), which Part II rewrites, and the
  record of their checks and the data files of their figures
  (\texttt{Revision\_SuperLoop/test2/data}). Their SHA-256 digests, as
  used for this stage:
\end{itemize}

\begin{verbatim}
4a7beee875494e11a5778aa9bd94fca841bc5ab0abff6a3666a6ad25e72ac731  Revision_SuperLoop/test2/rust/t2core/src/rng.rs
8c412f5703dc49c7ba51fb86ff9a4fd4bcc3ce4667a11216884a4e2e2b6a43a3  Revision_SuperLoop/test2/rust/t2core/src/c64.rs
ace63b498d0e0fcbca02f1b659aa6249343f1b9ef9b16a113349bb9023915447  Revision_SuperLoop/test2/rust/t2core/src/cmat.rs
aecd9e01a40e211133cc8ff2c4697ec62e0370b46e0d7ea2d31fb19613a9eeb4  Revision_SuperLoop/test2/rust/t2core/src/model.rs
046298ecaff5d17bcfe0c0747c34b4447ebef19291330757ee4058d5424b1ae7  Revision_SuperLoop/test2/rust/t2core/src/lattice.rs
790b8758ce9bf919ffd93eb7f3afdaf66b4b4abb640a1fbf205146feeae6a8c7  Revision_SuperLoop/test2/rust/t2core/src/sym.rs
26e7cc6bb7ac393f83bfaa222eeb0f2e0639f3ab878fbbd9e6e67d20c1376f68  Revision_SuperLoop/test2/rust/t2core/src/lib.rs
d7ed53e0635b63e4f84d23a7280bf8d3581086a6a64dd3df7cb6632f49063939  Revision_SuperLoop/test2/rust/t2fields/src/main.rs
033ddabf96d8a2e2e6fbf4d9bbd95ca32ea102b15c003848e9ee52ce8aabe0a7  Revision_SuperLoop/test2/data/t2_fields_checks.txt
5e95b089a62bbb22de0736a71128c6aea361a085aecd83569447cca0462d89d7  Revision_SuperLoop/test2/data/t2_dispersion_lattice.csv
2f2849f03bb13b8475272648b8ec80617437c28a51afed9b7c2c11fabac2db01  Revision_SuperLoop/test2/data/t2_dispersion_curves.csv
64544a7caeba4e37eae7f0754093f0e822d222e23ef07ce8663cb4b0e73e402c  Revision_SuperLoop/test2/data/t2_rk4_stability.csv
\end{verbatim}

\end{document}
